% Les mécanismes de la mondialisation — Géographie Tle Tech — Cours signé OnBuch+
% Programme officiel MINESEC. Compiler avec : tectonic N04-mecanismes-mondialisation.tex
\newcommand{\DOCMATIERE}{Géographie}
\newcommand{\DOCNIVEAU}{Tle Tech}
\newcommand{\DOCMODULE}{Géographie – classe de Terminale technique}
\newcommand{\DOCLECON}{N04}
\newcommand{\DOCTITRE}{Les mécanismes de la mondialisation}
% OnBuch+ — préambule « cours signés » ENRICHI (compilé avec tectonic / XeLaTeX)
% Système visuel « Pop » d'OnBuch+ : fond crème (jamais blanc), bordures encre
% épaisses, ombres « dures » décalées, titres Archivo Black en capitales.
% Version MATHÉMATIQUES : graphes (pgfplots/tikz), boîtes pédagogiques
% (définition, propriété, méthode, attention, le savais-tu, exercice résolu,
% exercice, TP, bilan), page de garde et signature OnBuch+ ancrée en bas.
%
% Macros attendues dans main.tex AVANT \input{preamble} :
%   \DOCMATIERE  \DOCNIVEAU  \DOCMODULE  \DOCLECON (numéro)  \DOCTITRE
\documentclass[11pt]{article}

\usepackage{fontspec}
\usepackage[french]{babel}
\usepackage[a4paper,margin=2cm,top=3.4cm,bottom=2.8cm,headheight=1.4cm,footskip=1.3cm]{geometry}
\usepackage{amsmath,amssymb,amsfonts}
\usepackage{mathtools} % \xmapsto, \coloneqq... souvent utilisés par les modèles
\usepackage{cancel} % simplifications barrées (bilans de forces)
% \abs{z} (module, valeur absolue) et \norm{u} (norme), utilisés par les modèles
\providecommand{\abs}{}\let\abs\relax\DeclarePairedDelimiter{\abs}{\lvert}{\rvert}
\providecommand{\norm}{}\let\norm\relax\DeclarePairedDelimiter{\norm}{\lVert}{\rVert}
\usepackage[table]{xcolor} % [table] : fournit aussi \rowcolors (alternance de lignes)
\usepackage{pagecolor}
\usepackage{tikz}
\usetikzlibrary{arrows.meta,positioning,calc,shapes.geometric,shapes.misc,shapes.arrows,decorations.pathmorphing,decorations.pathreplacing,decorations.markings,patterns,shadows,mindmap,trees,fit,backgrounds,matrix,intersections,angles,quotes}
\usepackage{pgfplots}
\usepackage[version=4]{mhchem} % équations nucléaires \ce{^{235}_{92}U}
\usepackage[european,siunitx]{circuitikz} % schémas électriques
\pgfplotsset{compat=1.18}
\usepgfplotslibrary{fillbetween,groupplots} % aires entre courbes (calcul intégral)
\usepackage{enumitem}
\usepackage{fancyhdr}
% [most] charge toutes les bibliothèques tcolorbox (listings, minted, raster,
% documentation...) et gonfle inutilement la mémoire TeX sur un document long
% (~300 boîtes) : on ne charge que ce qui est réellement utilisé.
\usepackage[skins,breakable]{tcolorbox}
\usepackage{graphicx}
\usepackage{colortbl}
\usepackage{booktabs}
\usepackage{tabularx}
\usepackage{multirow}
\usepackage{diagbox} % case d'en-tête diagonale des tableaux à double entrée
% Colonnes extensibles centrée / gauche / droite (C, L, R), souvent utilisées par les modèles
\newcolumntype{C}{>{\centering\arraybackslash}X}
\newcolumntype{L}{>{\raggedright\arraybackslash}X}
\newcolumntype{R}{>{\raggedleft\arraybackslash}X}
\usepackage{array}
\usepackage{multicol}
\usepackage{siunitx}
% parse-numbers=false : accepte aussi \SI{2,5 \times 10^{-3}}{...}
\sisetup{locale=FR,output-decimal-marker={,},group-minimum-digits=4,parse-numbers=false}
\DeclareSIUnit\ohm{\ensuremath{\symup{\Omega}}} % U+2126 absent de la police
\DeclareSIUnit\division{div} % réglages d'oscilloscope (ms/div, V/div)
\DeclareSIUnit\tr{tr} % tours (vitesse de rotation, ex. \SI{1500}{\tr\per\minute})
\DeclareSIUnit\molar{M} % concentration molaire (ex. \SI{3}{\molar}, \SI{1}{\micro\molar})
\DeclareSIUnit\an{an} % année (le modèle confond parfois avec \year, un compteur TeX) — ex. \SI{5}{\kilo\meter\per\an}
\DeclareSIUnit\FCFA{FCFA} % franc CFA (ex. \SI{500}{\FCFA\per\kilo\gram})
\DeclareSIUnit\degreeBrix{\textdegree Brix} % teneur en sucre d'un sirop/jus (réfractométrie)
\DeclareSIUnit\month{mois} % le modèle utilise parfois \month, un compteur TeX, comme unité de durée
\DeclareSIUnit\microliter{\micro\liter} % le modèle écrit parfois \microliter en un seul token
\DeclareSIUnit\million{millions} % dénombrement (ex. \SI{15}{\million\per\mL} spermatozoïdes)
\usepackage{qrcode}
% \degree : commande fréquemment utilisée par les modèles pour un angle ou une
% température en dehors de \SI{}{} (non fournie par les paquets chargés).
\providecommand{\degree}{^{\circ}}
% \qed (fin de démonstration), utilisé par les modèles sans amsthm
\providecommand{\qed}{\hfill$\square$}
% \barre : double barre des valeurs interdites dans un tableau de variations (tkz-tab)
\providecommand{\barre}{\ensuremath{\|}}
% environnement proof (amsthm non chargé) utilisé par les modèles
\ifdefined\proof\else\newenvironment{proof}[1][Démonstration]{\par\noindent\textit{#1.}\ }{\qed\par}\fi
\usepackage{lastpage}
\usepackage{hyperref}
\hypersetup{hidelinks}

% ============================================================
% Palette « Pop » OnBuch+ — mêmes hex que la classe P (plus_ui.dart)
% ============================================================
\definecolor{poporange}{HTML}{F59321}
\definecolor{popdark}{HTML}{E07A0C}
\definecolor{popcream}{HTML}{FFF6E8}
\definecolor{popink}{HTML}{1C1714}
\definecolor{poppurple}{HTML}{7A5AE0}
\definecolor{popgreen}{HTML}{2E9E6A}
\definecolor{poppink}{HTML}{E0457B}
\definecolor{popblue}{HTML}{2F7DE1}
\definecolor{popgold}{HTML}{C98A12}
\definecolor{popmuted}{HTML}{8A7F76}
\definecolor{popline}{HTML}{EADFD2}
\definecolor{popgray}{HTML}{8A7F76}
\colorlet{popgrey}{popgray}
% Alias tolérants (le modèle nomme parfois les couleurs différemment)
\colorlet{popwhite}{white}
\colorlet{popblack}{popink}
\colorlet{popred}{poppink}
\colorlet{popyellow}{popgold}
% Teintes claires (fonds de tableaux/encadrés) — le modèle nomme parfois
% les couleurs avec un suffixe L (ex. popblueL) sans les avoir définies.
\colorlet{poporangeL}{poporange!20}
\colorlet{popdarkL}{popdark!20}
\colorlet{poppurpleL}{poppurple!20}
\colorlet{popgreenL}{popgreen!20}
\colorlet{poppinkL}{poppink!20}
\colorlet{popblueL}{popblue!20}
\colorlet{popgoldL}{popgold!20}
\colorlet{popredL}{popred!20}
\colorlet{popgrayL}{popgray!20}
% Teintes claires pour fonds de tableaux / zones
\colorlet{popblueL}{popblue!10!white}
\colorlet{poporangeL}{poporange!14!white}
\colorlet{popgreenL}{popgreen!12!white}
\colorlet{poppurpleL}{poppurple!12!white}
\colorlet{poppinkL}{poppink!12!white}
\colorlet{popgoldL}{popgold!14!white}

\pagecolor{popcream}

% ============================================================
% Typographie
% ============================================================
\defaultfontfeatures{Ligatures=TeX}
\setmainfont{PlusJakartaSans-Medium.ttf}[Path=../../../fonts/,BoldFont=PlusJakartaSans-Bold.ttf]
% Maths en sans-serif (Fira Math) avec chiffres et lettres droites de Plus
% Jakarta, pour que les formules se fondent dans le texte.
\usepackage[math-style=ISO,bold-style=ISO]{unicode-math}
\setmathfont{FiraMath-Regular.otf}[Path=../../../fonts/]
\setmathfont{PlusJakartaSans-Medium.ttf}[Path=../../../fonts/,range={up/num,up/{Latin,latin},\mathup}]
\usepackage{icomma} % virgule décimale sans espace en mode math
% Caractères mathématiques Unicode absents des polices de texte (Plus Jakarta,
% Archivo Black) : rendus par leur équivalent mathématique, en texte comme en
% formule — sinon ils s'affichent en carré vide (ex. « Arithmétique dans ℤ »).
\usepackage{newunicodechar}
\newunicodechar{ℝ}{\ensuremath{\mathbb{R}}}
\newunicodechar{ℤ}{\ensuremath{\mathbb{Z}}}
\newunicodechar{ℂ}{\ensuremath{\mathbb{C}}}
\newunicodechar{ℕ}{\ensuremath{\mathbb{N}}}
\newunicodechar{ℚ}{\ensuremath{\mathbb{Q}}}
\newunicodechar{𝔻}{\ensuremath{\mathbb{D}}}
\newunicodechar{α}{\ensuremath{\alpha}}
\newunicodechar{β}{\ensuremath{\beta}}
\newunicodechar{γ}{\ensuremath{\gamma}}
\newunicodechar{δ}{\ensuremath{\delta}}
\newunicodechar{ε}{\ensuremath{\varepsilon}}
\newunicodechar{θ}{\ensuremath{\theta}}
\newunicodechar{λ}{\ensuremath{\lambda}}
\newunicodechar{μ}{\ensuremath{\mu}}
\newunicodechar{σ}{\ensuremath{\sigma}}
\newunicodechar{τ}{\ensuremath{\tau}}
\newunicodechar{φ}{\ensuremath{\varphi}}
\newunicodechar{ψ}{\ensuremath{\psi}}
\newunicodechar{ω}{\ensuremath{\omega}}
\newunicodechar{Σ}{\ensuremath{\Sigma}}
% majuscules produites par \MakeUppercase dans les titres (α-aminés -> Α)
\newunicodechar{Α}{\ensuremath{\alpha}}
\newunicodechar{Β}{\ensuremath{\beta}}
\newunicodechar{Γ}{\ensuremath{\Gamma}}
\newunicodechar{Δ}{\ensuremath{\Delta}}
\newunicodechar{Ε}{\ensuremath{\varepsilon}}
\newunicodechar{Θ}{\ensuremath{\Theta}}
\newunicodechar{Λ}{\ensuremath{\Lambda}}
\newunicodechar{Μ}{\ensuremath{\mu}}
\newunicodechar{Τ}{\ensuremath{\tau}}
\newunicodechar{Φ}{\ensuremath{\Phi}}
\newunicodechar{Ψ}{\ensuremath{\Psi}}
\newunicodechar{Ω}{\ensuremath{\Omega}}
\newunicodechar{↦}{\ensuremath{\mapsto}}
\newunicodechar{∈}{\ensuremath{\in}}
\newunicodechar{∉}{\ensuremath{\notin}}
\newunicodechar{∧}{\ensuremath{\wedge}}
\newunicodechar{⇔}{\ensuremath{\Leftrightarrow}}
\newunicodechar{⟺}{\ensuremath{\Longleftrightarrow}}
\newunicodechar{⇒}{\ensuremath{\Rightarrow}}
\newunicodechar{⟹}{\ensuremath{\Longrightarrow}}
\newunicodechar{∘}{\ensuremath{\circ}}
\newunicodechar{∪}{\ensuremath{\cup}}
\newunicodechar{∩}{\ensuremath{\cap}}
\newunicodechar{≡}{\ensuremath{\equiv}}
\newunicodechar{⊂}{\ensuremath{\subset}}
\newunicodechar{⊥}{\ensuremath{\perp}}
\newunicodechar{∃}{\ensuremath{\exists}}
\newunicodechar{∀}{\ensuremath{\forall}}
\newunicodechar{∛}{\ensuremath{\sqrt[3]{\,}}}
\newunicodechar{∜}{\ensuremath{\sqrt[4]{\,}}}
\newunicodechar{⃗}{\ensuremath{{}^{\to}}}
\newunicodechar{ⁿ}{\ifmmode^{n}\else\textsuperscript{\ensuremath{n}}\fi}
\newunicodechar{ˣ}{\ifmmode^{x}\else\textsuperscript{\ensuremath{x}}\fi}
\newunicodechar{ᵇ}{\ifmmode^{b}\else\textsuperscript{\ensuremath{b}}\fi}
\newunicodechar{ʳ}{\ifmmode^{r}\else\textsuperscript{\ensuremath{r}}\fi}
\newunicodechar{ᵘ}{\ifmmode^{u}\else\textsuperscript{\ensuremath{u}}\fi}
\newunicodechar{ʸ}{\ifmmode^{y}\else\textsuperscript{\ensuremath{y}}\fi}
\newunicodechar{ᵃ}{\ifmmode^{a}\else\textsuperscript{\ensuremath{a}}\fi}
\newunicodechar{ᵉ}{\ifmmode^{e}\else\textsuperscript{\ensuremath{e}}\fi}
\newunicodechar{ᵏ}{\ifmmode^{k}\else\textsuperscript{\ensuremath{k}}\fi}
\newunicodechar{ᵖ}{\ifmmode^{p}\else\textsuperscript{\ensuremath{p}}\fi}
\newunicodechar{ᴺ}{\ifmmode^{N}\else\textsuperscript{\ensuremath{N}}\fi}
\newunicodechar{ᶜ}{\ifmmode^{c}\else\textsuperscript{\ensuremath{c}}\fi}
\newunicodechar{ʰ}{\ifmmode^{h}\else\textsuperscript{\ensuremath{h}}\fi}
\newunicodechar{ᵠ}{\ifmmode^{q}\else\textsuperscript{\ensuremath{q}}\fi}
\newunicodechar{ₙ}{\ifmmode_{n}\else\textsubscript{\ensuremath{n}}\fi}
\newunicodechar{ₐ}{\ifmmode_{a}\else\textsubscript{\ensuremath{a}}\fi}
\newunicodechar{ᵢ}{\ifmmode_{i}\else\textsubscript{\ensuremath{i}}\fi}
\newunicodechar{ᵣ}{\ifmmode_{r}\else\textsubscript{\ensuremath{r}}\fi}
\newunicodechar{ₖ}{\ifmmode_{k}\else\textsubscript{\ensuremath{k}}\fi}
\newunicodechar{ᵦ}{\ifmmode_{\beta}\else\textsubscript{\ensuremath{\beta}}\fi}
\newunicodechar{ₓ}{\ifmmode_{x}\else\textsubscript{\ensuremath{x}}\fi}
\newfontfamily\popdisplay{ArchivoBlack-Regular.ttf}[Path=../../../fonts/]
\newfontfamily\poplogo{SpaceGrotesk-Bold.ttf}[Path=../../../fonts/]
\newfontfamily\popsemi{PlusJakartaSans-SemiBold.ttf}[Path=../../../fonts/]
\newfontfamily\popxbold{PlusJakartaSans-ExtraBold.ttf}[Path=../../../fonts/]
\newfontfamily\popxboldls{PlusJakartaSans-ExtraBold.ttf}[Path=../../../fonts/,LetterSpace=3.0]

\setlist[enumerate]{leftmargin=1.7em,itemsep=5pt,topsep=4pt}
\setlist[itemize]{leftmargin=1.7em,itemsep=4pt,topsep=4pt,label={\color{poporange}\textbullet}}
\setlength{\parindent}{0pt}
\setlength{\parskip}{5pt}
\renewcommand{\arraystretch}{1.25}

% pgfplots : style Pop par défaut
\pgfplotsset{
  every axis/.append style={axis line style={popink,line width=1pt},tick style={popink},
    grid style={popline,line width=0.5pt},label style={font=\small},tick label style={font=\footnotesize}},
  popcurve/.style={poporange,line width=1.8pt,no markers,smooth},
}
% Même style disponible dans un \draw TikZ simple (hors pgfplots)
\tikzset{popcurve/.style={poporange,line width=1.8pt}}
\tikzset{popgrid/.style={popline,very thin}}
\colorlet{popcurve}{poporange} % le nom du style est parfois utilisé comme couleur

% ============================================================
% Pill / squiggle
% ============================================================
\newcommand{\poppill}[3]{%
  \tikz[baseline=-0.55ex]\node[draw=popink,line width=1.2pt,rounded corners=2.6pt,%
    fill=#2,inner xsep=6.5pt,inner ysep=3pt]{%
    \popxboldls\fontsize{7}{7}\selectfont\color{#3}\MakeUppercase{#1}};%
}
\newcommand{\popsquiggle}[1]{%
  \tikz[baseline=-0.4ex]{
    \draw[#1,line width=2.6pt,line cap=round]
      plot [smooth] coordinates {(0,0.09) (0.34,0.01) (0.68,0.21) (1.02,0.01) (1.36,0.10)};
  }%
}
% Mot-clé surligné (marqueur orange sous le texte)
\newcommand{\cle}[1]{{\popxbold\color{popdark}#1}}

% ============================================================
% En-tête / pied
% ============================================================
\pagestyle{fancy}
\fancyhf{}
\renewcommand{\headrulewidth}{0pt}
\renewcommand{\footrulewidth}{0pt}
\lhead{\raisebox{-0.15ex}{\poplogo\fontsize{13}{13}\selectfont\color{popink}OnBuch\color{poporange}+}}
\rhead{\poppill{\DOCMATIERE\ \textbullet\ \DOCNIVEAU\ \textbullet\ Leçon \DOCLECON}{white}{popink}}
\lfoot{\popxbold\fontsize{7.6}{7.6}\selectfont\color{popmuted}\MakeUppercase{Cours signé OnBuch+}\ \textbullet\ {\popsemi\DOCTITRE}}
\rfoot{\tikz[baseline=-0.6ex]\node[rounded corners=8.5pt,fill=popink,minimum height=17pt,minimum width=17pt,inner xsep=5pt,inner sep=0]{\popxbold\fontsize{8}{8}\selectfont\color{popcream}\thepage\,/\,\pageref*{LastPage}};}

% ============================================================
% Titres
% ============================================================
\newcommand{\chaptitre}[2]{%
  \begin{tcolorbox}[enhanced,colback=poporange,colframe=popink,boxrule=2.6pt,arc=11pt,
      drop shadow=popink,left=15pt,right=15pt,top=12pt,bottom=11pt,before skip=3pt,after skip=18pt]
    \poppill{#2}{popink}{popcream}\par\vspace{9pt}
    {\raggedright\hyphenpenalty=10000\exhyphenpenalty=10000\popdisplay\fontsize{22}{24}\selectfont\color{popcream}\MakeUppercase{#1}\par}
  \end{tcolorbox}
}
% Section numérotée : pastille orange avec le numéro + titre Archivo
\newcounter{popsec}
\newcommand{\coursec}[1]{%
  \stepcounter{popsec}\setcounter{popsub}{0}%
  \par\vspace{16pt}\noindent
  \begin{minipage}{\linewidth}
  \tikz[baseline=-0.7ex]\node[draw=popink,line width=1.6pt,fill=poporange,rounded corners=4pt,minimum size=22pt,inner sep=0]{\popdisplay\fontsize{12}{12}\selectfont\color{popcream}\thepopsec};%
  \hspace{8pt}{\popdisplay\fontsize{15}{17}\selectfont\color{popink}\MakeUppercase{#1}}\par
  \vspace{4pt}\noindent{\color{poporange}\rule{36pt}{3.6pt}}
  \end{minipage}\par\nopagebreak\vspace{8pt}
}
\newcounter{popsub}
\newcommand{\courssub}[1]{%
  \stepcounter{popsub}%
  \par\vspace{9pt}\noindent{\popxbold\fontsize{11.5}{13}\selectfont\color{popink}{\color{poporange}\thepopsec.\thepopsub}\hspace{6pt}#1}\par\nopagebreak\vspace{4pt}
}

% ============================================================
% Boîtes pédagogiques — même recette : fond blanc, bordure épaisse colorée,
% ombre dure encre, badge plein. Titre optionnel [#1].
% ============================================================
\tcbset{popbase/.style={breakable,enhanced,colback=white,boxrule=2.3pt,arc=9pt,
  shadow={3pt}{-3pt}{0pt}{popink},left=14pt,right=14pt,top=11pt,bottom=11pt,before skip=11pt,after skip=15pt,
  fontupper=\popsemi\fontsize{10.6}{14.5}\selectfont\color{popink},
  fontlower=\popsemi\fontsize{10.6}{14.5}\selectfont\color{popink}}}
% Coche dessinée (le glyphe ✓ n'existe pas dans les polices du document)
\newcommand{\popcheck}[1]{\tikz[baseline=-0.5ex]\draw[#1,line width=1.6pt,line cap=round,line join=round](-0.13,0.0)--(-0.04,-0.09)--(0.13,0.1);}
\providecommand{\coche}{\popcheck{popgreen}}
\providecommand{\resultat}[1]{\par\smallskip\noindent\fcolorbox{popgreen}{popgreenL}{\parbox{\dimexpr\linewidth-2\fboxsep-2\fboxrule\relax}{\textbf{Résultat.} #1}}\par\smallskip}
\newcommand{\popboxhead}[3]{\poppill{#1}{#2}{white}\ifx\relax#3\relax\else\hspace{6pt}{\popxbold\fontsize{10.6}{12}\selectfont\color{popink}#3}\fi\par\vspace{7pt}}

\newtcolorbox{aretenir}[1][]{popbase,colframe=popgreen,before upper={\popboxhead{À retenir}{popgreen}{#1}}}
\newtcolorbox{exemplebox}[1][]{popbase,colframe=poporange,before upper={\popboxhead{Exemple}{poporange}{#1}}}
\newtcolorbox{definition}[1][]{popbase,colframe=popblue,before upper={\popboxhead{Définition}{popblue}{#1}}}
\newtcolorbox{propriete}[1][]{popbase,colframe=poppurple,before upper={\popboxhead{Propriété}{poppurple}{#1}}}
\newtcolorbox{methode}[1][]{popbase,colframe=popgold,before upper={\popboxhead{Méthode}{popgold}{#1}}}
\newtcolorbox{attention}[1][]{popbase,colframe=poppink,before upper={\popboxhead{Attention}{poppink}{#1}}}
\newtcolorbox{experience}[1][]{popbase,colframe=popblue,colback=popblueL,before upper={\popboxhead{Expérience}{popblue}{#1}}}
% « Le savais-tu ? » : carte violette pleine (contexte, culture, Cameroun)
\newtcolorbox{savaistu}[1][]{popbase,colback=poppurple,colframe=popink,
  fontupper=\popsemi\fontsize{10.4}{14.2}\selectfont\color{white},
  before upper={\poppill{Le savais-tu\,?}{white}{poppurple}\ifx\relax#1\relax\else\hspace{6pt}{\popxbold\color{white}#1}\fi\par\vspace{7pt}}}
% Exercice résolu : énoncé en haut, \tcblower puis la solution sur fond crème
\newcounter{popexo}\newcounter{popexr}
\newtcolorbox{exoresolu}[1][]{popbase,colframe=poporange,colbacklower=popcream,
  segmentation style={popink,line width=1.2pt,dashed},
  before upper={\stepcounter{popexr}\popboxhead{Exercice résolu \thepopexr}{poporange}{#1}},
  before lower={\poppill{Solution}{popink}{popcream}\par\vspace{6pt}}}
% Exercice (non résolu en place) — la correction est dans la partie Corrigés
\newtcolorbox{exercice}[1][]{popbase,colframe=popink,
  before upper={\stepcounter{popexo}\popboxhead{Exercice \thepopexo}{popink}{#1}}}
% Corrigé
\newtcolorbox{corrige}[1][]{popbase,colframe=popgreen,colback=popgreenL,
  before upper={\popboxhead{Corrigé}{popgreen}{#1}}}
% Objectifs / prérequis / bilan
\newtcolorbox{objectifs}{popbase,colframe=popink,colback=poporangeL,
  before upper={\popboxhead{Objectifs de la leçon}{popink}{}}}
\newtcolorbox{prerequis}{popbase,colframe=popink,colback=popblueL,
  before upper={\popboxhead{Prérequis}{popblue}{}}}
\newtcolorbox{bilan}{popbase,colframe=popink,colback=white,boxrule=2.8pt,
  before upper={\popboxhead{Fiche bilan}{poporange}{L'essentiel à connaître pour l'examen}}}
% Figure encadrée avec légende
\newtcolorbox{popfigure}[1][]{enhanced,breakable=false,colback=white,colframe=popline,boxrule=1.4pt,arc=8pt,
  left=10pt,right=10pt,top=10pt,bottom=8pt,before skip=10pt,after skip=12pt,halign=center,
  lower separated=false,halign lower=center,
  fontlower=\popsemi\fontsize{9}{11.5}\selectfont\color{popmuted},
  #1}
\newcommand{\legende}[1]{\par\vspace{6pt}{\popsemi\fontsize{9}{11.5}\selectfont\color{popmuted}#1}}

% ============================================================
% Signature OnBuch+
% ============================================================
\newcommand{\poplogoinline}[1]{{\poplogo\fontsize{#1}{#1}\selectfont\color{popink}OnBuch\color{poporange}+}}
\newcommand{\signatureonbuch}{%
  \begin{tcolorbox}[enhanced,colback=popink,colframe=popink,boxrule=2.2pt,arc=10pt,
      drop shadow=poporange,left=14pt,right=14pt,top=10pt,bottom=10pt,before skip=0pt,after skip=0pt]
    \tikz[baseline=-0.7ex]\node[circle,fill=popgreen,draw=popcream,line width=1.3pt,minimum size=22pt,inner sep=0]{\popcheck{white}};%
    \hspace{9pt}\begin{minipage}[c]{0.62\linewidth}
      {\popxbold\fontsize{12}{13}\selectfont\color{popcream}Signé\ \textbullet\ {\poplogo OnBuch\color{poporange}+}}\par\vspace{2pt}
      {\raggedright\popsemi\fontsize{8}{10.5}\selectfont\color{popline}\MakeUppercase{Équipe pédagogique OnBuch+}\\\MakeUppercase{Conforme au programme officiel MINESEC}\par}
    \end{minipage}\hfill
    {\poplogo\fontsize{22}{22}\selectfont\color{popcream}OnBuch\color{poporange}+}
  \end{tcolorbox}%
}

% ============================================================
% Page de garde
% #1 titre · #2 matière · #3 niveau · #4 module · #5 n° leçon · #6 durée ·
% #7 description / objectifs (texte ou itemize)
% ============================================================
\newcommand{\pagegarde}[7]{%
  \thispagestyle{empty}
  \newgeometry{a4paper,margin=1.7cm,top=1.6cm,bottom=1.4cm}%
  \begin{tikzpicture}[remember picture,overlay]
    \fill[poporange!18] ([xshift=-3.2cm,yshift=-3.6cm]current page.north east) circle (4.6cm);
    \fill[poppurple!14] ([xshift=2.4cm,yshift=7.6cm]current page.south west) circle (3.8cm);
  \end{tikzpicture}%
  {\poplogo\fontsize{30}{30}\selectfont\color{popink}OnBuch\color{poporange}+}\hfill
  \raisebox{0.6ex}{\poppill{Cours signé \textbullet\ Premium}{popink}{popcream}}\par
  \vspace{1pt}\popsquiggle{poppurple}\par
  \vspace{8pt}
  \begin{tcolorbox}[enhanced,colback=poporange,colframe=popink,boxrule=2.8pt,arc=13pt,
      drop shadow=popink,left=18pt,right=18pt,top=12pt,bottom=13pt,after skip=10pt]
    \poppill{#2 \textbullet\ #3}{popink}{popcream}\hspace{5pt}\poppill{#4}{white}{popink}\par\vspace{8pt}
    {\popdisplay\fontsize{14}{14}\selectfont\color{popink}LEÇON #5}\par\vspace{4pt}
    {\raggedright\hyphenpenalty=10000\exhyphenpenalty=10000\popdisplay\fontsize{25}{27}\selectfont\color{popcream}\MakeUppercase{#1}\par}
    \vspace{8pt}
    {\popxbold\fontsize{9.5}{9.5}\selectfont\color{popink}\MakeUppercase{Durée conseillée : #6}}
  \end{tcolorbox}
  \begin{tcolorbox}[enhanced,colback=white,colframe=popink,boxrule=2.2pt,arc=10pt,
      drop shadow=popink,left=16pt,right=16pt,top=10pt,bottom=10pt,after skip=8pt]
    \poppill{Dans cette leçon}{poppurple}{white}\par\vspace{5pt}
    \popsemi\fontsize{9.3}{12.6}\selectfont\color{popink}#7
  \end{tcolorbox}
  \noindent\begin{minipage}[t]{0.487\linewidth}
    \begin{tcolorbox}[enhanced,colback=popgreen,colframe=popink,boxrule=2.2pt,arc=10pt,
        drop shadow=popink,left=11pt,right=11pt,top=10pt,bottom=10pt,height=3.1cm,valign=center]
      \tcbox[colback=white,colframe=popink,boxrule=1.3pt,arc=5pt,boxsep=0pt,nobeforeafter,
        left=3pt,right=3pt,top=3pt,bottom=3pt,valign=center]{\qrcode[height=1.9cm]{https://play.google.com/store/apps/details?id=com.luvvix.onbuch&hl=fr}}%
      \hspace{8pt}\begin{minipage}[c]{0.5\linewidth}
        {\popdisplay\fontsize{11}{12.5}\selectfont\color{white}\MakeUppercase{Télécharge OnBuch}}\par\vspace{4pt}
        {\raggedright\popsemi\fontsize{8.4}{11}\selectfont\color{white}Gratuit \textbullet\ Google Play\\Tuteur IA, annales, concours, résultats\par}
      \end{minipage}
    \end{tcolorbox}
  \end{minipage}\hfill
  \begin{minipage}[t]{0.487\linewidth}
    \begin{tcolorbox}[enhanced,colback=poppurple,colframe=popink,boxrule=2.2pt,arc=10pt,
        drop shadow=popink,left=11pt,right=11pt,top=10pt,bottom=10pt,height=3.1cm,valign=center]
      \tikz[baseline=-0.7ex]\node[circle,fill=white,draw=popink,line width=1.3pt,minimum size=1.6cm,inner sep=0]{\popdisplay\fontsize{24}{24}\selectfont\color{poppurple}+};%
      \hspace{8pt}\begin{minipage}[c]{0.6\linewidth}
        {\popdisplay\fontsize{11}{12.5}\selectfont\color{white}\MakeUppercase{Passe à OnBuch+}}\par\vspace{4pt}
        {\raggedright\popsemi\fontsize{8.4}{11}\selectfont\color{white}Tous les cours signés, Tuteur IA illimité, fiches PDF — onglet OnBuch+ de l'app\par}
      \end{minipage}
    \end{tcolorbox}
  \end{minipage}
  \vspace{8pt}
  \popcontenu
  \vfill
  \signatureonbuch
  \restoregeometry
  \newpage
}

% Rangée « Ce cours contient » (page de garde)
\newcommand{\poptile}[3]{% #1 couleur #2 gros texte #3 légende
  \begin{tcolorbox}[enhanced,colback=#1,colframe=popink,boxrule=2pt,arc=9pt,shadow={2.5pt}{-2.5pt}{0pt}{popink},
      left=6pt,right=6pt,top=9pt,bottom=9pt,halign=center,valign=center,height=1.75cm,width=\linewidth,nobeforeafter]
    {\popdisplay\fontsize{12.5}{13}\selectfont\color{white}\MakeUppercase{#2}}\par\vspace{3pt}
    {\centering\popsemi\fontsize{7.8}{9.5}\selectfont\color{white}#3\par}
  \end{tcolorbox}}
\newcommand{\popcontenu}{%
  \par\noindent{\popxboldls\fontsize{7.5}{7.5}\selectfont\color{popmuted}CE COURS SIGNÉ CONTIENT}\par\vspace{7pt}
  \noindent\begin{minipage}[t]{0.235\linewidth}\poptile{popblue}{Cours}{complet, illustré, pas à pas}\end{minipage}\hfill
  \begin{minipage}[t]{0.235\linewidth}\poptile{poporange}{Méthodes}{et exercices résolus}\end{minipage}\hfill
  \begin{minipage}[t]{0.235\linewidth}\poptile{poppink}{Exercices}{type examen + corrigés}\end{minipage}\hfill
  \begin{minipage}[t]{0.235\linewidth}\poptile{popgreen}{Fiche bilan}{l'essentiel à retenir}\end{minipage}\par}

% Sommaire
\newtcolorbox{sommaire}{enhanced,colback=white,colframe=popink,boxrule=2.2pt,arc=10pt,
  drop shadow=popink,left=18pt,right=18pt,top=13pt,bottom=13pt,before skip=6pt,after skip=20pt,
  before upper={\poppill{Sommaire}{poppurple}{white}\par\vspace{9pt}\popsemi\fontsize{10.6}{14.5}\selectfont\color{popink}}}

% Signature de fin de cours — poussée tout en bas de la dernière page
\newcommand{\findecours}{%
  \par\vfill\nopagebreak
  \begin{center}\popsquiggle{poporange}\end{center}\vspace{4pt}
  \signatureonbuch
}
\AtBeginDocument{\def\llbracket{\mathopen{[\mkern-3.5mu[}}\def\rrbracket{\mathclose{]\mkern-3.5mu]}}} % intervalles d'entiers (glyphe ⟦ absent de la police)
% Glyphes absents de Fira Math : redéfinis à partir de glyphes disponibles
\AtBeginDocument{%
  \def\setminus{\mathbin{\backslash}}\let\smallsetminus\setminus
  \def\vdots{\mathord{\vbox{\baselineskip3.2pt\lineskiplimit0pt\kern4pt\hbox{.}\hbox{.}\hbox{.}}}}%
  \def\ddots{\mathinner{\mkern1mu\raise6.5pt\hbox{.}\mkern2mu\raise3.6pt\hbox{.}\mkern2mu\raise0.7pt\hbox{.}\mkern1mu}}%
  \def\triangle{\mathord{\scalebox{0.9}{$\symup{\Delta}$}}}%
  \def\nabla{\mathord{\raisebox{\depth}{\scalebox{1}[-1]{$\symup{\Delta}$}}}}%
  \def\checkmark{\popcheck{popdark}}%
  \def\star{\mathbin{\ast}}\def\bigstar{\mathord{\ast}}%
  \def\diamond{\mathbin{\scalebox{0.6}{$\Diamond$}}}%
  \def\bigcup{\mathop{\mathchoice{\raisebox{-0.3em}{\scalebox{1.7}{$\cup$}}}{\scalebox{1.25}{$\cup$}}{\cup}{\cup}}\displaylimits}%
  \def\bigcap{\mathop{\mathchoice{\raisebox{-0.3em}{\scalebox{1.7}{$\cap$}}}{\scalebox{1.25}{$\cap$}}{\cap}{\cap}}\displaylimits}%
  \def\lesseqqgtr{\mathrel{\lessgtr}}\def\bigoplus{\mathop{\oplus}\displaylimits}%
}
\providecommand{\ssi}{si et seulement si} % abréviation fréquente
\AtBeginDocument{\providecommand{\wideparen}{\overparen}} % arc de cercle

%% FIGFIX-BEGIN v3 — correctifs globaux des figures (docs/onbuch-generation/tools/figfix)
\usepackage{adjustbox}\usepackage{etoolbox}
% toute figure TikZ/pgfplots plus large que la ligne est réduite à la largeur disponible
\BeforeBeginEnvironment{tikzpicture}{\begin{adjustbox}{max width=\linewidth}}
\AfterEndEnvironment{tikzpicture}{\end{adjustbox}}
% légendes pgfplots lisibles par-dessus les courbes
\pgfplotsset{every axis/.append style={legend style={fill=white,fill opacity=0.92,text opacity=1,draw=black!15,inner sep=3pt}}}
% étiquettes d'arêtes (syntaxe edge["..."]) avec fond blanc
\tikzset{every edge quotes/.append style={fill=white,inner sep=1.2pt}}
% bulles de cartes mentales : texte à la ligne dans la bulle, bulle un peu plus grande
\tikzset{every concept/.append style={text width=2.0cm,align=center,minimum size=2.3cm,inner sep=2pt}}
%% FIGFIX-END
\begin{document}

\pagegarde{Les mécanismes de la mondialisation}{Géographie}{Tle Tech}{Géographie – classe de Terminale technique}{N04}{2 séances (fiche de progression harmonisée)}{Cette leçon explique les facteurs qui ont rendu possible la mondialisation (révolutions des transports, des télécommunications, libéralisation des échanges) et analyse son fonctionnement concret : flux de biens, de capitaux, d'informations et d'hommes organisés par les firmes transnationales et les grandes places mondiales. Elle s'appuie sur des cartes, des données chiffrées et des exemples camerounais pour construire des croquis et des diagrammes.
\vspace{4pt}
\begin{multicols}{2}\small
\begin{itemize}
\item À la fin de cette leçon, je sais définir la mondialisation et identifier ses principaux facteurs
\item À la fin de cette leçon, je sais expliquer le rôle des transports et des télécommunications dans la mise en relation des territoires
\item À la fin de cette leçon, je sais décrire les principaux flux mondiaux (marchandises, capitaux, informations, migrations) à partir d'une carte
\item À la fin de cette leçon, je sais expliquer le rôle des firmes transnationales et des États dans le fonctionnement de la mondialisation
\item À la fin de cette leçon, je sais construire un croquis ou un diagramme simple représentant un flux mondial
\item À la fin de cette leçon, je sais appliquer les notions de la leçon à des situations concrètes de la vie courante au Cameroun
\end{itemize}\end{multicols}}

\begin{sommaire}
\begin{multicols}{2}
\begin{enumerate}
\item La mondialisation : définition et mise en évidence
\item Les facteurs techniques de la mondialisation
\item Les facteurs économiques et politiques de la mondialisation
\item Le fonctionnement de la mondialisation : les flux
\item Les acteurs du fonctionnement : FTN, métropoles et États
\item TP guidé : croquis et diagrammes des mécanismes de la mondialisation
\item Activité d'intégration
\item Méthodes et exercices résolus
\item Exercices
\item Corrigés des exercices
\item Fiche bilan
\end{enumerate}
\end{multicols}
\end{sommaire}

\begin{objectifs}
\begin{itemize}
\item À la fin de cette leçon, je sais définir la mondialisation et identifier ses principaux facteurs
\item À la fin de cette leçon, je sais expliquer le rôle des transports et des télécommunications dans la mise en relation des territoires
\item À la fin de cette leçon, je sais décrire les principaux flux mondiaux (marchandises, capitaux, informations, migrations) à partir d'une carte
\item À la fin de cette leçon, je sais expliquer le rôle des firmes transnationales et des États dans le fonctionnement de la mondialisation
\item À la fin de cette leçon, je sais construire un croquis ou un diagramme simple représentant un flux mondial
\item À la fin de cette leçon, je sais appliquer les notions de la leçon à des situations concrètes de la vie courante au Cameroun
\item À la fin de cette leçon, je sais rédiger et justifier une démarche, une réponse ou une conclusion à partir de documents
\item À la fin de cette leçon, je sais montrer que la mondialisation produit une interdépendance inégale entre les territoires
\end{itemize}
\end{objectifs}

\begin{prerequis}
\begin{itemize}
\item La notion de mondialisation vue en classe de Première (acteurs et manifestations)
\item Les grandes puissances et les pôles de l'économie mondiale (Triade)
\item La lecture d'une carte et la construction d'un croquis géographique
\item Les notions de flux, de réseau et d'échange vues en Géographie de 4ème/3ème
\item La localisation des principaux continents, océans et grandes villes mondiales
\end{itemize}
\end{prerequis}

\newpage

\coursec{La mondialisation : définition et mise en évidence}

\courssub{Situation d'accroche et problématique}

À Douala, un jeune commerçant du quartier Deïdo commande un lot de téléphones chinois sur la plateforme Alibaba, paye par virement bancaire en quelques secondes et reçoit sa marchandise au port autonome de Douala trois semaines plus tard. Le même jour, sa sœur regarde une série américaine sur Netflix depuis son smartphone, tandis que leur oncle, émigré aux États-Unis, envoie de l'argent à la famille par Western Union. À l'autre bout de la ville, un chauffeur de taxi écoute à la radio les cours du pétrole à New York, car il sait que ces cours influencent le prix du carburant à la pompe à Bonabéri.

Cette scène, banale en apparence, est extraordinaire : en une seule journée, une famille camerounaise ordinaire est connectée à la Chine, aux États-Unis et aux marchés financiers mondiaux. Comment des produits, des capitaux, des informations et des hommes circulent-ils ainsi à l'échelle de la planète ? Quels mécanismes rendent possible cette interconnexion du monde, et pourquoi le Cameroun en bénéficie-t-il inégalement ? Telle est la problématique de cette leçon.

\courssub{Définition de la mondialisation}

\textbf{Intuition.} Observez vos vêtements, votre téléphone, votre nourriture : le tissu vient peut-être du Nigeria, le téléphone de Chine, le riz de Thaïlande. Aucun pays ne vit plus seul, replié sur lui-même. Le monde fonctionne comme un immense réseau où chaque territoire est relié aux autres.

\begin{definition}[La mondialisation]
La \cle{mondialisation} est le processus d'\cle{interconnexion} croissante des économies et des sociétés à l'échelle planétaire. Elle se traduit par l'intensification des \cle{flux} (de marchandises, de capitaux, d'informations, de personnes) entre les différentes parties du monde, et par l'interdépendance grandissante des territoires.
\end{definition}

Analysons cette définition mot à mot :
\begin{itemize}
\item C'est un \textbf{processus} : la mondialisation n'est pas un état figé, mais un mouvement historique qui s'accélère depuis plusieurs décennies.
\item Elle concerne les \textbf{économies} (échanges, investissements, entreprises) mais aussi les \textbf{sociétés} (cultures, modes de vie, migrations).
\item Elle s'exerce à l'échelle \textbf{planétaire} : presque aucun territoire n'échappe à ces connexions, même si tous n'y participent pas de la même façon.
\end{itemize}

\begin{attention}[Mondialisation ou occidentalisation ?]
Erreur fréquente à l'examen : confondre la \cle{mondialisation} et l'\cle{occidentalisation}. La mondialisation est un processus global d'interconnexion qui touche tous les domaines (économie, culture, politique). L'occidentalisation (adoption des modes de vie européens et nord-américains : langue anglaise, fast-food, habillement) n'est qu'\textbf{une seule dimension culturelle} de la mondialisation, et encore contestée : la culture chinoise, indienne ou nigériane (Nollywood) rayonne aussi dans le monde. Ne jamais écrire que la mondialisation se résume à l'américanisation du monde.
\end{attention}

\courssub{Les manifestations de la mondialisation dans la vie quotidienne camerounaise}

La mondialisation n'est pas une notion abstraite : elle est observable chaque jour au Cameroun.

\textbf{Les produits importés.} Les marchés de Douala, Yaoundé ou Bafoussam regorgent de marchandises venues du monde entier : téléphones et textiles chinois, riz thaïlandais, huile de palme malaisienne, voitures d'occasion européennes (les « au revoir Calais »), médicaments indiens. En sens inverse, le Cameroun exporte cacao, café, coton, pétrole et bois vers l'Union européenne et la Chine.

\textbf{La téléphonie mobile et internet.} Le téléphone portable est devenu l'outil quotidien de connexion au monde : appels internationaux, réseaux sociaux (WhatsApp, Facebook, TikTok), information en temps réel. Les câbles sous-marins à fibre optique (SAT-3, WACS, ACE) qui relient directement le Cameroun (Douala, Kribi) au réseau mondial ont fait entrer le pays dans l'ère numérique.

\textbf{Les transferts d'argent de la diaspora.} Les Camerounais vivant à l'étranger (Europe, Amérique du Nord, Afrique du Sud) envoient chaque année des centaines de milliards de FCFA à leurs familles par Western Union, MoneyGram ou les applications mobiles. Ces transferts financent la construction de maisons, la scolarité des enfants et les petits commerces : c'est un flux financier majeur de la mondialisation.

\begin{savaistu}[Le téléphone mobile, outil de connexion au monde]
Plus de 80 \% des Camerounais urbains possèdent un téléphone mobile. Grâce au \emph{mobile money} (Orange Money, MTN Mobile Money), même les habitants des quartiers non bancarisés paient, épargnent et reçoivent de l'argent de la diaspora sans compte bancaire. Le téléphone mobile est ainsi devenu le principal outil quotidien de connexion du Camerounais ordinaire à la mondialisation.
\end{savaistu}

\courssub{Les trois dimensions de la mondialisation}

La mondialisation n'est pas seulement économique. Elle possède trois dimensions complémentaires et interconnectées.

\begin{itemize}
\item \textbf{La dimension économique} : libre-échange, multiplication des échanges de marchandises et de services, mobilité des capitaux, implantation mondiale des firmes transnationales (FTN). C'est la dimension la plus visible.
\item \textbf{La dimension culturelle} : diffusion planétaire des langues (anglais), des musiques (afrobeat, hip-hop), des séries, des sports (la Coupe du monde de football suivie sur tous les continents), mais aussi échanges et métissages culturels.
\item \textbf{La dimension politique} : multiplication des organisations internationales (ONU, OMC, Union africaine, CEMAC), adoption de règles communes (droits de l'homme, accords sur le climat), diplomatie mondiale.
\end{itemize}

\begin{popfigure}
\begin{center}
\begin{tikzpicture}
\draw[fill=popblue!35,draw=popblue,thick] (-1.3,0) circle (1.7);
\draw[fill=poporange!35,draw=poporange,thick] (1.3,0) circle (1.7);
\draw[fill=popgreen!35,draw=popgreen,thick] (0,1.7) circle (1.7);
\node[font=\bfseries] at (-2.1,-0.9) {Économique};
\node[font=\bfseries] at (2.1,-0.9) {Culturelle};
\node[font=\bfseries] at (0,2.9) {Politique};
\node[font=\small\bfseries,text=popdark,align=center] at (0,0.7) {Mondialisation};
\node[font=\scriptsize,align=center] at (-1.9,0.35) {flux de biens,\\capitaux, FTN};
\node[font=\scriptsize,align=center] at (1.9,0.35) {langues,\\séries, sports};
\node[font=\scriptsize,align=center] at (0,1.35) {ONU, OMC,\\accords};
\end{tikzpicture}
\end{center}
\legende{Schéma 1 : les trois dimensions interconnectées de la mondialisation. Chaque dimension recouvre partiellement les deux autres : par exemple, l'OMC est une organisation politique qui régit les échanges économiques.}
\end{popfigure}

\courssub{L'accélération de la mondialisation depuis les années 1980}

Si les échanges internationaux existent depuis l'Antiquité (routes de la soie, commerce transatlantique), la mondialisation connaît une \textbf{accélération spectaculaire depuis les années 1980}. Deux indicateurs le prouvent : la valeur du commerce mondial de marchandises et les investissements directs à l'étranger (IDE).

\begin{exemplebox}[L'explosion du commerce mondial en chiffres]
La valeur du commerce mondial de marchandises est passée d'environ \textbf{2 000 milliards de dollars en 1980} à plus de \textbf{25 000 milliards de dollars en 2022}, soit une multiplication par plus de 12 en un peu plus de quarante ans. Dans le même temps, les stocks d'IDE (capitaux investis par les entreprises à l'étranger) sont passés de moins de 1 000 milliards de dollars à plus de 40 000 milliards. Ces chiffres illustrent l'intensification sans précédent des flux économiques planétaires.
\end{exemplebox}

\begin{popfigure}
\begin{center}
\begin{tikzpicture}
\begin{axis}[
width=0.95\linewidth,height=6.2cm,
xlabel={Ann\'ee},ylabel={Milliers de milliards de dollars},
xmin=1978,xmax=2025,ymin=0,ymax=27,
xtick={1980,1990,2000,2010,2020},
ytick={0,5,10,15,20,25},
grid=major,grid style={popline!40},
legend pos=north west,
legend style={font=\scriptsize}
]
\addplot[popcurve,mark=*,mark size=1.5pt] coordinates {
(1980,2.0) (1990,3.5) (2000,6.4) (2008,16.1) (2010,15.2) (2015,16.5) (2020,17.6) (2022,25.0) (2023,23.8)
};
\legend{Commerce mondial de marchandises (exportations)}
\end{axis}
\end{tikzpicture}
\end{center}
\legende{Graphique 1 : évolution de la valeur du commerce mondial de marchandises de 1980 à 2023, en milliers de milliards de dollars (ordres de grandeur, source : OMC). On observe une croissance forte, avec des reculs ponctuels lors des crises (2009, 2020).}
\end{popfigure}

La lecture du graphique appelle deux remarques. D'une part, la courbe monte fortement et régulièrement : la mondialisation est bien un processus d'intensification. D'autre part, la courbe fléchit lors des crises (2009, 2020) : les économies sont tellement liées qu'une crise née dans un pays freine les échanges du monde entier.

\courssub{L'interdépendance des territoires}

\begin{definition}[L'interdépendance]
L'\cle{interdépendance} désigne la situation dans laquelle des territoires dépendent les uns des autres : un événement local (crise financière, épidémie, conflit, catastrophe) produit des effets à l'échelle mondiale, et chaque territoire dépend des autres pour ses approvisionnements, ses débouchés et son financement.
\end{definition}

Trois exemples récents démontrent cette interdépendance :

\begin{itemize}
\item \textbf{La crise financière de 2008.} Née de la faillite de banques américaines (subprimes), elle s'est propagée en quelques semaines aux bourses européennes et asiatiques, puis a ralenti l'économie mondiale entière, réduisant la demande de pétrole et de matières premières exportées par le Cameroun.
\item \textbf{La pandémie de Covid-19 (2020).} Un virus apparu en Chine a paralysé la planète : fermeture des frontières, rupture des chaînes d'approvisionnement, pénuries de médicaments et de composants électroniques. Le commerce mondial a chuté, comme le montre le graphique 1.
\item \textbf{La guerre en Ukraine (2022).} La Russie et l'Ukraine fournissent environ un quart des exportations mondiales de blé. Le conflit a bloqué les ports de la mer Noire : le prix mondial du blé a flambé, et au Cameroun, très dépendant du blé importé, le prix de la baguette de pain est passé de 125 à 150 FCFA en 2022. Un conflit à 6 000 km a ainsi touché le petit-déjeuner des Camerounais.
\end{itemize}

\begin{exoresolu}[Calculer et interpréter une évolution]
Énoncé. La valeur du commerce mondial de marchandises est passée d'environ 2 000 milliards de dollars en 1980 à 25 000 milliards de dollars en 2022. 1) Calculez le taux de variation entre 1980 et 2022. 2) Que révèle ce résultat ?
\tcblower
Solution détaillée.
1) On applique la formule du taux de variation :
\[ t = \frac{V_{2022}-V_{1980}}{V_{1980}}\times 100 = \frac{25\,000-2\,000}{2\,000}\times 100 = \frac{23\,000}{2\,000}\times 100 = 1\,150\ \% \]
2) Le commerce mondial de marchandises a augmenté de 1 150 \% entre 1980 et 2022 : sa valeur a été multipliée par 12,5. Cette croissance spectaculaire prouve l'accélération de la mondialisation depuis les années 1980 : les économies échangent de plus en plus et deviennent de plus en plus interdépendantes.
\end{exoresolu}

\begin{aretenir}[L'essentiel de la section]
\begin{itemize}
\item La mondialisation est le processus d'interconnexion croissante des économies et des sociétés à l'échelle planétaire.
\item Elle se manifeste au Cameroun par les produits importés, la téléphonie mobile, internet et les transferts d'argent de la diaspora.
\item Elle possède trois dimensions : économique, culturelle et politique.
\item Elle s'accélère depuis les années 1980 : le commerce mondial est passé de 2 000 à plus de 25 000 milliards de dollars (1980-2022).
\item Elle crée une interdépendance des territoires : un événement local (crise de 2008, Covid-19, guerre en Ukraine) a des effets mondiaux, jusque sur le prix du pain au Cameroun.
\end{itemize}
\end{aretenir}

\begin{exercice}[Pour s'entraîner]
1) Définissez la mondialisation en une phrase complète. 2) Citez deux manifestations de la mondialisation observables dans votre quartier ou votre ville. 3) Expliquez, avec un exemple, pourquoi on dit que les territoires sont interdépendants. 4) Pourquoi est-il faux d'affirmer que la mondialisation est synonyme d'occidentalisation ?
\end{exercice}

\coursec{Les facteurs techniques de la mondialisation}

Dans la section précédente, nous avons défini la mondialisation comme un processus d'intensification des échanges et d'interdépendance des espaces à l'échelle planétaire, et nous avons montré qu'elle se manifeste par des flux croissants de marchandises, de capitaux, d'informations et de personnes. Mais une question demeure : comment ces flux sont-ils devenus possibles à une telle échelle et à une telle vitesse ? La réponse tient en deux mots : les \cle{innovations techniques}. Sans la révolution des transports et celle des télécommunications, la mondialisation serait restée un rêve. C'est ce que nous allons étudier maintenant.

\courssub{La révolution des transports maritimes}

\textbf{Une intuition simple.} Pour vendre un tee-shirt fabriqué au Bangladesh à 5 000 FCFA à Douala, il faut que le transport de ce tee-shirt coûte quelques francs seulement. Cela n'est possible que si l'on transporte des milliers de tonnes de marchandises d'un coup, dans des boîtes standardisées, sur des navires gigantesques. C'est exactement ce qu'a permis la \cle{conteneurisation}.

\begin{definition}[La conteneurisation]
La \cle{conteneurisation} est un mode de transport des marchandises dans des boîtes métalliques normalisées appelées \cle{conteneurs}. L'unité de mesure est l'\cle{EVP} (Équivalent Vingt Pieds), c'est-à-dire un conteneur de 20 pieds de long (environ 6 mètres). Le conteneur passe du camion au train, puis du train au navire, sans que la marchandise soit déballée : c'est le \cle{transport multimodal}. Inventée par l'Américain Malcom McLean en 1956, elle a réduit de façon spectaculaire les coûts et les délais de manutention dans les ports.
\end{definition}

\textbf{Le mécanisme.} Avant 1956, le chargement d'un navire exigeait des centaines de dockers manipulant sacs, caisses et fûts un par un : un navire restait des semaines à quai. Avec le conteneur, une grue portique charge des centaines de boîtes par heure. Conséquences :

\begin{itemize}
\item la \emph{baisse des coûts de fret} : le coût du transport maritime a été divisé par plus de dix depuis 1950, si bien que le transport représente souvent moins de 2 \% du prix final d'un produit manufacturé ;
\item le \emph{gigantisme des navires} : pour rentabiliser les trajets, les armateurs construisent des porte-conteneurs de plus en plus grands ;
\item la \emph{spécialisation des ports} : les grands ports (Shanghai, Singapour, Rotterdam) deviennent des \cle{hubs} (plaques tournantes) où les flux se concentrent.
\end{itemize}

\begin{exemplebox}[Un géant des mers]
Les plus grands porte-conteneurs du monde mesurent près de 400 mètres de long (soit quatre terrains de football bout à bout) et transportent plus de \textbf{24 000 EVP}. Un seul de ces navires peut emporter l'équivalent du contenu de plusieurs milliers de camions. À cette échelle, le transport d'un téléviseur de la Chine vers l'Europe coûte moins de 10 dollars.
\end{exemplebox}

\textbf{Les routes maritimes et les points de passage obligés.} Le transport maritime assure environ \textbf{80 \% du commerce mondial en volume}. Les flux se concentrent sur quelques grandes routes qui relient les trois pôles de la Triade (Amérique du Nord, Europe, Asie de l'Est). Ces routes passent par des \cle{détroits} et des \cle{canaux} stratégiques, véritables goulots d'étranglement du commerce mondial.

\begin{popfigure}
\begin{tikzpicture}[scale=1.0, font=\small]
% Océan
\fill[popblueL] (0,0) rectangle (14,6.2);
% Continents très simplifiés
\fill[popgreenL, draw=popmuted] (0.2,2.6) -- (1.6,5.6) -- (3.0,5.4) -- (3.4,3.6) -- (2.4,2.4) -- (1.2,2.2) -- cycle; % Amérique N
\fill[popgreenL, draw=popmuted] (2.6,0.4) -- (3.6,2.4) -- (4.4,2.2) -- (4.0,0.2) -- cycle; % Amérique S
\fill[popgreenL, draw=popmuted] (6.6,3.4) -- (7.4,5.6) -- (9.4,5.6) -- (9.6,4.2) -- (8.6,3.6) -- (7.4,3.2) -- cycle; % Europe
\fill[popgreenL, draw=popmuted] (6.8,0.6) -- (7.4,3.2) -- (8.8,3.4) -- (9.6,2.2) -- (9.0,0.4) -- cycle; % Afrique
\fill[popgreenL, draw=popmuted] (9.8,2.6) -- (9.8,4.6) -- (11.6,5.6) -- (13.4,5.2) -- (13.2,3.4) -- (12.0,2.6) -- (11.0,2.2) -- cycle; % Asie
\fill[popgreenL, draw=popmuted] (11.8,0.4) -- (12.8,1.2) -- (13.6,0.8) -- (13.2,0.2) -- cycle; % Océanie
% Routes maritimes
\draw[line width=2pt, poporange] (3.2,4.4) -- (6.8,4.6); % Atlantique N
\draw[line width=2pt, poporange] (4.2,1.6) -- (7.0,1.8); % Atlantique S
\draw[line width=2pt, poporange] (8.8,3.4) -- (9.9,3.1) -- (10.6,2.6) -- (11.4,2.3); % Suez -> océan Indien
\draw[line width=2pt, poporange] (11.4,2.3) -- (12.3,2.9); % Malacca
\draw[line width=2pt, poporange] (12.3,2.9) -- (12.6,4.2); % vers Chine/Japon
\draw[line width=2pt, poporange] (3.0,3.0) -- (3.4,3.3); % Panama
% Détroits (points roses)
\fill[poppink] (3.2,3.15) circle (0.09);
\fill[poppink] (7.6,4.55) circle (0.09);
\fill[poppink] (9.9,3.1) circle (0.09);
\fill[poppink] (12.3,2.9) circle (0.09);
% Étiquettes
\node[above, popdark] at (3.2,3.3) {\footnotesize Panama};
\node[above, popdark] at (7.6,4.7) {\footnotesize Gibraltar};
\node[above right, popdark] at (9.9,3.15) {\footnotesize Suez};
\node[below, popdark] at (12.3,2.75) {\footnotesize Malacca};
\node[popdark] at (5.0,4.8) {\footnotesize Route Atlantique Nord};
\node[popdark] at (10.6,1.9) {\footnotesize Route de l'océan Indien};
% Nord
\draw[->, thick, popdark] (13.6,5.4) -- (13.6,6.0);
\node[above, popdark] at (13.6,6.0) {\footnotesize N};
\end{tikzpicture}
\legende{Fig. 1 — Les grandes routes maritimes mondiales (traits orange épais) et les points de passage obligés (points roses) : Panama, Gibraltar, Suez, Malacca. Croquis simplifié, sans échelle précise. Le détroit de Malacca voit passer environ 30 \% du commerce maritime mondial.}
\end{popfigure}

\begin{methode}[Lire une carte de flux maritimes]
Pour analyser correctement une carte de flux maritimes, procéder en trois étapes :
\begin{enumerate}
\item \textbf{Identifier les routes} : repérer les lignes de flux et comparer leur épaisseur (plus le trait est épais, plus le trafic est important) ; noter les espaces reliés (généralement les pôles de la Triade).
\item \textbf{Repérer les nœuds} : localiser les grands ports (hubs) où les flux convergent et se redistribuent, comme Singapour, Shanghai ou Rotterdam.
\item \textbf{Localiser les points de passage obligés} : détroits (Malacca, Gibraltar, Ormuz) et canaux (Suez, Panama) ; expliquer leur importance stratégique : leur fermeture obligerait les navires à des détours de plusieurs semaines.
\end{enumerate}
Conclure toujours en reliant la carte à l'idée de hiérarchie des espaces : toutes les mers ne se valent pas.
\end{methode}

\courssub{Les transports aériens : la vitesse pour les hommes et les produits de valeur}

Si le navire transporte le volume, l'avion transporte la \emph{valeur} et les \emph{personnes}. Le transport aérien est beaucoup plus coûteux au kilogramme, mais il est incomparablement plus rapide : un avion relie Douala à Paris en moins de 7 heures, là où un navire met plus de deux semaines.

\begin{itemize}
\item \textbf{Le transport des passagers} : plus de 4 milliards de passagers par an dans le monde avant la crise de 2020. L'avion est l'outil des migrations internationales, du tourisme de masse et des déplacements d'affaires des cadres des firmes transnationales.
\item \textbf{Le fret aérien} : il ne représente qu'environ 1 \% du commerce mondial en volume, mais près de \textbf{35 \% en valeur}. Il transporte les \cle{produits à forte valeur ajoutée} (téléphones, composants électroniques, médicaments, pièces détachées) et les produits périssables (fleurs coupées du Kenya, fruits et légumes, poissons).
\end{itemize}

Les grands aéroports (Dubaï, Londres-Heathrow, Paris-Charles-de-Gaulle, Addis-Abeba) fonctionnent aussi comme des hubs : ils concentrent les flux de passagers et de marchandises et renforcent la puissance des métropoles qui les accueillent.

\courssub{La révolution des télécommunications}

\textbf{Une intuition simple.} En 1950, téléphoner du Cameroun en France relevait de l'exploit. Aujourd'hui, un élève de Yaoundé envoie en une seconde une photo à un correspondant de Tokyo. L'information circule désormais \emph{instantanément} et presque \emph{gratuitement}. C'est la troisième révolution technique de la mondialisation.

\begin{itemize}
\item \textbf{Les câbles sous-marins de fibre optique} : posés au fond des océans, ils relient les continents et transportent l'essentiel du trafic mondial de données sous forme de signaux lumineux.
\item \textbf{Les satellites} : ils assurent la télédiffusion, le positionnement (GPS) et les télécommunications vers les zones non câblées.
\item \textbf{Internet} : réseau mondial interconnecté, il permet les échanges de données, le commerce électronique, la visioconférence et les transactions financières en temps réel.
\item \textbf{La téléphonie mobile} : avec plus de 5 milliards d'utilisateurs, elle a permis aux pays en développement de « sauter » l'étape du téléphone fixe. Au Cameroun, le mobile money illustre cette irruption du numérique dans la vie quotidienne.
\end{itemize}

\begin{savaistu}[Des autoroutes invisibles au fond des océans]
Plus de \textbf{95 \% du trafic internet mondial} transite par environ 500 câbles sous-marins de fibre optique posés au fond des océans, et non par satellite comme on le croit souvent. Un seul câble moderne peut transporter des dizaines de téraoctets de données par seconde. Ces câbles, longs parfois de plus de 10 000 km, sont vulnérables : une ancre de navire ou un séisme sous-marin peut priver un pays entier de connexion pendant plusieurs jours.
\end{savaistu}

\begin{attention}[Internet n'est pas « immatériel »]
Erreur fréquente dans les copies : présenter internet comme un « nuage » immatériel, présent partout et nulle part. En réalité, internet repose sur des \textbf{infrastructures physiques localisées} : câbles sous-marins, centres de données (data centers), serveurs, antennes relais. Ces infrastructures sont concentrées dans les pays riches et les grandes métropoles. Celui qui contrôle les câbles et les serveurs contrôle une partie de l'économie mondiale : c'est un enjeu de \textbf{puissance}.
\end{attention}

\courssub{La compression de l'espace-temps : le monde « rétrécit »}

La conséquence majeure de ces révolutions techniques est ce que les géographes appellent la \cle{compression de l'espace-temps} : les distances n'ont pas changé, mais le \emph{temps} nécessaire pour les franchir s'est effondré. Le monde « rétrécit ».

\begin{center}
\begin{tabularx}{\linewidth}{|l|X|X|}
\hline
\rowcolor{popblueL} \textbf{Liaison} & \textbf{Vers 1850} & \textbf{Aujourd'hui} \\
\hline
Londres -- New York (marchandises) & 3 à 4 semaines (voilier) & 7 jours (porte-conteneurs) ou 7 heures (avion) \\
\hline
Paris -- Douala (passagers) & plusieurs semaines (paquebot) & environ 7 heures (avion) \\
\hline
Envoi d'un message intercontinental & plusieurs semaines (courrier) & instantané (internet) \\
\hline
Coût du fret maritime (indice) & 100 (en 1950) & inférieur à 10 \\
\hline
\end{tabularx}
\end{center}

\begin{popfigure}
\begin{tikzpicture}
\begin{axis}[
  ybar, bar width=0.9cm,
  width=0.85\linewidth, height=6.5cm,
  ymin=0, ymax=115,
  ylabel={Indice de coût (base 100 en 1950)},
  symbolic x coords={1950,1970,1990,2010,2020},
  xtick=data,
  nodes near coords,
  every node near coord/.append style={font=\small},
  legend style={at={(0.5,-0.18)}, anchor=north, legend columns=2},
  ymajorgrids=true, grid style={popline!40},
]
\addplot[fill=popblue, draw=popdark] coordinates {(1950,100) (1970,60) (1990,30) (2010,12) (2020,8)};
\addplot[fill=poporange, draw=popdark] coordinates {(1950,100) (1970,45) (1990,15) (2010,3) (2020,1)};
\legend{Fret maritime, Communications internationales}
\end{axis}
\end{tikzpicture}
\legende{Fig. 2 — Baisse des coûts du transport maritime et des communications internationales depuis 1950 (indices, base 100 en 1950, valeurs indicatives). L'effondrement du coût des communications est encore plus spectaculaire que celui du fret : c'est le moteur de la mondialisation financière et numérique.}
\end{popfigure}

\begin{exoresolu}[Calculer et interpréter une baisse de coût]
Énoncé : l'indice du coût du fret maritime passe de 100 en 1950 à 8 en 2020. 1) Calculez le taux de variation du coût entre 1950 et 2020. 2) Expliquez en quelques lignes une conséquence de cette baisse pour la mondialisation.
\tcblower
Solution détaillée :
1) Le taux de variation se calcule par :
\[ T = \frac{V_{2020} - V_{1950}}{V_{1950}} \times 100 = \frac{8 - 100}{100} \times 100 = -92 \% \]
Le coût du fret maritime a donc diminué de \textbf{92 \%} entre 1950 et 2020.
2) Conséquence : le transport ne représente plus qu'une part infime du prix des marchandises. Il devient donc rentable de produire à l'autre bout du monde (par exemple au Bangladesh ou en Chine) pour vendre en Europe ou en Afrique. Cette baisse a permis la \emph{délocalisation} des usines et la fragmentation mondiale des chaînes de production : un même téléphone est conçu aux États-Unis, assemblé en Chine avec des composants venus de dix pays, puis vendu au Cameroun.
\end{exoresolu}

\courssub{Le raccordement du Cameroun aux réseaux mondiaux}

Le Cameroun n'est pas à l'écart de ces révolutions techniques : il cherche au contraire à s'insérer dans les réseaux mondiaux.

\textbf{Les câbles sous-marins.} Le pays est raccordé à plusieurs câbles de fibre optique qui débarquent sur son littoral :

\begin{itemize}
\item \textbf{SAT-3} : câble reliant l'Europe à l'Afrique du Sud le long de la côte atlantique africaine, avec un point d'atterrissement à \textbf{Douala} (2002) ;
\item \textbf{ACE} (\emph{Africa Coast to Europe}) : câble entré en service en 2012, débarquant également à \textbf{Douala} ;
\item \textbf{SAIL} : câble reliant le Cameroun au Brésil, débarquant à \textbf{Kribi} (2018), première liaison directe entre l'Afrique centrale et l'Amérique du Sud.
\end{itemize}

Ces câbles sont prolongés par une dorsale nationale de fibre optique qui relie les grandes villes (Douala, Yaoundé, Bafoussam, Garoua) et les pays voisins enclavés (Tchad, Centrafrique).

\textbf{Le port en eau profonde de Kribi.} Inauguré en 2015 puis agrandi, le \cle{port en eau profonde de Kribi} peut accueillir de grands porte-conteneurs que le port de Douala, ensablé et situé sur l'estuaire du Wouri, ne peut recevoir. Il vise à faire du Cameroun une \emph{porte d'entrée maritime} de l'Afrique centrale, au service du Tchad et de la Centrafrique. Il complète le port de Douala, qui reste le premier port du pays et concentre l'essentiel du trafic.

\courssub{Les limites : fractures numériques et inégalités d'infrastructures}

Ces révolutions techniques ne profitent pas à tous les espaces de la même manière. Le monde « rétrécit » pour ceux qui sont connectés, mais reste immense pour les autres.

\begin{definition}[La fracture numérique]
La \cle{fracture numérique} désigne les inégalités d'accès aux technologies de l'information et de la communication (internet, ordinateurs, téléphonie mobile) entre les pays, et à l'intérieur des pays entre les villes et les campagnes, les riches et les pauvres. Elle crée un monde à deux vitesses : les espaces connectés participent pleinement à la mondialisation, les espaces déconnectés en sont marginalisés.
\end{definition}

\textbf{À l'échelle mondiale}, les infrastructures sont concentrées dans les pays développés et les émergents : les grands ports, les hubs aéroportuaires, les centres de données et les points d'atterrissement des câbles se situent d'abord dans la Triade. Une part importante de la population mondiale, surtout en Afrique subsaharienne et dans les campagnes d'Asie du Sud, n'a pas un accès régulier à internet haut débit.

\textbf{À l'échelle du Cameroun}, la fracture numérique est également visible :

\begin{itemize}
\item les villes (Douala, Yaoundé) disposent de la 4G et de la fibre, tandis que de nombreuses zones rurales, surtout dans l'Extrême-Nord, l'Est et l'Adamaoua, sont mal couvertes ;
\item le coût des données reste élevé par rapport aux revenus, ce qui limite l'usage d'internet ;
\item l'électricité, condition de tout usage numérique, fait elle-même défaut dans de nombreuses localités rurales.
\end{itemize}

\begin{exercice}[Pour s'entraîner]
1) Définis la conteneurisation et explique en quatre lignes pourquoi elle a fait baisser les coûts du commerce mondial. 2) Sur un croquis du monde, place et nomme quatre points de passage obligés du trafic maritime et justifie le choix de l'un d'eux. 3) Cite deux câbles sous-marins qui desservent le Cameroun et leur ville d'atterrissement. 4) Explique en cinq lignes ce qu'est la fracture numérique au Cameroun.
\end{exercice}

\begin{corrige}[Exercice]
1) La conteneurisation est le transport des marchandises dans des conteneurs normalisés (EVP) passant du camion au navire sans déballage. Elle réduit le temps de manutention dans les ports, donc les coûts de main-d'œuvre et d'immobilisation des navires ; elle permet aussi le gigantisme des porte-conteneurs, qui répartit le coût du voyage sur des milliers de conteneurs. Le fret devient si bon marché que produire loin pour vendre près devient rentable.
2) Panama (entre Pacifique et Caraïbes), Gibraltar (entre Atlantique et Méditerranée), Suez (entre Méditerranée et mer Rouge), Malacca (entre océan Indien et mer de Chine). Malacca, par exemple, est le passage le plus court entre l'Asie de l'Est et l'Europe : environ 30 \% du commerce maritime mondial y transite.
3) SAT-3 et ACE débarquent à Douala ; SAIL débarque à Kribi.
4) Au Cameroun, la fracture numérique oppose les grandes villes bien connectées (4G, fibre, points d'atterrissement des câbles à Douala et Kribi) aux zones rurales mal couvertes, sans électricité stable ni réseau performant ; elle oppose aussi les ménages aisés, qui peuvent payer les données, aux ménages pauvres. Les espaces déconnectés participent moins aux échanges et aux services numériques (mobile money, formation en ligne, commerce électronique).
\end{corrige}

\begin{aretenir}[L'essentiel de la section]
La mondialisation repose sur trois révolutions techniques : la \cle{conteneurisation} et le gigantisme naval, qui ont divisé par plus de dix le coût du fret maritime (80 \% du commerce mondial en volume) ; le \cle{transport aérien}, rapide, réservé aux passagers et aux produits à forte valeur ajoutée ; et les \cle{télécommunications} (câbles sous-marins de fibre optique, satellites, internet, mobile), qui transmettent l'information instantanément. Il en résulte une \cle{compression de l'espace-temps} : le monde « rétrécit ». Le Cameroun s'y raccorde par les câbles SAT-3, ACE et SAIL (Douala, Kribi) et par le port en eau profonde de Kribi. Mais ces infrastructures sont inégalement réparties : les \cle{fractures numériques} creusent les écarts entre espaces connectés et espaces marginalisés.
\end{aretenir}

\coursec{Les facteurs économiques et politiques de la mondialisation}

Dans la section précédente, nous avons vu que la mondialisation repose d'abord sur des \cle{facteurs techniques} : la révolution des transports (conteneurisation, aviation), des télécommunications et d'Internet a réduit les distances et abaissé les coûts de circulation des biens, des personnes et des informations. Mais une technique, aussi performante soit-elle, ne suffit pas : encore faut-il que les hommes et les États \emph{autorisent} et \emph{organisent} ces circulations. C'est là qu'interviennent les \cle{facteurs économiques et politiques} : la libéralisation des échanges, la dérégulation financière, la construction d'organisations régionales intégrées, l'action des États et des institutions internationales. Autrement dit, si la technique a fourni les \emph{tuyaux} de la mondialisation, les choix économiques et politiques en ont \emph{ouvert les vannes}.

\courssub{La libéralisation des échanges}

\textbf{Intuition.} Imaginez un commerçant de Douala qui importe des motos de Chine. Si l'État camerounais prélève une taxe élevée à la frontière (un \emph{droit de douane}), la moto coûtera cher au consommateur et peu de motos seront importées. Si cette taxe baisse ou disparaît, les échanges s'intensifient. C'est exactement ce qui s'est passé à l'échelle du monde depuis 1945 : la baisse générale des barrières douanières a fait exploser le volume du commerce mondial.

\begin{definition}[Le libre-échange]
Le \cle{libre-échange} est une doctrine et une pratique économiques qui consistent à supprimer ou à réduire les obstacles aux échanges internationaux de biens et de services : droits de douane (taxes à l'importation), quotas (limitations de quantités) et autres barrières commerciales. Une \cle{zone de libre-échange} est un ensemble de pays qui suppriment entre eux ces obstacles.
\end{definition}

\textbf{Le mécanisme historique.} Dès 1947, le GATT (\emph{General Agreement on Tariffs and Trade}, Accord général sur les tarifs douaniers et le commerce) engage des cycles de négociations (les « rounds ») pour abaisser progressivement les droits de douane. Ce mouvement culmine avec la création de l'\cle{OMC} (Organisation mondiale du commerce) en \textbf{1995}, qui succède au GATT. L'OMC fixe les règles du commerce mondial, arbitre les litiges entre États et pousse à la poursuite de la libéralisation.

\begin{exemplebox}[L'OMC, arbitre du commerce mondial]
L'OMC compte \textbf{164 membres} qui réalisent à eux seuls \textbf{plus de 98 \% du commerce mondial}. Le Cameroun en est membre depuis 1995. Le taux moyen des droits de douane dans les pays développés est passé d'environ 40 \% en 1947 à moins de 5 \% aujourd'hui : en un demi-siècle, le volume du commerce mondial de marchandises a été multiplié par plus de 40.
\end{exemplebox}

Les \cle{accords commerciaux} bilatéraux (entre deux pays) ou multilatéraux (entre plusieurs pays) complètent ce dispositif : ils organisent la suppression réciproque des barrières et créent un cadre juridique stable qui rassure les entreprises et les investisseurs.

\courssub{La dérégulation financière}

\textbf{Intuition.} Depuis les années 1980, l'argent circule entre les places financières du monde aussi vite qu'un message sur Internet. Un opérateur à New York peut acheter des titres à Tokyo en quelques secondes. Cette fluidité est le résultat d'un choix politique : la \cle{dérégulation}.

\begin{definition}[La dérégulation financière]
La \cle{dérégulation financière} désigne l'ensemble des mesures prises par les États, à partir des années 1980 (sous l'impulsion des États-Unis et du Royaume-Uni), pour supprimer les contrôles sur les mouvements de capitaux : fin du contrôle des changes, libéralisation des taux d'intérêt, ouverture des marchés financiers aux investisseurs étrangers.
\end{definition}

Ses conséquences sont majeures pour la mondialisation :

\begin{itemize}
\item \textbf{L'interconnexion des bourses} : grâce aux réseaux informatiques, les places financières (New York, Londres, Tokyo, Hong Kong, Francfort) fonctionnent en relais les unes des autres ; quand l'une ferme, une autre ouvre. La finance mondiale travaille ainsi \textbf{24 heures sur 24}.
\item \textbf{La circulation instantanée des capitaux} : les sommes échangées chaque jour sur les marchés des changes dépassent \num{7500} milliards de dollars, soit plusieurs dizaines de fois la valeur du commerce réel de biens et de services. Les capitaux se déplacent en un clic vers les placements les plus rentables.
\item \textbf{Le financement des firmes transnationales} : cette fluidité permet aux entreprises d'investir massivement à l'étranger (les \emph{investissements directs étrangers}, IDE), moteur essentiel de la mondialisation productive.
\end{itemize}

\begin{attention}[Ne pas confondre flux réels et flux financiers]
Une erreur fréquente consiste à présenter les flux financiers comme des échanges de marchandises. Or l'essentiel des capitaux qui circulent chaque jour correspond à des opérations \emph{financières} (achat-vente de devises, d'actions, de titres), et non au commerce de biens. Au Bac, distinguez toujours les \cle{flux commerciaux} (marchandises) des \cle{flux financiers} (capitaux).
\end{attention}

\courssub{Les organisations régionales intégrées}

Parallèlement à la libéralisation mondiale, les États se regroupent à l'échelle \emph{régionale} pour supprimer les barrières entre voisins : c'est la \cle{mondialisation des échanges par blocs régionaux}.

\begin{itemize}
\item \textbf{L'Union européenne (UE)} : l'intégration la plus poussée au monde. Marché unique (libre circulation des biens, des capitaux, des services et des personnes), union douanière, et même une monnaie commune, l'euro, pour 20 de ses 27 membres.
\item \textbf{L'ALENA devenu l'ACEUM} : l'Accord de libre-échange nord-américain (États-Unis, Canada, Mexique), signé en 1994, a été remplacé en 2020 par l'ACEUM (Accord Canada--États-Unis--Mexique), qui modernise les règles commerciales entre les trois pays.
\item \textbf{L'ASEAN} : l'Association des nations de l'Asie du Sud-Est (10 membres, dont l'Indonésie, la Thaïlande, le Vietnam) organise une zone de libre-échange dynamique, au cœur des ateliers industriels du monde.
\item \textbf{En Afrique} : la \cle{CEMAC} (Communauté économique et monétaire de l'Afrique centrale : Cameroun, Gabon, Congo, Tchad, Centrafrique, Guinée équatoriale) et la \cle{CEDEAO} (Communauté économique des États de l'Afrique de l'Ouest, 15 membres autour du Nigeria) organisent l'intégration régionale du continent.
\end{itemize}

\begin{definition}[La ZLECAf]
La \cle{ZLECAf} (Zone de libre-échange continentale africaine) est un accord entré en vigueur en \textbf{2021} qui vise à créer un marché unique africain en supprimant progressivement les droits de douane sur la quasi-totalité des biens échangés entre les pays du continent. Elle regroupe potentiellement plus de 1,3 milliard de consommateurs, ce qui en fait la plus grande zone de libre-échange du monde par le nombre de pays participants. Son siège est à Accra (Ghana).
\end{definition}

\begin{savaistu}[Le Cameroun et l'intégration africaine]
Le Cameroun est membre de la \cle{CEMAC} et a ratifié la \cle{ZLECAf}, dont le siège se trouve à Accra (Ghana). Objectif affiché : vendre davantage de produits camerounais (cacao transformé, bois, produits agricoles) aux autres pays africains, au lieu d'exporter presque uniquement des matières premières brutes vers l'Europe et l'Asie. Le commerce intra-africain ne représente encore qu'environ 15 \% du commerce total du continent, contre plus de 60 \% en Europe : la marge de progression est immense.
\end{savaistu}

\begin{popfigure}
\begin{center}
\begin{tikzpicture}[scale=1,font=\small]
% Planisphère très simplifié en blocs
\fill[popcream] (-7.2,-3.2) rectangle (7.2,3.2);
% Amérique du Nord
\draw[fill=popblue!50,draw=popdark] (-6.6,0.6) -- (-4.6,2.6) -- (-2.6,2.4) -- (-2.4,1.0) -- (-3.6,0.4) -- (-5.2,0.2) -- cycle;
% Amérique du Sud
\draw[fill=popmuted!40,draw=popdark] (-4.4,0.2) -- (-2.8,0.0) -- (-2.4,-1.4) -- (-3.2,-2.8) -- (-4.0,-1.2) -- cycle;
% Europe
\draw[fill=popgreen!55,draw=popdark] (-0.6,1.4) -- (1.0,2.6) -- (2.0,2.2) -- (1.6,1.2) -- (0.4,1.0) -- cycle;
% Afrique
\draw[fill=popgold!60,draw=popdark] (-0.6,0.8) -- (1.4,0.9) -- (2.0,-0.6) -- (1.2,-2.4) -- (0.2,-2.6) -- (-0.8,-0.8) -- cycle;
% Asie
\draw[fill=popmuted!40,draw=popdark] (2.2,2.4) -- (6.6,2.6) -- (6.8,0.8) -- (4.6,0.4) -- (2.4,1.0) -- cycle;
% Asie du Sud-Est
\draw[fill=poporange!60,draw=popdark] (3.4,0.2) -- (4.8,0.2) -- (5.0,-0.8) -- (3.6,-0.6) -- cycle;
% Océanie
\draw[fill=popmuted!40,draw=popdark] (5.2,-1.6) -- (6.6,-1.4) -- (6.4,-2.4) -- (5.4,-2.4) -- cycle;
% Étiquettes
\node[font=\small\bfseries,text=popdark] at (-4.5,1.6) {ACEUM};
\node[font=\small\bfseries,text=popdark] at (0.7,1.9) {UE};
\node[font=\small\bfseries,text=popdark] at (4.2,-0.2) {ASEAN};
\node[font=\small\bfseries,text=popdark] at (0.6,-0.9) {ZLECAf};
\node[font=\scriptsize,text=popdark] at (0.6,-1.4) {CEMAC, CEDEAO};
% Nord
\draw[-{Stealth},thick,popdark] (-6.8,-2.6) -- (-6.8,-1.6);
\node[font=\scriptsize,popdark] at (-6.8,-2.9) {N};
\end{tikzpicture}
\end{center}
\legende{Carte schématique des principales organisations régionales d'intégration dans le monde : ACEUM (Amérique du Nord), Union européenne, ASEAN (Asie du Sud-Est) et, en Afrique, ZLECAf, CEMAC et CEDEAO. Croquis simplifié, sans précision cartographique.}
\end{popfigure}

\courssub{Le rôle des États : ouverture et résistances}

Contrairement à une idée répandue, les États n'ont pas disparu avec la mondialisation : ils en sont des \emph{acteurs} essentiels, mais aussi parfois des \emph{freins}.

\textbf{Les politiques d'ouverture.} De nombreux États mènent des politiques d'\cle{attractivité} pour attirer les investissements étrangers : création de \cle{zones franches} et de zones économiques spéciales (où les entreprises bénéficient de taxes réduites), simplification administrative, construction d'infrastructures (ports, autoroutes, fibre optique). La Chine, avec ses zones économiques spéciales comme Shenzhen dès 1980, est l'exemple le plus célèbre de réussite de cette stratégie. Au Cameroun, le port en eau profonde de Kribi s'inscrit dans cette logique d'attractivité.

\textbf{Les protections et les résistances.} Mais les États continuent aussi de \emph{protéger} certains secteurs jugés stratégiques : subventions à l'agriculture (aux États-Unis comme dans l'Union européenne), droits de douane maintenus sur les produits sensibles, protection des \emph{industries naissantes} dans les pays en développement.

\begin{attention}[Erreur fréquente : la libéralisation n'est pas totale]
Ne présentez jamais la libéralisation des échanges comme achevée ou totale. De nombreux États protègent encore fortement des secteurs entiers : l'\textbf{agriculture} (subventions et tarifs élevés dans les pays riches), les \textbf{industries naissantes} des pays en développement (qui ne pourraient pas survivre face à la concurrence mondiale sans protection temporaire), ou encore les services publics. Le monde actuel combine libre-échange \emph{et} protections.
\end{attention}

\begin{savaistu}[Complément utile au Bac : le retour des tensions commerciales]
Depuis 2018, la \cle{guerre commerciale entre les États-Unis et la Chine} illustre les limites du libre-échange : les deux puissances se sont mutuellement imposé des droits de douane punitifs sur des centaines de milliards de dollars de marchandises (acier, aluminium, produits électroniques, soja...). Cet épisode montre que la mondialisation n'est pas un processus irréversible : les rivalités entre grandes puissances peuvent freiner, voire partiellement inverser, la libéralisation des échanges. Citer cet exemple dans une copie du Bac témoigne d'une culture actualisée.
\end{savaistu}

\courssub{Le rôle des institutions internationales}

Enfin, la mondialisation est encadrée par de grandes \cle{institutions internationales} créées pour l'essentiel après 1945 (système dit « de Bretton Woods ») :

\begin{itemize}
\item \textbf{Le FMI} (Fonds monétaire international, 1944) : il veille à la stabilité financière mondiale et vient en aide aux États en difficulté de paiement, en échange de réformes économiques (les « plans d'ajustement structurel » des années 1980-1990, qui ont fortement marqué les pays africains, dont le Cameroun).
\item \textbf{La Banque mondiale} (1944) : elle finance le développement (routes, écoles, barrages, santé) par des prêts aux pays pauvres et émergents.
\item \textbf{L'ONU} (Organisation des Nations unies, 1945) : cadre politique de la coopération mondiale ; ses agences spécialisées (PNUD, FAO, UNICEF...) fixent des objectifs communs à l'humanité, comme les Objectifs de développement durable (ODD).
\item \textbf{L'OMC} (1995), déjà étudiée, complète ce dispositif pour le commerce.
\end{itemize}

Ces institutions fixent les \emph{règles du jeu} de la mondialisation : elles encouragent l'ouverture des économies, financent les infrastructures qui rendent les flux possibles et arbitrent les conflits. Mais elles sont aussi critiquées : les pays riches y disposent d'un poids décisionnel plus important, et leurs recettes libérales sont parfois jugées inadaptées aux réalités des pays en développement.

\begin{popfigure}
\begin{center}
\begin{tikzpicture}[font=\small,node distance=1.1cm,
actor/.style={draw=popdark,thick,rounded corners,fill=popblueL,minimum width=3.2cm,minimum height=0.9cm,align=center},
rel/.style={-{Stealth},thick,popdark}]
\node[actor,fill=popgreenL, align=center] (etats){\textbf{États}\\ lois, ouverture, protections};
\node[actor,fill=poporangeL,right=2.6cm of etats, align=center] (ftn){\textbf{FTN}\\ investissements, production};
\node[actor,fill=poppurpleL,below=1.6cm of etats, align=center] (oi){\textbf{Organisations internationales}\\ OMC, FMI, Banque mondiale, ONU};
\node[actor,fill=popgoldL,below=1.6cm of ftn, align=center] (ong){\textbf{ONG et société civile}\\ contestation, régulation sociale};
\draw[rel] (etats) -- node[above,font=\scriptsize]{zones franches, IDE} (ftn);
\draw[rel] (ftn) -- node[right,font=\scriptsize,align=center]{pression,\\ lobbying} (ong);
\draw[rel] (oi) -- node[left,font=\scriptsize,align=center]{règles,\\ financements} (etats);
\draw[rel] (oi) -- node[above,font=\scriptsize]{cadre juridique} (ftn);
\draw[rel] (ong) -- node[below,font=\scriptsize,align=center]{critiques,\\ contre-pouvoirs} (oi);
\draw[rel] (ong) -- node[above,font=\scriptsize,align=center,sloped]{dénonciation} (etats);
\end{tikzpicture}
\end{center}
\legende{Schéma des acteurs de la mondialisation et de leurs relations : les États fixent les politiques, les organisations internationales établissent les règles, les firmes transnationales (FTN) organisent la production mondiale, les ONG exercent des contre-pouvoirs.}
\end{popfigure}

\courssub{Exercice résolu et synthèse}

\begin{exoresolu}[Analyse de document : mesurer la libéralisation]
\textbf{Énoncé.} Le taux moyen des droits de douane à l'importation dans les pays développés est passé d'environ 40 \% en 1947 à environ 4 \% aujourd'hui. Dans le même temps, le volume du commerce mondial de marchandises a été multiplié par plus de 40. 1) Calculez la baisse en points de pourcentage du taux de droits de douane. 2) Quel facteur de la mondialisation cette évolution illustre-t-elle ? Justifiez. 3) Citez l'organisation qui encadre cette évolution depuis 1995.
\tcblower
\textbf{Solution détaillée.}
1) Baisse $= 40 - 4 = 36$ points de pourcentage. Le taux a été divisé par 10.
2) Cette évolution illustre la \cle{libéralisation des échanges}, facteur économique et politique de la mondialisation : la suppression progressive des barrières douanières a réduit le coût des échanges et permis l'explosion du commerce mondial (volume multiplié par plus de 40).
3) Depuis 1995, c'est l'\cle{OMC} (Organisation mondiale du commerce, 164 membres réalisant plus de 98 \% du commerce mondial) qui fixe les règles du commerce international et arbitre les litiges. Avant elle, le GATT (1947) assurait ce rôle.
\end{exoresolu}

\begin{aretenir}[L'essentiel de la section]
\begin{itemize}
\item La \cle{libéralisation des échanges} (baisse des droits de douane, accords commerciaux, OMC créée en 1995) est le principal facteur économique de la mondialisation.
\item La \cle{dérégulation financière} depuis les années 1980 permet la circulation instantanée des capitaux entre des bourses interconnectées 24 h/24.
\item Les \cle{organisations régionales intégrées} (UE, ACEUM, ASEAN ; CEMAC, CEDEAO et ZLECAf en Afrique, entrée en vigueur en 2021) créent des blocs de libre-échange.
\item Les \cle{États} restent centraux : politiques d'ouverture et d'attractivité, mais aussi protections (agriculture, industries naissantes) et tensions (guerre commerciale USA--Chine).
\item Les \cle{institutions internationales} (FMI, Banque mondiale, ONU, OMC) fixent les règles et financent la mondialisation.
\end{itemize}
\end{aretenir}

\coursec{Le fonctionnement de la mondialisation : les flux}

Nous avons vu dans la section précédente que la mondialisation repose sur des facteurs économiques (libéralisation des échanges, rôle des firmes transnationales) et politiques (accords commerciaux, organisations internationales). Mais ces facteurs ne sont que des \emph{conditions} : ils rendent la mondialisation possible. Comment fonctionne-t-elle concrètement, au quotidien ? La réponse tient en un mot : par les \cle{flux}, c'est-à-dire les circulations permanentes de marchandises, de capitaux, d'informations et d'hommes à l'échelle de la planète. Cette section décrit ces flux, les mesure et montre qu'ils dessinent un monde profondément hiérarchisé.

\courssub{La notion de flux : comprendre le fonctionnement de la mondialisation}

\begin{definition}[Flux]
En géographie, un \cle{flux} est un déplacement de biens, de capitaux, d'informations ou de personnes entre deux lieux, matérialisé par un itinéraire (route maritime, câble, liaison aérienne) et mesurable par une quantité (tonnes, dollars, octets, individus). La mondialisation fonctionne par la mise en relation permanente des territoires grâce à ces flux.
\end{definition}

L'intuition est simple : lorsqu'un habitant de Douala achète un téléphone fabriqué en Chine avec des composants coréens, payé en francs CFA convertis en dollars, et qu'il s'en sert pour appeler son frère installé en France, il mobilise en un seul geste quatre types de flux : un flux de marchandises, un flux de capitaux, un flux d'informations et un flux humain (la migration de son frère). La mondialisation, c'est cela : un enchevêtrement de circulations.

On distingue classiquement quatre grandes catégories de flux : les flux de \cle{marchandises}, les flux de \cle{capitaux}, les flux d'\cle{informations} et les flux \cle{humains}.

\courssub{Les flux de marchandises : le commerce mondial}

\courssub{Des échanges massifs et croissants}

Le commerce mondial de marchandises représente environ \num{25000} milliards de dollars par an (données OMC, années 2020-2023). Il a été multiplié par plus de 40 depuis 1960 : c'est le moteur historique de la mondialisation. Deux grandes catégories de biens circulent :

\begin{itemize}
\item les \cle{matières premières} : pétrole et gaz (le produit le plus échangé au monde en valeur), minerais (fer, cuivre, bauxite), produits agricoles tropicaux (cacao, café, coton, bananes) et céréales (blé, maïs, riz) ;
\item les \cle{produits manufacturés} : automobiles, machines, équipements électroniques, textiles, médicaments. Ils représentent aujourd'hui plus de 70\,\% de la valeur du commerce mondial.
\end{itemize}

Les flux de marchandises empruntent principalement les \cle{routes maritimes} : environ 80\,\% du volume du commerce mondial transite par la mer, à bord de porte-conteneurs géants. Les détroits stratégiques (Malacca, Suez, Ormuz, Panama) sont les goulots d'étranglement de la mondialisation.

\courssub{Une orientation dominante Nord-Nord, mais un Sud qui monte}

L'observation des flux montre une structure claire :

\begin{itemize}
\item les échanges \cle{Nord-Nord} (entre pays développés : Amérique du Nord, Europe occidentale, Asie orientale) restent dominants en valeur ;
\item les échanges \cle{Nord-Sud} opposent souvent des exportations de matières premières brutes (Sud) à des exportations de produits manufacturés (Nord) ;
\item les échanges \cle{Sud-Sud} progressent rapidement depuis les années 2000, tirés par la Chine, devenue premier partenaire commercial de l'Afrique et de l'Amérique latine. Le commerce Sud-Sud représente désormais plus de la moitié des exportations des pays en développement.
\end{itemize}

\begin{aretenir}[L'essentiel sur les flux de marchandises]
Le commerce mondial est dominé par les produits manufacturés et par les échanges entre pays riches, mais l'émergence de la Chine et des pays du Sud modifie la carte : les flux Sud-Sud sont la grande nouveauté du XXI\textsuperscript{e} siècle.
\end{aretenir}

\courssub{Les flux de capitaux : l'argent qui circule}

Les capitaux circulent plus vite encore que les marchandises, car ils sont dématérialisés : un ordre de bourse met quelques millisecondes pour aller de New York à Tokyo. On distingue trois types de flux de capitaux.

\begin{definition}[Les IDE]
Les \cle{IDE} (investissements directs à l'étranger, ou investissements directs étrangers) sont des capitaux investis par une entreprise d'un pays dans un autre pays pour y créer, acheter ou développer une activité productive durable (usine, filiale, plantation, mine). Ils se distinguent des placements spéculatifs car ils impliquent un contrôle durable sur l'entreprise (généralement plus de 10\,\% du capital). Les IDE mondiaux représentent environ \num{1300} à \num{1600} milliards de dollars par an selon les années.
\end{definition}

Les \cle{flux financiers spéculatifs} (ou de portefeuille) sont des mouvements de capitaux à court terme (achat et vente d'actions, de devises, de titres de dette) qui recherchent un profit rapide. Leur montant quotidien sur le marché des changes dépasse \num{7500} milliards de dollars par jour (données BRI, 2022) : c'est le flux le plus volumineux de la mondialisation, mais aussi le plus instable (crises financières de 1997 en Asie, de 2008).

Enfin, les \cle{transferts d'argent des migrants} constituent un flux souvent sous-estimé mais vital.

\begin{definition}[Diaspora et transferts de fonds]
La \cle{diaspora} est l'ensemble des personnes originaires d'un même pays ou d'une même communauté, dispersées à l'étranger mais qui entretiennent des liens avec leur pays d'origine. Les \cle{transferts de fonds} (ou \emph{remittances}) sont les sommes d'argent que les migrants envoient à leurs familles restées au pays. À l'échelle mondiale, ces transferts dépassent 800 milliards de dollars par an vers les pays à revenu faible ou intermédiaire (données Banque mondiale, 2023), soit souvent plus que l'aide publique au développement et parfois plus que les IDE reçus par ces pays.
\end{definition}

\begin{exemplebox}[La diaspora camerounaise, un acteur économique majeur]
La diaspora camerounaise est estimée à plusieurs centaines de milliers de personnes, installées surtout en France, aux États-Unis, en Allemagne et au Canada. Ses transferts de fonds vers le Cameroun dépassent \textbf{300 millions de dollars par an} (plus de 180 milliards de FCFA, données Banque mondiale/BCEAO récentes). Cet argent finance la scolarité des enfants, la construction de maisons, les soins de santé et de petits commerces : c'est un flux vital pour de nombreuses familles, et un apport en devises non négligeable pour l'économie nationale.
\end{exemplebox}

\courssub{Les flux d'informations : l'instantanéité}

Les flux d'informations sont la grande révolution récente de la mondialisation. Grâce à \cle{internet}, aux câbles sous-marins de fibre optique (qui transportent plus de 95\,\% des données intercontinentales), aux satellites et aux réseaux de téléphonie mobile, l'information circule de manière \emph{instantanée} et \emph{dématérialisée}.

\begin{itemize}
\item Plus de 5,3 milliards de personnes utilisent internet dans le monde (UIT, 2023) ; les \cle{réseaux sociaux} (Facebook, WhatsApp, TikTok, X) relient des milliards d'individus en temps réel.
\item Les \cle{médias mondiaux} (CNN, BBC, Al-Jazeera, TV5 Monde, France 24) diffusent les mêmes images simultanément sur toute la planète : un événement devient mondial en quelques minutes.
\item Au Cameroun, le téléphone mobile est devenu le premier outil d'accès à l'information et aux services financiers (\emph{mobile money} : MTN MoMo, Orange Money), illustrant l'intégration du pays dans les flux numériques mondiaux.
\end{itemize}

\begin{attention}[Piège : l'instantanéité ne signifie pas l'égalité]
Tous les territoires ne sont pas connectés de la même façon. Les débits, le coût d'accès et la couverture restent très inégaux : on parle de \cle{fracture numérique}. Dire « le monde entier est connecté » est une erreur fréquente à l'examen : il faut toujours préciser que la connexion est inégale selon les pays et, à l'intérieur des pays, entre villes et campagnes.
\end{attention}

\courssub{Les flux humains : migrations et tourisme}

Contrairement aux capitaux et aux informations, les hommes circulent moins librement : les frontières et les visas filtrent les mobilités. Pourtant, les flux humains n'ont jamais été aussi importants.

\begin{itemize}
\item Les \cle{migrations internationales} concernent environ \textbf{281 millions de personnes} (environ 3,6\,\% de la population mondiale, ONU 2020). Elles vont principalement du Sud vers le Nord (Amérique du Nord, Europe, Golfe persique), mais aussi du Sud vers le Sud (migrations intra-africaines, très nombreuses : vers l'Afrique du Sud, la Côte d'Ivoire, le Gabon).
\item Les \cle{migrations de travail} répondent aux besoins de main-d'œuvre des pays riches (travailleurs qualifiés : médecins, ingénieurs — on parle alors de \emph{fuite des cerveaux} ou \emph{mobilité des talents} ; et travailleurs peu qualifiés : bâtiment, agriculture, services).
\item Le \cle{tourisme international} dépasse un milliard de touristes par an (avant la crise sanitaire de 2020, reprise forte depuis 2023) : c'est un flux humain temporaire mais massif, source de devises pour de nombreux pays du Sud.
\item La \cle{diaspora camerounaise} illustre ces mobilités : étudiants partis se former en Europe ou en Amérique du Nord, travailleurs qualifiés, sportifs et artistes de renommée mondiale.
\end{itemize}

\courssub{Les acteurs du fonctionnement : FTN, métropoles et États}

Le fonctionnement des flux ne se fait pas tout seul ; il est piloté, organisé et régulé par des acteurs majeurs.

\begin{itemize}
\item Les \cle{firmes transnationales (FTN)} sont les \textbf{premiers organisateurs} des flux. Elles conçoivent les chaînes de valeur mondiales : elles décident où produire (usines), quoi acheter (matières premières), où vendre (marchés) et comment financer (IDE, flux intra-groupe). Exemples : Apple, TotalEnergies, Samsung, Dangote Group.
\item Les \cle{métropoles mondiales} (New York, Londres, Tokyo, Shanghai, Singapour, Dubaï, Paris, Francfort) sont les \textbf{nœuds de commande} où les flux se croisent : sièges de FTN, places boursières, grands ports/aéroports, centres de décision politique et culturelle. Elles forment un \emph{archipel métropolitain} au-dessus des États.
\item Les \cle{États} et les \cle{organisations internationales} (OMC, FMI, Banque mondiale, ONU, unions régionales comme l'UE, l'ALENA/USMCA, la ZLECAf) fixent le \textbf{cadre réglementaire} : droits de douane, normes, accords de libre-échange, politiques de change, lois sur l'immigration. Ils peuvent faciliter ou entraver les flux.
\end{itemize}

\begin{aretenir}[Acteurs et flux]
FTN = organisatrices (production, investissement) ; Métropoles = nœuds de connexion et de commande ; États/OI = régulateurs (règles, frontières, accords).
\end{aretenir}

\courssub{La hiérarchie des flux : un monde inégalement connecté}

L'observation cartographique des flux révèle une loi fondamentale : les flux ne se répartissent pas uniformément, ils se \emph{concentrent}.

\begin{itemize}
\item La \cle{Triade} (Amérique du Nord, Europe occidentale, Asie orientale) concentre l'essentiel des flux de marchandises, de capitaux et d'informations. Les trois grandes façades maritimes (Asie de l'Est, Europe du Nord-Ouest, Amérique du Nord-Est) polarisent le commerce mondial.
\item Les \cle{métropoles} mondiales sont les nœuds où ces flux se croisent et se commandent : bourses, sièges de FTN, grands ports et aéroports.
\item À l'inverse, de nombreux territoires sont \emph{marginalisés} : l'intérieur de l'Afrique, les zones rurales enclavées, les pays sans littoral (pays enclavés) ne captent qu'une faible part des flux mondiaux.
\end{itemize}

\begin{attention}[Erreur fréquente : « l'Afrique est hors de la mondialisation »]
C'est faux. L'Afrique est bien \emph{intégrée} à la mondialisation : elle exporte massivement du pétrole, des minerais, du cacao, du café ; elle reçoit des IDE, des transferts de migrants et des produits manufacturés. Mais elle y est intégrée de façon \textbf{inégale et subordonnée} : elle fournit surtout des matières premières \emph{brutes}, peu transformées, dont les prix sont fixés sur les marchés mondiaux, et elle importe des produits manufacturés à forte valeur ajoutée. Le problème de l'Afrique n'est donc pas l'exclusion, mais une \emph{position défavorable} dans la hiérarchie des flux.
\end{attention}

\begin{popfigure}
\begin{center}
\begin{tikzpicture}[scale=1.1, font=\small]
% Fond océan
\fill[popblueL!50] (-0.5,-2.8) rectangle (13.0,4.6);
% Continents très simplifiés (polygones fermés)
\fill[popgreenL] (0.5,1.5) -- (0.8,3.8) -- (2.5,4.3) -- (3.8,3.5) -- (3.4,1.6) -- (2.0,0.8) -- cycle; % Am Nord
\fill[popgreenL] (2.8,0.8) -- (3.4,0.6) -- (3.8,-1.0) -- (3.3,-2.6) -- (2.7,-1.4) -- cycle; % Am Sud
\fill[popgreenL] (5.5,2.4) -- (5.8,3.9) -- (7.5,4.1) -- (7.8,2.8) -- (6.8,2.2) -- cycle; % Europe
\fill[popgreenL] (5.3,2.1) -- (7.0,2.0) -- (7.4,0.3) -- (6.5,-1.8) -- (5.7,-0.6) -- cycle; % Afrique
\fill[popgreenL] (8.0,2.8) -- (8.5,4.2) -- (11.8,3.9) -- (12.2,2.4) -- (10.5,1.8) -- (8.8,2.2) -- cycle; % Asie
\fill[popgreenL] (10.5,0.3) -- (11.5,0.5) -- (11.7,-1.0) -- (10.7,-1.2) -- cycle; % Australie
% Labels continents
\node at (2.2,2.9) {\textbf{Amérique du Nord}};
\node at (3.2,-0.2) {\textbf{Amérique du Sud}};
\node at (6.7,3.5) {\textbf{Europe}};
\node at (6.2,0.5) {\textbf{Afrique}};
\node at (10.2,3.3) {\textbf{Asie orientale}};
% Flèches MARCHANDISES (Triade) - Bleues, ÉPAISSES
\draw[line width=3.5pt, popblue, -{Stealth[length=4mm]}] (3.8,3.0) to[bend left=10] (5.6,3.4); % AmN -> Eur
\draw[line width=3.5pt, popblue, -{Stealth[length=4mm]}] (7.8,3.2) to[bend left=10] (9.2,3.3); % Eur -> As
\draw[line width=3.5pt, popblue, -{Stealth[length=4mm]}] (9.5,2.2) to[bend left=20] (3.8,2.4); % As -> AmN (pacifique/panama)
% Flèches MATIÈRES PREMIÈRES (Afrique -> Triade) - Vertes, POINTILLÉES, moyennes
\draw[line width=1.8pt, popgreen!60!black, dashed, -{Stealth[length=3mm]}] (6.0,1.2) to[bend left=15] (6.5,2.6); % Afr -> Eur
\draw[line width=1.8pt, popgreen!60!black, dashed, -{Stealth[length=3mm]}] (6.8,0.6) to[bend left=25] (9.2,2.0); % Afr -> As
% Flèches IDE (Triade -> Afrique/Sud) - Orange, FINES
\draw[line width=1.2pt, poporange, -{Stealth[length=3mm]}] (5.8,2.8) to[bend right=15] (6.2,1.0); % Eur -> Afr
\draw[line width=1.2pt, poporange, -{Stealth[length=3mm]}] (9.0,2.5) to[bend right=20] (6.5,0.8); % As -> Afr
\draw[line width=1.2pt, poporange, -{Stealth[length=3mm]}] (3.0,2.0) to[bend right=20] (3.2,-0.5); % AmN -> AmS
% Nord
\draw[-{Stealth}, popdark, thick] (0.5,-1.8) -- (0.5,-1.0); \node[below] at (0.5,-1.8) {\textbf{N}};
\end{tikzpicture}
\end{center}
\legende{Croquis schématique des principaux flux mondiaux. \textbf{Flèches bleues épaisses} : grands flux de marchandises (produits manufacturés) au sein de la Triade. \textbf{Flèches vertes pointillées} : exportations de matières premières de l'Afrique vers l'Europe et l'Asie. \textbf{Flèches orange fines} : flux d'IDE (investissements) de la Triade vers le Sud. L'épaisseur des flèches est proportionnelle à l'importance quantitative/valorique du flux.}
\end{popfigure}

\begin{methode}[Construire une carte de flux (Croquis de géographie)]
Pour réaliser un croquis de flux au Baccalauréat technique / Probatoire, suivre rigoureusement les étapes :
\begin{enumerate}
\item \textbf{Fond de carte} : Tracer le contour simplifié des continents ou pays concernés. Repérer et nommer les lieux de départ et d'arrivée des flux (villes, ports, pays, régions).
\item \textbf{Flèches orientées} : Représenter chaque flux par une flèche partant du lieu d'origine vers le lieu de destination.
\item \textbf{Proportionnalité} : Faire varier l'\textbf{épaisseur} de la flèche de façon \textbf{proportionnelle} à la valeur du flux (flux majeur = flèche épaisse, flux secondaire = flèche fine).
\item \textbf{Distinction visuelle} : Différencier les natures de flux (marchandises, capitaux, informations, hommes) par des \textbf{couleurs} distinctes et/ou des \textbf{types de traits} (plein, pointillé, traitillé).
\item \textbf{Éléments obligatoires} : \textbf{Titre} précis, \textbf{Légende} explicative (couleur/épaisseur = nature/quantité), \textbf{Flèche du Nord}, \textbf{Échelle indicative} (ex: 1 cm = 5000 km ou « Échelle non respectée »).
\end{enumerate}
\end{methode}

\courssub{Étude de cas : les flux commerciaux du Cameroun}

Le Cameroun illustre parfaitement l'insertion subordonnée d'un pays africain dans les flux mondiaux.

\textbf{Les exportations} sont dominées par les matières premières brutes ou peu transformées (données INS / Douanes, ordres de grandeur années récentes) :

\begin{itemize}
\item le \cle{pétrole brut} (environ 30 à 40\,\% des exportations en valeur), extrait dans le bassin du Rio del Rey (offshore) et exporté par le terminal de Kole ;
\item le \cle{cacao} (environ 12 à 15\,\%), dont le Cameroun est le 5\textsuperscript{e} producteur mondial, exporté en fèves via Douala ;
\item le \cle{bois} (environ 8 à 10\,\%), exporté encore trop souvent en grumes (bois ronds) depuis le port de Douala, malgré la volonté de transformation locale ;
\item le \cle{coton} (environ 4 à 6\,\%), produit dans le Nord (régions de l'Adamaoua, Nord, Extrême-Nord) par la SODECOTON, exporté en fibre ;
\item s'ajoutent le café (arabica/robusta), la banane (plantations du Mungo), l'aluminium (Alucam, énergie hydroélectrique) et le caoutchouc.
\end{itemize}

\textbf{Les importations} sont composées de produits à forte valeur ajoutée : produits manufacturés (machines, véhicules, équipements électriques, pharmaceutiques), céréales (blé, riz — dépendance alimentaire) et \emph{carburants raffinés} (essence, gasoil) — paradoxe d'un pays producteur de pétrole brut qui importe ses produits raffinés, faute de capacités de raffinage suffisantes (projet de raffinerie de Kribi en cours).

Les principaux partenaires sont la Chine (premier client et premier fournisseur), l'Union européenne (France, Pays-Bas, Espagne, Italie), l'Inde et les pays voisins de la CEMAC/CEEAC. Le port autonome de \cle{Douala} concentre l'essentiel de ces flux (95\,\% du trafic portuaire), relayé par le port en eau profonde de \cle{Kribi} (terminal minéralier, conteneurs, futur hub logistique).

\begin{popfigure}
\begin{center}
\begin{tikzpicture}
\begin{axis}[
    ybar, bar width=18pt,
    width=0.95\linewidth, height=6.5cm,
    ymin=0, ymax=45,
    ylabel={Part des exportations (\% de la valeur)},
    symbolic x coords={{Pétrole brut},{Cacao},{Bois},{Coton},{Autres produits}},
    xtick=data, 
    x tick label style={rotate=20, anchor=east, font=\small},
    nodes near coords, nodes near coords style={font=\small, popdark},
    ymajorgrids, grid style={popline!40},
    axis line style={popdark},
    enlarge x limits=0.15
]
\addplot[fill=popblue, draw=popdark] coordinates {({Pétrole brut},35) ({Cacao},14) ({Bois},9) ({Coton},5) ({Autres produits},37)};
\end{axis}
\end{tikzpicture}
\end{center}
\legende{Structure des exportations du Cameroun (parts approximatives en \% de la valeur totale, moyennes récentes ; source : INS, Douanes, BCEAO). La catégorie « Autres produits » regroupe café, bananes, aluminium, caoutchouc, produits manufacturés divers.}
\end{popfigure}

\begin{exoresolu}[Analyser les flux commerciaux du Cameroun]
\textbf{Énoncé.} À partir du graphique ci-dessus et de vos connaissances : 1) Calculez la part cumulée des quatre principaux produits d'exportation du Cameroun. 2) Montrez que cette structure révèle une insertion subordonnée du Cameroun dans la mondialisation. 3) Proposez une solution concrète pour améliorer cette insertion.
\tcblower
\textbf{Solution détaillée.}
\begin{enumerate}
\item \textbf{Part cumulée} : $35 + 14 + 9 + 5 = 63\,\%$. Les quatre principaux produits (pétrole, cacao, bois, coton) représentent donc environ 63\,\% des exportations camerounaises : le commerce extérieur du pays repose sur une base très étroite (faible diversification).
\item \textbf{Insertion subordonnée} : Ces quatre produits sont tous des \emph{matières premières brutes ou peu transformées} (pétrole brut, cacao en fèves, bois en grumes, coton fibre). Leur prix est fixé sur les marchés mondiaux (cours de Bourse : Brent, Londres, New York, Cotlook A) et varie fortement (volatilité). En échange, le Cameroun importe des produits manufacturés (machines, véhicules, médicaments) et même des carburants raffinés, dont les prix sont plus stables et la valeur ajoutée élevée. Cette asymétrie structurelle (exporter du brut à bas prix, importer du transformé à prix fort) est la marque d'une insertion \emph{subordonnée} : le pays capte une faible part de la valeur ajoutée mondiale et subit la volatilité des cours.
\item \textbf{Solution} : Mettre en œuvre une politique de \emph{transformation locale} (industrialisation) : raffinage du pétrole sur place (projet Kribi), création d'unités de broyage de cacao (pâte, beurre, poudre), usines de sciage/contreplaqué/bois meuble, filatures et tissages de coton. Cela permet d'exporter des produits à plus forte valeur ajoutée, de créer des emplois locaux, de réduire la facture d'importation et de stabiliser les recettes d'exportation. C'est l'objectif de la \cle{Stratégie Nationale de Développement (SND30)} et de la \cle{ZLECAf}.
\end{enumerate}
\end{exoresolu}

\begin{exercice}[Pour s'entraîner -- Type Probatoire/Bac Tech]
1) Définissez les IDE et distinguez-les clairement des flux financiers spéculatifs (de portefeuille). 
2) Expliquez pourquoi les transferts de la diaspora constituent un flux essentiel et structurel pour l'économie camerounaise. 
3) \textbf{Croquis} : Sur un fond de carte du monde simplifié (fourni ou à tracer), représentez \textbf{trois flux majeurs de marchandises} (dont un Sud-Sud) et \textbf{deux flux d'IDE} en respectant la méthode de la carte de flux (titre, légende, flèches proportionnelles, nord, distinction marchandises/IDE).
\end{exercice}

\begin{corrige}[Exercice 1 à 3]
1) \textbf{IDE} : Investissements directs à l'étranger = capitaux investis durablement (\textgreater 10\,\% capital) pour créer/contrôler une activité productive (usine, filiale, mine). \textbf{Flux spéculatifs} : Achats/ventes d'actifs financiers (actions, obligations, devises) à court terme pour gain de cours, sans contrôle de l'entreprise. Volatiles, massifs (7500 Mds \$/jour).
2) Les transferts (300 Mds \$/an env.) financent directement la consommation des ménages (santé, éducation, logement), l'investissement productif (petits commerces, agriculture) et apportent des devises stables (contre-cycliques), souvent supérieures à l'APD. Ils réduisent la pauvreté et soutiennent le taux de change.
3) \textbf{Attendus du croquis} : 
\begin{itemize}
\item Fond de carte : contours continents, noms océans, Triade, Afrique, Chine.
\item Marchandises (flèches \textbf{bleues épaisses}) : Asie $\rightarrow$ AmNord (électronique), Asie $\rightarrow$ Europe (biens de conso), AmNord $\leftrightarrow$ Europe (machines, chimie), \textbf{Chine $\rightarrow$ Afrique} (manufacturés) ou \textbf{Afrique $\rightarrow$ Chine} (minerais/pétrole) pour le flux Sud-Sud.
\item IDE (flèches \textbf{orange fines}) : Triade $\rightarrow$ Afrique (pétrole, mines, télécoms), Triade $\rightarrow$ Asie du Sud-Est (usines délocalisées).
\item Légende : code couleur/épaisseur. Titre : « Principaux flux de marchandises et d'IDE dans la mondialisation ». Nord + Échelle.
\end{itemize}
\end{corrige}

\begin{aretenir}[Bilan de la section : Les mécanismes de la mondialisation]
La mondialisation fonctionne par quatre grands flux interconnectés : \textbf{marchandises} (commerce, 80\,\% mer), \textbf{capitaux} (IDE durables, spéculation volatile, transferts migrants vitaux), \textbf{informations} (instantanées, câbles, fracture numérique), \textbf{hommes} (migrations travail/études, tourisme). Ces flux sont \textbf{organisés par les FTN}, \textbf{commandés depuis les métropoles} et \textbf{régulés par les États/OI}. Ils dessinent une \textbf{hiérarchie mondiale} : concentration dans la Triade, marginalisation relative de l'Afrique (position subordonnée : export brut / import transformé), illustrée par le cas du Cameroun.
\end{aretenir}

\coursec{Les acteurs du fonctionnement : FTN, métropoles et États}

\courssub{Transition : des flux aux acteurs qui les commandent}

Dans la section précédente, nous avons analysé les \textbf{flux} (marchandises, capitaux, informations, hommes) qui irriguent la mondialisation. Ces flux ne circulent pas spontanément : ils sont organisés, orientés, accélérés ou freinés par des \textbf{acteurs} puissants qui détiennent le pouvoir de décision. Comprendre la mondialisation, c'est identifier \emph{qui} commande, \emph{d'où} partent les ordres et \emph{comment} s'articulent les stratégies des firmes, la concentration métropolitaine et l'action des États. C'est l'objet de cette section.

\courssub{Les firmes transnationales (FTN) : moteurs et architectes de la division internationale du travail}

\begin{definition}[Firme transnationale (FTN)]
Une \cle{firme transnationale (FTN)} (ou multinationale) est une entreprise qui possède ou contrôle des unités de production ou de services (filiales, succursales) dans au moins deux pays. Elle se caractérise par une \textbf{stratégie unifiée} pilotée depuis son siège social (pays d'origine) et une \textbf{recherche permanente d'optimisation mondiale} de ses coûts et de son accès aux marchés.
\end{definition}

\begin{methode}[Analyser la stratégie spatiale d'une FTN]
\begin{enumerate}
    \item \textbf{Localiser le siège social} (pays d'origine, fonctions de commandement : R\&D, finance, marketing global).
    \item \textbf{Identifier les implantations} : usines (production), centres de R\&D délocalisés, plateformes logistiques, filiales commerciales.
    \item \textbf{Qualifier la Division Internationale du Travail (DIT)} : quel pays fait quoi ? (conception, composants, assemblage, vente, recyclage).
    \item \textbf{Expliquer les facteurs de localisation} : coûts du travail, compétences, proximité des marchés, stabilité politique, incitations fiscales.
\end{enumerate}
\end{methode}

\begin{definition}[Division internationale du travail (DIT)]
La \cle{Division internationale du travail (DIT)} désigne la spécialisation des pays ou des régions dans des tâches précises au sein d'un même processus de production mondialisé. Elle distingue classiquement : les \textbf{pays du Nord} (conception, R\&D, haute technologie, finance, services à haute valeur ajoutée) et les \textbf{pays du Sud/émergents} (extraction matières premières, assemblage, manufacturing bas/moyen de gamme). La \emph{Nouvelle DIT} (depuis les années 1990) voit des pays émergents (Chine, Inde, Vietnam) monter en gamme (R\&D, composants high-tech).
\end{definition}

\begin{exemplebox}[La chaîne de production mondialisée de l'iPhone (Apple)]
Apple (USA) est l'archétype de la FTN \emph{sans usine} (\emph{fabless}).
\begin{itemize}
    \item \textbf{Conception, design, OS, marketing} : Cupertino, Californie (USA) — \emph{commandement, haute valeur ajoutée}.
    \item \textbf{Composants stratégiques} : Processeurs (TSMC, Taïwan / Samsung, Corée du Sud), écrans (Corée du Sud, Japon, Chine), capteurs (Japon, Europe).
    \item \textbf{Assemblage final} : Zhengzhou (Chine, Foxconn — « iPhone City »), de plus en plus en Inde (Tamil Nadu) et au Brésil.
    \item \textbf{Vente \& Services} : Réseau mondial (Apple Stores, revendeurs, App Store).
    \item \textbf{Recyclage} : Robot « Daisy » (USA, Europe) + filières informelles (Afrique, Asie).
\end{itemize}
Apple capte la majeure partie de la valeur (marge, propriété intellectuelle) ; l'assemblage ne représente que $\approx 3\text{--}5\%$ du prix final.
\end{exemplebox}

\begin{popfigure}
\begin{tikzpicture}[scale=0.85, every node/.style={font=\small}]
    % Styles
    \tikzset{
        hub/.style={rectangle, rounded corners, draw=popblue, thick, fill=popblueL, minimum width=2.5cm, minimum height=1.2cm, align=center},
        prod/.style={rectangle, rounded corners, draw=poporange, thick, fill=poporangeL, minimum width=2.2cm, minimum height=1cm, align=center},
        comp/.style={ellipse, draw=popgreen, thick, fill=popgreenL, minimum width=2cm, minimum height=0.9cm, align=center},
        arrow/.style={-Stealth, thick, draw=popdark},
        labelbox/.style={rectangle, draw=poppurple, fill=poppurpleL, rounded corners, font=\tiny, align=center}
    }
    % Nœuds principaux
    \node[hub, align=center] (siege) at (0, 3.5){\textbf{SIÈGE / COMMANDEMENT}\\(Cupertino, USA)\\\textit{Conception, Design, iOS,}\\\textit{Marketing, Finance, Juridique}};
    \node[labelbox, align=center] at (0, 4.8){\textbf{VALEUR AJOUTÉE MAXIMALE}\\(Propriété intellectuelle, Marque)};
    % Composants (Est)
    \node[comp, align=center] (proc) at (-4.5, 1.5){Processeurs\\(TSMC, Taïwan)};
    \node[comp, align=center] (ecran) at (-2, 1.5){Écrans OLED\\(Samsung/LG, Corée)};
    \node[comp, align=center] (capteurs) at (0.5, 1.5){Capteurs / Caméras\\(Sony, Japon)};
    \node[comp, align=center] (autres) at (3, 1.5){Autres puces\\Mémoires (Corée/USA)};
    % Assemblage (Centre)
    \node[prod, align=center] (assemblage) at (-0.5, -0.5){\textbf{ASSEMBLAGE FINAL}\\(Zhengzhou, Chine\\Foxconn)\\\textit{+ Inde (Tamil Nadu)}\\\textit{Brésil}};
    % Distribution (Sud)
    \node[prod, align=center] (dist) at (0, -2.5){\textbf{DISTRIBUTION MONDIALE}\\(Apple Stores, Web,\\Revendeurs agréés)};
    \node[prod, align=center] (cmr) at (3.5, -2.5){\textbf{CAMEROUN}\\(Douala, Yaoundé\\Revendeurs, Market)};
    % Flèches
    \draw[arrow] (siege) -- node[left, font=\tiny] {Spécifs, Ordres} (proc);
    \draw[arrow] (siege) -- (ecran);
    \draw[arrow] (siege) -- (capteurs);
    \draw[arrow] (siege) -- (autres);
    \draw[arrow] (proc) -- node[above, font=\tiny] {Composants} (assemblage);
    \draw[arrow] (ecran) -- (assemblage);
    \draw[arrow] (capteurs) -- (assemblage);
    \draw[arrow] (autres) -- (assemblage);
    \draw[arrow, very thick] (assemblage) -- node[left, font=\tiny] {Produits finis} (dist);
    \draw[arrow] (dist) -- node[above, font=\tiny] {Export} (cmr);
    \draw[arrow, dashed, draw=poppink] (cmr) .. controls (5,0) and (5,3) .. node[right, font=\tiny, align=center] {Recyclage /\\Économie circulaire} (siege);
    % Légende intégrée
    \node[draw=popline, fill=white, rounded corners, anchor=south east, font=\tiny, align=center] at (6.5, -3){
        \textbf{Légende :}\\
        \textcolor{popblue}{$\blacksquare$} Commandement / R\&D (USA)\\
        \textcolor{popgreen}{$\blacksquare$} Composants High-Tech (Asie du Nord-Est)\\
        \textcolor{poporange}{$\blacksquare$} Assemblage / Manufacturing (Chine/Inde)\\
        \textcolor{popmuted}{$\blacksquare$} Commercialisation (Monde)\\
        \textcolor{poppink}{$\blacksquare$} Boucle retour (Recyclage)
    };
\end{tikzpicture}
\legende{Schéma simplifié de la chaîne de valeur mondiale d'un smartphone (modèle Apple). La flèche épaisse symbolise le flux principal de valeur physique ; les flèches fines, les flux de composants et d'informations. La flèche rose en pointillés représente la boucle de l'économie circulaire encore marginale.}
\end{popfigure}

\begin{aretenir}[FTN et DIT : le couple indissociable]
\begin{itemize}
    \item La FTN est l'\textbf{acteur principal} de la DIT : elle \textbf{découpe} le processus productif et \textbf{localise} chaque maillon là où le rapport productivité/coût est optimal.
    \item Deux stratégies majeures : \textbf{délocalisation verticale} (sous-traitance en cascade) et \textbf{délocalisation horizontale} (usines « jumelles » pour servir des marchés régionaux).
    \item La FTN organise des \textbf{chaînes de valeur mondiales (CVM)} dont elle capte la rente (propriété intellectuelle, marque, finance) aux deux extrémités (\emph{smiling curve}).
\end{itemize}
\end{aretenir}

\courssub{FTN au Cameroun : ancrage local, stratégie globale}

Le Cameroun accueille de nombreuses FTN, notamment dans les secteurs des télécoms, de l'énergie, de l'agro-industrie, de la logistique portuaire et de la grande distribution. Leur présence illustre la \textbf{compétition pour l'accès aux marchés} (27 millions de consommateurs, hub CEMAC) et aux \textbf{ressources} (pétrole, bois, cacao, coton).

\begin{center}
\begin{tabularx}{\linewidth}{|l|X|X|X|}
\hline
\rowcolor{popblueL}
\textbf{FTN (Siège)} & \textbf{Secteur} & \textbf{Implantation / Activité au Cameroun} & \textbf{Illustration de la DIT} \\
\hline
\textbf{MTN} (Afrique du Sud) & Télécoms & Leader mobile (part de marché $\approx 50\%$), Mobile Money (MoMo), fibre optique. & Siège régional à Johannesburg ; R\&D en Afrique du Sud/Europe ; infrastructure locale (antennes, data centers) ; services financiers innovants adaptés au contexte. \\
\hline
\textbf{Orange} (France) & Télécoms & 2$^{\text{ème}}$ opérateur mobile, Orange Money, Orange Energy, fibre. & Stratégie multi-pays pilotée de Paris ; innovation locale (Orange Digital Center Yaoundé) ; partenariats États/ONG. \\
\hline
\textbf{TotalEnergies} (France) & Pétrole / Énergies & Amont : exploitation offshore (Rio del Rey, Bongkot). Aval : réseau stations, lubrifiants, solaire. & Siège technique Paris ; ingénierie France/UK ; production offshore Cameroun ; raffinage (SONARA, participation) ; distribution locale. \\
\hline
\textbf{Bolloré Logistics / Africa Logistics} (France) & Logistique / Portuaire & \textbf{Concessionnaire terminal conteneurs Port de Douala} (jusqu'à 2024/25 transition), chemin de fer Camrail, logistique hydrocarbures. & \textbf{Intégration verticale} : mer (navires), port (terminaux), rail (Camrail), route, dépôts. Contrôle de la chaîne logistique import/export Afrique Centrale. \\
\hline
\textbf{Nestlé} (Suisse) & Agro-alimentaire & Usines (Douala : Nescafé, Maggi, céréales ; Ngaoundéré : lait). Achats locaux (cacao, café, lait). & Siège Vevey (CH) ; R\&D Suisse/USA ; transformation locale matière première locale ; marques globales adaptées (Maggi cube, Nescafé 3en1). \\
\hline
\textbf{Castel / Sodecoton} (France) & Brasserie / Coton & Brasseries du Cameroun (leader bière/boissons), Sodecoton (filière coton Nord). & Intégration amont-aval : culture coton (petits paysans encadrés) $\to$ égrenage $\to$ export / textile ; brassage local matières premières (maïs, sucre) + marques globales. \\
\hline
\end{tabularx}
\end{center}

\begin{savaistu}[Bolloré et le Port de Douala : un cas d'école d'intégration verticale]
Le consortium \textbf{Bolloré Africa Logistics} (désormais \textbf{MSC - Mediterranean Shipping Company} depuis 2022) a exploité pendant 15 ans le \textbf{terminal à conteneurs du Port de Douala} (concession 2005-2025). Ce terminal traite \textbf{> 90\% du trafic conteneurisé} du Cameroun et sert de porte d'entrée/sortie pour le \textbf{Tchad} et la \textbf{RCA} (pays enclavés). Bolloré gérait aussi \textbf{Camrail} (chemin de fer Douala-Ngaoundéré) et des flottes routières. C'est un exemple parfait de FTN \textbf{intégrée verticalement} qui contrôle la \textbf{chaîne logistique complète} (mer–port–rail–route) pour capter la rente de transit. L'État camerounais a lancé la construction du \textbf{Port en eau profonde de Kribi} (exploité par Kribi Conteneur Terminal, filiale du port de Singapour PSA) pour casser ce monopole et retrouver de la marge de manœuvre.
\end{savaistu}

\begin{exoresolu}[Analyse de la stratégie de Nestlé Cameroun]
\textbf{Énoncé :} À partir de l'exemple de Nestlé, montrer comment une FTN articule ancrage local et stratégie globale.
\textbf{Solution :}
\begin{enumerate}
    \item \textbf{Commandement global :} Siège à Vevey (Suisse) définit la stratégie marques (Nescafé, Maggi, Cerelac), normes qualité, R\&D (centres en Suisse, USA, Chine).
    \item \textbf{Adaptation locale (Glocalisation) :} Produits adaptés au pouvoir d'achat et goûts : \emph{Maggi cube} (bouillon, format petit, prix bas), \emph{Nescafé 3en1} (café+sucre+lait, dose unique), \emph{Cerelac} (céréales infantiles enrichies fer).
    \item \textbf{Ancrage productif :} 2 usines (Douala, Ngaoundéré). Création d'emploi direct ($\approx 1000$) et indirect (distributeurs, transporteurs).
    \item \textbf{Intégration amont (Création de valeur partagée - CSV) :} Programme \emph{Nestlé Cocoa Plan} (cacao durable, plants améliorés, formation planteurs), \emph{Nescafé Plan} (café robusta), collecte lait local (Nord) via centres de collecte + vétérinaires. Sécurise l'approvisionnement, améliore rendements planteurs.
    \item \textbf{Insertion dans la DIT :} Le Cameroun fournit matières premières (cacao, café, lait) + main-d'œuvre transformation ; Nestlé apporte technologie, marques, accès marché global. La valeur ajoutée la plus forte (marketing, R\&D, finance) reste au siège.
\end{enumerate}
\end{exoresolu}

\begin{attention}[Piège : opposer FTN et États de façon absolue]
\textbf{Erreur fréquente :} « Les FTN sont tout-puissantes, les États sont impuissants » OU « L'État doit chasser les FTN pour protéger la souveraineté ».
\textbf{Réalité :} \textbf{Interdépendance asymétrique}.
\begin{itemize}
    \item \textbf{FTN $\to$ État :} Besoin de \textbf{sécurité juridique}, \textbf{infrastructures} (routes, électricité, port), \textbf{main-d'œuvre formée}, \textbf{stabilité politique}, \textbf{accès au marché} (réglementation favorable). Une FTN peut partir si l'environnement se dégrade (coût de sortie vs coût de séjour).
    \item \textbf{État $\to$ FTN :} Besoin d'\textbf{investissements} (IDE), d'\textbf{emplois}, de \textbf{transfert de technologie}, de \textbf{recettes fiscales}, d'\textbf{accès aux marchés mondiaux} via les réseaux de la FTN.
    \item \textbf{Négociation :} L'État \textbf{réglemente} (code investissement, travail, environnement, fiscalité) et \textbf{négocie} (conventions d'établissement). La FTN \textbf{fait du lobbying} pour des conditions favorables. Le \textbf{rapport de force} dépend de l'attractivité du pays (taille du marché, ressources, stabilité) et de la mobilité du capital de la FTN.
\end{itemize}
\end{attention}

\begin{aretenir}[Chiffres clés : la puissance économique des FTN]
\begin{itemize}
    \item Les \textbf{100 plus grandes FTN} réalisent un chiffre d'affaires cumulé \textbf{supérieur au PIB de l'ensemble des pays d'Afrique subsaharienne} (hors Afrique du Sud) réunis.
    \item Environ \textbf{80\% du commerce mondial} est intra-firme (entre maison-mère et filiales ou entre filiales d'une même FTN) ou piloté par des FTN.
    \item Les \textbf{Investissements Directs Étrangers (IDE)} mondiaux (flux annuels) $\approx$ 1300 milliards USD (2023, CNUCED), dont \textbf{moins de 5\% pour l'Afrique} (majoritairement pétrole/mines). Les FTN sont les vecteurs quasi exclusifs des IDE.
\end{itemize}
\end{aretenir}

\courssub{Les métropoles et villes mondiales : les nœuds de commandement du système}

\begin{definition}[Ville mondiale (Métropole mondiale)]
Une \cle{ville mondiale} (ou \emph{global city} selon Saskia Sassen) est une métropole qui concentre des \textbf{fonctions de commandement et de contrôle} de l'économie mondiale : sièges de FTN, places financières, marchés spécialisés (assurance, maritime, matières premières), organisations internationales, pôles de R\&D et d'innovation, hubs aériens/numériques. Elle n'est pas seulement grande par sa population, mais par sa \textbf{connectivité} et son \textbf{pouvoir de décision} sur les réseaux mondiaux.
\end{definition}

\begin{propriete}[La hiérarchie des villes mondiales (Réseau GaWC - Loughborough)]
Le groupe de recherche \textbf{GaWC} classe les villes selon leur \textbf{connectivité} dans le réseau des services aux entreprises (finance, droit, conseil, pub, comptabilité) :
\begin{itemize}
    \item \textbf{Alpha ++ (Top) :} \textbf{Londres}, \textbf{New York} (les deux « capitales du monde » historiques, finance, droit, culture, médias).
    \item \textbf{Alpha + :} \textbf{Shanghai}, \textbf{Tokyo}, \textbf{Dubaï}, \textbf{Singapour}, \textbf{Hong Kong}, \textbf{Paris}, \textbf{Pékin}, \textbf{Sydney}.
    \item \textbf{Alpha / Alpha - :} Francfort, Toronto, Chicago, Séoul, Madrid, Milan, Moscou, São Paulo, Mumbai, Johannesburg (seule ville africaine Alpha -), etc.
    \item \textbf{Bêta / Gamma :} Villes mondiales de second rang (Lyon, Marseille, Casablanca, Lagos, Nairobi, Douala/Yaoundé \emph{au niveau régional}).
\end{itemize}
\end{propriete}

\begin{popfigure}
\begin{tikzpicture}[scale=0.65, every node/.style={font=\small}]
    % Fond planisphère très simplifié (blocs continents)
    \fill[popcream] (-10,-5) rectangle (15,6); % Océans
    % Amériques
    \fill[popgreenL] (-10,-5) .. controls (-8,-2) and (-6,1) .. (-7,4) .. controls (-9,5) .. (-10,6) -- (-10,-5) -- cycle;
    \fill[popgreenL] (-8,-5) .. controls (-6,-3) .. (-4,-5) -- cycle; % Am Sud
    % Afrique
    \fill[poporangeL] (0,-5) .. controls (1,-2) and (3,1) .. (2,4) .. controls (0,5) .. (-1,4) .. controls (-2,1) .. (0,-5) -- cycle;
    % Europe
    \fill[popblueL] (1,3) .. controls (3,4) and (5,4) .. (6,3) .. controls (5,2) .. (3,2) -- (1,3) -- cycle;
    % Asie
    \fill[poppurpleL] (6,0) .. controls (8,2) and (12,3) .. (13,1) .. controls (12,-1) .. (8,-2) .. controls (6,-1) .. (6,3) -- cycle;
    % Océanie
    \fill[poppinkL] (12,-4) ellipse ({1.5} and {1});
    % Villes Mondiales (Alpha ++)
    \node[circle, fill=popink, draw=white, thick, minimum size=6pt] (ny) at (-7.5, 2.5) {};
    \node[above=1pt of ny, font=\tiny\bfseries, text=popink] {New York};
    \node[circle, fill=popink, draw=white, thick, minimum size=6pt] (ldn) at (2.5, 3.5) {};
    \node[above=1pt of ldn, font=\tiny\bfseries, text=popink] {Londres};
    % Villes Alpha +
    \foreach \pos/\name/\anchor in {
        (9.5, 1.5)/Tokyo/left,
        (8.5, 0.5)/Shanghai/left,
        (5.5, 1.5)/Dubaï/above,
        (7.5, 2.5)/Pékin/above,
        (6.5, 3.5)/Séoul/above,
        (6.5, -0.5)/Singapour/below,
        (6.5, 0.5)/Hong Kong/below,
        (3.5, 3.5)/Paris/above,
        (4.5, -1)/Johannesburg/below,
        (-5.5, -2.5)/São Paulo/right,
        (12, -3.5)/Sydney/left} {
        \node[circle, fill=poporange, draw=white, thick, minimum size=5pt] at \pos {};
        \node[\anchor=1pt of \pos, font=\tiny, text=poporange] {\name};
    }
    % Places financières majeures (étoiles)
    \foreach \pos in {(-7.5,2.5), (2.5,3.5), (9.5,1.5), (8.5,0.5), (5.5,1.5), (7.5,2.5), (3.5,3.5), (6.5,0.5), (-5.5,-2.5)} {
        \node[star, star points=5, star point ratio=2.25, fill=popgold, draw=popdark, thick, minimum size=5pt, inner sep=0] at \pos {};
    }
    % Flux financiers principaux (flèches courbes)
    \draw[popgold, thick, -Stealth, bend left=30] (ny) to (ldn);
    \draw[popgold, thick, -Stealth, bend left=30] (ldn) to (ny);
    \draw[popgold, thick, -Stealth, bend left=20] (ldn) to (8.5,0.5); % Ldn -> Shanghai
    \draw[popgold, thick, -Stealth, bend left=20] (ny) to (5.5,1.5); % NY -> Dubai
    \draw[popgold, thick, -Stealth, bend right=20] (5.5,1.5) to (8.5,0.5); % Dubai -> Shanghai
    \draw[popgold, thick, -Stealth, bend left=15] (2.5,3.5) to (4.5,-1); % Ldn -> Joburg
    \draw[popgold, thick, -Stealth, bend right=15] (-7.5,2.5) to (-5.5,-2.5); % NY -> Sao Paulo
    % Légende
    \node[draw=popline, fill=white, rounded corners, anchor=south west, font=\tiny, align=left] at (-9.5, -4.5) {
        \textbf{Légende :}\\
        \textcolor{popink}{\Large $\bullet$} \textbf{Alpha ++} (Commandement suprême)\\
        \textcolor{poporange}{\Large $\bullet$} \textbf{Alpha +} (Commandement régional/secteur)\\
        \textcolor{popgold}{\Large $\star$} \textbf{Places financières majeures}\\
        \textcolor{popgold}{\rightarrow} Flux financiers dominants\\
        \textit{Base : Réseau GaWC / Classements financiers (GFCI)}
    };
    % Nord
    \node[draw=popdark, fill=popcream, circle, minimum size=1cm] at (-9, 5) {\textbf{N}};
\end{tikzpicture}
\legende{Carte schématique des principales villes mondiales (classement GaWC) et places financières. Les flèches dorées illustrent les circuits financiers dominants (New York $\leftrightarrow$ Londres comme cœur, extensions vers Asie et marchés émergents). Johannesburg est la seule ville africaine classée Alpha -. Douala/Yaoundé sont des métropoles \emph{régionales} (niveau Gamma/Sufficiency), non mondiales.}
\end{popfigure}

\begin{exemplebox}[Spécialisation des places financières mondiales]
\begin{itemize}
    \item \textbf{New York (Wall Street) :} Marché d'actions (NYSE, Nasdaq), capital-investissement, hedge funds, dollar US (monnaie de réserve).
    \item \textbf{Londres (City) :} Marché des changes (Forex $\approx 40\%$ mondial), euro-obligations, assurance (Lloyd's), droit commercial international, finance verte.
    \item \textbf{Tokyo :} Marché obligataire japonais, banque commerciale géante, porte d'entrée capitaux asiatiques.
    \item \textbf{Shanghai / Hong Kong / Shenzhen :} Internationalisation du Yuan (RMB), introduction en bourse tech chinoises, hub Asie-Pacifique.
    \item \textbf{Dubaï (DIFC) / Singapour :} Hubs offshore/onshore, gestion de fortune (private banking), finance islamique (Dubaï), droit common law, fiscalité attractive $\to$ \textbf{paradis fiscaux / places offshore}.
\end{itemize}
\end{exemplebox}

\begin{savaistu}[Paradis fiscaux et optimisation : la fuite des capitaux]
Les \textbf{paradis fiscaux} (ou \emph{places financières offshore}) — îles Caïmans, Bermudes, Luxembourg, Irlande, Pays-Bas, Singapour, Delaware (USA) — offrent \textbf{impôt sur les sociétés $\approx 0\%$}, \textbf{secret bancaire} (ou opacité bénéficiaires effectifs) et \textbf{absence d'activité économique réelle} (boîtes aux lettres). Les FTN y logent leurs \textbf{propriétés intellectuelles} (brevets, marques) et font payer des \textbf{redevances} par leurs filiales dans les pays à forte imposition (France, Cameroun, etc.), déplaçant artificiellement les bénéfices (\emph{profit shifting}). \textbf{Conséquence :} L'Afrique perd \textbf{$\approx$ 50 à 80 milliards USD/an} (fuite de capitaux illicites, PNUD/CEA), soit plus que l'aide publique au développement reçue. L'OCDE/G20 (projet BEPS, impôt minimum mondial 15\%) tente de contrer cela, mais la résistance est forte.
\end{savaistu}

\courssub{Le rôle des États : souveraineté mise à l'épreuve, mais acteurs stratégiques}

\begin{definition}[Souveraineté dans la mondialisation]
La \cle{souveraineté} étatique (pouvoir exclusif sur un territoire : lois, monnaie, frontières, force publique) est \textbf{relativisée} par la mondialisation : mobilité des capitaux (échappent à l'impôt), normes internationales (OMC, accords de libre-échange, traités d'investissement) qui limitent la marge de manœuvre législative, pouvoir des agences de notation (Standard \& Poor's, Moody's) sur le coût de la dette. Cependant, l'État reste \textbf{l'acteur indispensable} : seul garant de l'ordre juridique, de la sécurité, des infrastructures lourdes, de la cohésion sociale, de la régulation environnementale.
\end{definition}

\begin{methode}[Analyser l'action de l'État face à la mondialisation (Grille d'observation)]
\begin{enumerate}
    \item \textbf{Politiques d'attractivité (Offre) :} Zones économiques spéciales (ZES), exonérations fiscales, guichets uniques, formation ciblée, infrastructures dédiées (ex: Zone industrielle de Kribi).
    \item \textbf{Régulation / Protection (Demande) :} Droit du travail, code de l'environnement, clauses de contenu local (ex: \% main-d'œuvre locale, transformation sur place), contrôle des IDE stratégiques (sécurité nationale).
    \item \textbf{Négociation internationale :} Accords commerciaux (ALECA, ZLECAf, EPA UE), traités bilatéraux d'investissement, participation OMC, G20, ONU.
    \item \textbf{Grand projets d'infrastructure :} Ports, corridors, énergie, numérique (backbone fibre) pour s'insérer dans les réseaux.
    \item \textbf{Redistribution / Cohésion :} Fiscalité prélevée sur les gagnants (FTN, secteurs exportateurs) pour financer services publics, filets sociaux, aménagement du territoire.
\end{enumerate}
\end{methode}

\begin{exemplebox}[Le Port en eau profonde de Kribi (Cameroun) : l'État comme acteur d'infrastructure stratégique]
\begin{itemize}
    \item \textbf{Contexte :} Port de Douala saturé (tirant d'eau max 11m, congestion, monopole terminal conteneurs).
    \item \textbf{Projet État :} Port de Kribi (Phase 1 livrée 2018, Phase 2 en cours). Tirant d'eau \textbf{16--18m} $\to$ accueille navires \textbf{Post-Panamax / New Panamax} (capacité > 10 000 EVP).
    \item \textbf{Partenariat Public-Privé (PPP) :} Concession terminal conteneurs à \textbf{Kribi Conteneur Terminal (KCT)} — filiale de \textbf{PSA International (Singapour)}, 1$^{\text{er}}$ opérateur portuaire mondial. Concession mine de fer Mbalam-Nabeba (Sundance/Australie puis AustSino/Chine) $\to$ terminal minéralier dédié.
    \item \textbf{Corridor :} Route Kribi-Lolodorf-Ebolowa-Yaoundé + projet ferroviaire (Kribi-Mbalam) $\to$ désenclavement Sud \& Est + export minerais.
    \item \textbf{Objectif souverain :} \textbf{Réduire la dépendance à Douala/Bolloré}, capter trafic transit Tchad/RCA, créer pôle industriel (ZES Kribi : aluminium, engrais, gaz/GNL), affermir la \textbf{souveraineté portuaire} et négocier en position de force avec les FTN logistiques.
\end{itemize}
\end{exemplebox}

\begin{attention}[Le piège de la « course au fond » (Race to the bottom)]
Pour attirer les FTN, les États (surtout en développement) sont tentés de \textbf{baisser les standards} : fiscalité (exonérations totales 10-20 ans), droit du travail (zones franches sans syndicats), environnement (tolérance pollution). \textbf{Résultat :} \textbf{Moins de recettes fiscales}, \textbf{précarisation de l'emploi}, \textbf{dégradation écologique}, \textbf{absence de transfert de technologie} (usines « boîtes noires » : entrée matières + main-d'œuvre, sortie produit, pas de R\&D locale). \textbf{Leçon :} L'attractivité durable repose sur la \textbf{qualité} (infrastructures, main-d'œuvre qualifiée, stabilité, État de droit) et non sur le \textbf{coût bas} seul. La \textbf{ZLECAf} (Zone de libre-échange continentale africaine) vise à mutualiser les marchés pour renforcer le pouvoir de négociation des États africains face aux FTN.
\end{attention}

\courssub{Acteurs de régulation et de contestation : le contre-pouvoir mondial}

La mondialisation n'est pas un jeu à somme nulle sans règles. Des acteurs tentent de la \textbf{réguler}, la \textbf{corriger} ou la \textbf{contester}.

\begin{center}
\begin{tabularx}{\linewidth}{|l|X|X|}
\hline
\rowcolor{popblueL}
\textbf{Type d'acteur} & \textbf{Exemples majeurs} & \textbf{Rôle / Levier d'action} \\
\hline
\textbf{Organisations internationales (Régulation formelle)} & \textbf{OMC} (commerce), \textbf{FMI / BM} (finance, dette, politiques structurelles), \textbf{OCDE} (normes fiscales BEPS, lignes directrices FTN), \textbf{OIT} (normes travail), \textbf{ONU / CNUCED} (développement, IDE), \textbf{OMI} (transport maritime), \textbf{Union Africaine / CEMAC} (régional). & Fixation de \textbf{règles contraignantes ou souples} (soft law), surveillance, conditionnalité aide/prêts, arbitrage différends (OMC, CIRDI). \\
\hline
\textbf{ONG internationales (Plaidoyer, expertise, terrain)} & \textbf{Oxfam, Transparency International, Greenpeace, Amnesty International, WWF, ActionAid, Publiez Ce Que Vous Payez (PCQVP).} & Enquêtes, rapports, campagnes médiatiques, lobbying États/OMC/OCDE, actions juridiques, soutien communautés locales, certification (label commerce équitable, FSC bois, RSPO huile de palme). \\
\hline
\textbf{Mouvements altermondialistes / Sociaux (Contestation, alternatives)} & \textbf{Forum Social Mondial (FSM)}, \textbf{ATTAC}, \textbf{Via Campesina} (paysans), \textbf{Syndicats internationaux (CSI)}, \textbf{Mouvements climat (Fridays for Future, Extinction Rebellion)}, \textbf{Collectifs locaux anti-mines/barrages/accaparement terres}. & Manifestations (contre-sommets G7/G20/OMC), désobéissance civile, construction d'alternatives (économie solidaire, circuits courts, agroécologie, monnaies locales), pression pour \textbf{justice fiscale, climatique, sociale}. \\
\hline
\textbf{Acteurs hybrides / Multipartites} & \textbf{Global Compact (ONU)}, \textbf{Initiative pour la Transparence dans les Industries Extractives (ITIE)}, \textbf{Principes directeurs ONU Entreprises et Droits de l'Homme}, \textbf{Label B-Corp}. & \textbf{Autorégulation volontaire} (souvent critiquée comme \emph{greenwashing/socialwashing}), reporting extra-financier (ESG), diligence raisonnable (devoir de vigilance - loi France 2017, directive UE CS3D 2024). \\
\hline
\end{tabularx}
\end{center}

\begin{exoresolu}[Le devoir de vigilance : un nouvel outil juridique contre l'impunité des FTN]
\textbf{Énoncé :} Expliquer en quoi la loi française sur le devoir de vigilance (2017) et la directive européenne CS3D (2024) changent la donne pour les FTN opérant au Cameroun.
\textbf{Solution :}
\begin{enumerate}
    \item \textbf{Principe :} Obligation pour les grandes entreprises (siège en France/UE, seuils effectifs/chiffre d'affaires) d'élaborer, publier et mettre en œuvre un \textbf{plan de vigilance} identifiant les risques graves (droits humains, santé, environnement) liés à \textbf{l'ensemble de leur chaîne de valeur} (filiales, sous-traitants, fournisseurs directs/indirects).
    \item \textbf{Sanction :} Responsabilité civile : réparation du préjudice si absence de plan ou mise en œuvre défaillante ayant causé un dommage. Juridiction du siège social (Paris) compétente.
    \item \textbf{Impact au Cameroun :} Une FTN française (TotalEnergies, Orange, Bolloré/MSC, Castel, etc.) \textbf{doit} surveiller ses fournisseurs camerounais (ex: planteurs cacao/café, sous-traitants BTP, sécurité privée, transport). Si pollution, travail forcé, expulsion illégale $\to$ victime peut saisir juge français.
    \item \textbf{Limites :} Preuve du lien de causalité difficile, coût procédure, seules très grandes entreprises concernées, pas de sanction pénale. Mais \textbf{effet de levier majeur} : les FTN imposent clauses contractuelles, audits, formations à leurs partenaires camerounais pour se protéger juridiquement.
\end{enumerate}
\end{exoresolu}

\begin{aretenir}[Synthèse : L'architecture des acteurs de la mondialisation]
\begin{center}
\begin{tikzpicture}[scale=0.9, every node/.style={font=\small}]
    \tikzset{
        actor/.style={rectangle, rounded corners, draw=popblue, thick, fill=popblueL, minimum width=3cm, minimum height=1.5cm, align=center},
        flux/.style={-Stealth, thick, draw=popdark},
        doubleflux/.style={<->, thick, draw=poporange},
    }
    % Nœuds
    \node[actor, align=center] (ftn) at (0, 3){\textbf{FIRMES TRANSNATIONALES (FTN)}\\\textit{Moteurs : Investissement, Technologie,}\\\textit{Marque, Chaînes de valeur mondiales}};
    \node[actor, align=center] (met) at (6, 3){\textbf{MÉTROPOLES MONDIALES}\\\textit{Nœuds : Finance, Sièges, R\&D, Hubs,}\\\textit{Commandement réseau (NY, Londres, Shanghai...)}};
    \node[actor, align=center] (etat) at (3, 0){\textbf{ÉTATS}\\\textit{Régulateurs : Lois, Infrastructures,}\\\textit{Fiscalité, Négociation, Cohésion,}\\\textit{Souveraineté relative}};
    \node[actor, align=center] (reg) at (3, -2.5){\textbf{RÉGULATION / CONTESTATION}\\\textit{OMC, FMI, OCDE, ONU, ONG, Syndicats,}\\\textit{Mouvements altermondialistes,}\\\textit{Devoir de vigilance, Justice climatique}};
    % Relations
    \draw[doubleflux] (ftn) -- node[above, font=\tiny] {Implantation / Lobbying / Ordres} (etat);
    \draw[doubleflux] (met) -- node[right, font=\tiny] {Sièges / Finance / Talents} (ftn);
    \draw[doubleflux] (met) -- node[above right, font=\tiny] {Places financières / Normes} (reg);
    \draw[flux] (reg) -- node[left, font=\tiny] {Règles / Sanctions / Pression} (ftn);
    \draw[flux] (reg) -- node[left, font=\tiny] {Normes / Conditionnalité} (etat);
    \draw[flux] (etat) -- node[below left, font=\tiny] {Infrastructures / Marché / Main-d'œuvre} (ftn);
    \draw[flux, bend left=15] (ftn) to node[above, font=\tiny] {Flux IDE, Tech, Emplois} (met);
    \draw[flux, bend right=15] (met) to node[below, font=\tiny] {Services aux entreprises} (ftn);
\end{tikzpicture}
\end{center}
\legende{Schéma des relations entre les quatre pôles d'acteurs. Les doubles flèches indiquent des relations d'interdépendance forte ; les flèches simples, des relations d'autorité, de contrainte ou de service.}
\end{aretenir}

\begin{exercice}[Application : Acteurs et territoires au Cameroun]
\begin{enumerate}
    \item En vous appuyant sur l'exemple du \textbf{Port de Kribi}, expliquez comment l'État camerounais tente de \textbf{reprendre la main} sur la chaîne logistique face à la FTN Bolloré/MSC. Citez au moins trois leviers d'action étatique.
    \item Une FTN agro-industrielle (type Socapalm/Safacam — groupe Bolloré ou Wilmar/Olam) veut étendre sa plantation de palmiers à huile dans la région du Sud. Identifiez \textbf{quatre acteurs de régulation/contestation} (internationaux, nationaux, locaux) qui peuvent intervenir et précisez pour chacun son levier d'action.
    \item \textbf{Croquis :} Réalisez un croquis simple (schéma) montrant la \textbf{Division Internationale du Travail} pour le \textbf{cacao camerounais} (de la parcelle du paysan à la tablette de chocolat en Europe). Légendez : acteurs, lieux, flux, valeur ajoutée.
\end{enumerate}
\end{exercice}

\begin{corrige}[Exercice n°5 - Acteurs et territoires]
\begin{enumerate}
    \item \textbf{Port de Kribi — Leviers étatiques :}
    \begin{itemize}
        \item \textbf{Infrastructure publique :} Financement/construction du port en eau profonde (tirant d'eau 18m) et des corridors routier/ferroviaire (biens publics que la FTN ne ferait pas).
        \item \textbf{Cadre juridique / Concession :} Choix d'un \textbf{autre} opérateur mondial (PSA Singapore) pour le terminal conteneurs $\to$ \textbf{concurrence} / fin de monopole Bolloré à Douala.
        \item \textbf{ZES (Zone Économique Spéciale) Kribi :} Cadre fiscal/réglementaire incitatif pour attirer des industries de transformation (aluminium, engrais, GNL) $\to$ création de valeur locale, pas seulement transit.
        \item \textbf{Souveraineté portuaire :} L'État (via Port Autonome de Kribi - PAK) reste propriétaire des infrastructures, maître d'ouvrage, régulateur des tarifs et d'accès.
    \end{itemize}
    \item \textbf{Acteurs régulation/contestation plantation palmier :}
    \begin{tabularx}{\linewidth}{|l|X|X|}
    \hline
    \textbf{Acteur} & \textbf{Type} & \textbf{Levier} \\
    \hline
    \textbf{RSPO (Table ronde huile de palme durable)} & ONG/Initiative multipartite internationale & Certification durable (zéro déforestation, droits travailleurs, communautés). Refus certification = perte marchés UE/USA. \\
    \hline
    \textbf{MINEPDED / MINFOF (Ministères Environnement/Forêts)} & État national (Cameroun) & Étude d'impact environnemental (EIES), titre de concession, contrôle respect cahier des charges, sanctions administratives. \\
    \hline
    \textbf{Communautés villageoises / Chefferies / ONG locales (ex: CED, RELUFA)} & Acteurs locaux / Société civile & Droit au consentement libre, préalable et éclairé (CLPE), recours judiciaires (tribunaux), médiatisation, blocage physique, négociation redevances/terres. \\
    \hline
    \textbf{Banques de développement / Investisseurs (ex: BAD, Proparco, banques commerciales)} & Acteurs financiers & \textbf{Conditionnalité ESG} : refus financement si non-conformité normes environnementales/sociales (Performance Standards SFI). \\
    \hline
    \end{tabularx}
    \item \textbf{Croquis Cacao Cameroun $\to$ Chocolat Europe (Description pour réalisation) :}
    \begin{itemize}
        \item \textbf{Espace 1 (Cameroun - Sud/Centre/Est) :} \textcolor{popgreen}{Parcelles paysannes} (petits exploitants, 0,5--2 ha, rendement faible $\approx$ 300 kg/ha). \textcolor{poporange}{Collecte/Égrenage} (coopératives, pisteurs, exportateurs locaux type Telcar/Cicam/Olam). \textbf{Export fèves brutes} (90\%+), port Douala/Kribi. \textit{Valeur ajoutée : Faible (prix bord champ $\approx$ 1000--1500 FCFA/kg)}.
        \item \textbf{Flux :} Navire vrac $\to$ \textbf{Ports européens} (Amsterdam, Anvers, Hambourg).
        \item \textbf{Espace 2 (Europe - Allemagne, Pays-Bas, France, Belgique, Suisse) :} \textcolor{popblue}{Transformation industrielle} (torréfaction, broyage, pressage $\to$ beurre/poudre de cacao) — \textbf{Barry Callebaut, Cargill, Olam, Cémoi}. \textcolor{poppurple}{Fabrication chocolat} (industriels : Nestlé, Lindt, Ferrero, Mars, Mondelez + artisans). \textcolor{poppink}{Distribution/Conso}. \textit{Valeur ajoutée : Élevée (transformation, marque, marketing, R\&D, finance)}.
        \item \textbf{Acteurs transversaux :} \textbf{Bourses} (Londres ICE, New York NYSE - cours mondial), \textbf{Spéculateurs}, \textbf{Certificateurs} (Rainforest Alliance, Fairtrade, Bio), \textbf{ONG} (suivi travail enfants, déforestation), \textbf{États} (réglementation déforestation importée UE - EUDR).
    \end{itemize}
    \textit{Flèche de valeur :} Large flèche « Fèves brutes (Prix bas) » vers Nord ; Fine flèche « Revenus paysans / Taxes » vers Sud. Annotation : « 70\% de la valeur finale captée en aval (transformation + marque) ».
\end{enumerate}
\end{corrige}

\coursec{TP guidé : croquis et diagrammes des mécanismes de la mondialisation}

Dans la section précédente, nous avons identifié les \cle{acteurs} qui font fonctionner la mondialisation : les \cle{FTN}, les \cle{métropoles} et les \cle{États}. Mais au Baccalauréat technique, il ne suffit pas de connaître ces acteurs : il faut savoir les \emph{représenter} et les \emph{exploiter}. C'est l'objet de ce travail pratique guidé : transformer les savoirs de l'unité en productions graphiques (croquis, diagrammes, schémas) et en raisonnements rédigés, conformément aux savoir-faire officiels : \emph{appliquer les notions de l'unité à des situations concrètes} et \emph{rédiger et justifier une démarche, une réponse ou une conclusion}.

\courssub{Rappels de méthode : les règles d'or du croquis et du diagramme}

Avant de produire, rappelons les quatre exigences communes à toute production graphique en géographie.

\begin{aretenir}[Les quatre règles d'or]
\begin{enumerate}
\item Un \cle{titre} précis, placé au-dessus du document, qui reprend le sujet, l'espace et l'échelle.
\item Une \cle{légende} complète, organisée et hiérarchisée (du plus important au moins important), placée sous le document.
\item Une \cle{orientation} (flèche du nord) et, si possible, une échelle indicative sur les croquis.
\item La \cle{proportionnalité} des figurés : l'épaisseur d'une flèche ou la hauteur d'une barre doit être proportionnelle à la valeur représentée.
\end{enumerate}
\end{aretenir}

\begin{attention}[Erreur fréquente à l'examen]
Un croquis \textbf{sans légende ni titre} est \textbf{sanctionné systématiquement} au Baccalauréat technique, même si le dessin est correct : il peut perdre jusqu'à la moitié des points. De même, une flèche dont l'épaisseur ne correspond pas à l'importance du flux est considérée comme une erreur de méthode. Vérifie toujours : titre, légende, nord, proportionnalité.
\end{attention}

\begin{methode}[Réussir un croquis géographique au Baccalauréat technique : 5 étapes]
\begin{enumerate}
\item \textbf{Analyser la consigne} : souligner le sujet, l'espace concerné, l'échelle et ce qu'il faut montrer (flux, acteurs, inégalités). En déduire le titre définitif.
\item \textbf{Tracer le fond de carte} : reproduire rapidement et simplement les contours (continents, frontières, littoraux) au crayon, sans détails inutiles ; placer la flèche du nord.
\item \textbf{Placer les figurés} : localiser d'abord les éléments principaux (ports, métropoles, routes), puis les flux par des flèches d'épaisseur proportionnelle aux volumes.
\item \textbf{Construire la légende} : chaque figuré utilisé sur la carte doit apparaître dans la légende, et rien d'autre ; ordonner la légende logiquement.
\item \textbf{Titrer et vérifier} : écrire le titre au-dessus, relire la consigne et contrôler que tout ce qui était demandé est visible sur le croquis.
\end{enumerate}
\end{methode}

\courssub{TP 1 — Croquis : les grandes routes maritimes et les places portuaires mondiales}

\textbf{Consigne.} Sur le fond de carte muet du monde fourni, représente les grandes routes maritimes mondiales, les détroits et canaux stratégiques, ainsi que les principaux ports mondiaux, dont Douala et Kribi. N'oublie ni le titre, ni la légende, ni le nord.

\textbf{Travail préparatoire (analyse de la consigne).} Trois familles de figurés sont attendues : des \emph{flèches} pour les routes maritimes (plus épaisses sur les axes majeurs), des \emph{points} pour les ports (plus gros pour les plus grandes places portuaires), des \emph{symboles} pour les détroits et canaux. Les routes majeures à retenir : la route \cle{Europe -- Amérique du Nord} (Atlantique Nord), la route \cle{Europe -- Asie} par le canal de Suez et l'océan Indien, la route \cle{Asie -- Amérique du Nord} (Pacifique Nord), et la route du golfe de Guinée qui dessert les ports camerounais.

\begin{popfigure}
\begin{center}
\begin{tikzpicture}[scale=1.0]
% Océans (fond)
\fill[popblueL!45] (-7.2,-3.2) rectangle (7.2,3.2);
% Amérique du Nord
\fill[popcream] (-6.8,0.4) -- (-5.6,2.9) -- (-3.4,3.0) -- (-2.6,1.9) -- (-3.1,0.9) -- (-4.4,0.3) -- (-5.6,-0.2) -- cycle;
% Amérique du Sud
\fill[popcream] (-4.4,-0.2) -- (-3.0,-0.3) -- (-2.4,-1.3) -- (-2.9,-2.9) -- (-3.6,-2.8) -- (-4.2,-1.3) -- cycle;
% Europe
\fill[popcream] (-0.6,1.5) -- (0.9,2.9) -- (2.0,2.7) -- (1.7,1.7) -- (0.5,1.3) -- cycle;
% Afrique
\fill[popcream] (-0.7,1.2) -- (1.5,1.2) -- (2.0,0.2) -- (1.5,-1.5) -- (0.8,-2.7) -- (0.1,-2.6) -- (-0.5,-0.9) -- cycle;
% Asie
\fill[popcream] (2.1,2.6) -- (5.6,2.9) -- (6.6,1.6) -- (5.9,0.6) -- (4.9,0.9) -- (4.2,-0.2) -- (3.6,0.7) -- (2.6,0.5) -- (2.1,1.5) -- cycle;
% Asie du Sud-Est
\fill[popcream] (4.6,0.5) -- (5.3,0.4) -- (5.2,-0.6) -- (4.7,-0.6) -- cycle;
% Australie
\fill[popcream] (4.9,-1.4) -- (6.3,-1.3) -- (6.4,-2.3) -- (5.0,-2.4) -- cycle;
% Routes maritimes (flèches d'épaisseur proportionnelle)
\draw[line width=2.2pt,poporange,->] (-2.5,1.9) -- (-0.8,1.9);
\draw[line width=2.2pt,poporange,->] (0.4,1.3) -- (1.1,0.9) -- (2.3,0.4) -- (3.6,0.1) -- (4.9,0.3);
\draw[line width=2.2pt,poporange,->] (5.6,1.3) .. controls (4.2,2.2) and (-1.5,2.6) .. (-3.0,2.2);
\draw[line width=1.2pt,poporange,->] (0.2,0.6) -- (-0.1,-0.4) -- (-0.3,-1.6);
\draw[line width=1.2pt,poporange,->] (-0.3,-0.6) -- (-2.6,-0.8);
% Détroits et canaux
\node[star,star points=6,fill=popink,inner sep=1.4pt] at (1.15,0.85) {};
\node[star,star points=6,fill=popink,inner sep=1.4pt] at (5.05,0.15) {};
\node[star,star points=6,fill=popink,inner sep=1.4pt] at (-3.05,0.55) {};
% Ports
\fill[popink] (-2.4,1.75) circle (2.4pt); \node[above,font=\scriptsize] at (-2.4,1.8) {New York};
\fill[popink] (-0.55,1.55) circle (2.4pt); \node[above,font=\scriptsize] at (-0.55,1.6) {Rotterdam};
\fill[popink] (5.35,1.15) circle (2.4pt); \node[above,font=\scriptsize] at (5.35,1.2) {Shanghai};
\fill[popink] (4.85,0.35) circle (2.4pt); \node[below,font=\scriptsize] at (4.85,0.3) {Singapour};
\fill[popgreen] (0.35,0.15) circle (2.2pt); \node[right,font=\scriptsize\bfseries] at (0.42,0.15) {Douala};
\fill[popgreen] (0.15,-0.25) circle (2.2pt); \node[right,font=\scriptsize\bfseries] at (0.22,-0.3) {Kribi};
% Nord
\draw[->,thick] (-6.9,-2.2) -- (-6.9,-1.5); \node[above,font=\scriptsize] at (-6.9,-1.5) {N};
% Titre
\node[font=\small\bfseries,align=center] at (0,3.7) {Les grandes routes maritimes et les places portuaires mondiales};
\end{tikzpicture}
\end{center}
\legende{Croquis modèle (très simplifié). Flèches orange : routes maritimes (épaisseur proportionnelle au trafic) ; points noirs : grandes places portuaires mondiales ; points verts : ports camerounais de Douala et Kribi ; étoiles : passages stratégiques (canal de Suez, détroit de Malacca, canal de Panama). Échelle indicative : carte mondiale schématique.}
\end{popfigure}

\textbf{Correction collective.} Le croquis est projeté au tableau et comparé au modèle. Chaque élève s'auto-évalue à l'aide de la grille suivante.

\begin{center}
\begin{tabularx}{\linewidth}{|X|c|c|}
\hline
\rowcolor{popblueL} \textbf{Critère d'auto-évaluation} & \textbf{Réussi} & \textbf{À revoir} \\
\hline
Titre complet (sujet + espace + échelle) au-dessus du croquis & $\square$ & $\square$ \\
\hline
Fond de carte lisible et flèche du nord présente & $\square$ & $\square$ \\
\hline
Au moins 3 routes maritimes majeures correctement tracées & $\square$ & $\square$ \\
\hline
Épaisseur des flèches proportionnelle à l'importance du flux & $\square$ & $\square$ \\
\hline
Au moins 4 grands ports placés, dont Douala et Kribi & $\square$ & $\square$ \\
\hline
Détroits/canaux stratégiques localisés (Suez, Malacca, Panama) & $\square$ & $\square$ \\
\hline
Légende complète, organisée, sans élément absent de la carte & $\square$ & $\square$ \\
\hline
\end{tabularx}
\end{center}

\begin{savaistu}[Le port de Kribi, futur hub de l'Afrique centrale ?]
Inauguré en 2018, le port en eau profonde de \cle{Kribi} (région du Sud) peut accueillir de très grands navires, contrairement au port de Douala encombré et ensablé. Il ambitionne de devenir un \cle{hub régional} pour l'Afrique centrale, desservant aussi le Tchad et la Centrafrique. C'est un exemple parfait de la manière dont un État (le Cameroun) s'insère dans les flux de la mondialisation en investissant dans les infrastructures portuaires.
\end{savaistu}

\courssub{TP 2 — Diagramme en barres : les principaux partenaires commerciaux du Cameroun}

\textbf{Consigne.} À partir du tableau de données ci-dessous (parts des principaux partenaires dans les exportations du Cameroun, en \%), construis un diagramme en barres avec titre, axes légendés et barres proportionnelles.

\begin{center}
\begin{tabularx}{\linewidth}{|l|X|}
\hline
\rowcolor{popblueL} \textbf{Partenaire commercial} & \textbf{Part dans les exportations du Cameroun (en \%, ordre de grandeur récent)} \\
\hline
Chine & 25 \% \\
\hline
Union européenne (dont France, Pays-Bas, Italie) & 30 \% \\
\hline
Inde & 10 \% \\
\hline
Tchad et CEMAC (voisins régionaux) & 8 \% \\
\hline
États-Unis & 5 \% \\
\hline
Autres partenaires & 22 \% \\
\hline
\end{tabularx}
\end{center}

\begin{popfigure}
\begin{center}
\begin{tikzpicture}
\begin{axis}[
    ybar,
    width=0.95\linewidth,
    height=7cm,
    bar width=16pt,
    ymin=0, ymax=35,
    ylabel={Part dans les exportations (en \%)},
    symbolic x coords={UE, Chine, Inde, CEMAC, USA, Autres},
    xtick=data,
    x tick label style={font=\scriptsize},
    ytick={0,5,10,15,20,25,30,35},
    nodes near coords,
    nodes near coords style={font=\scriptsize\bfseries, /pgf/number format/fixed, /pgf/number format/precision=0},
    axis lines=left,
    ymajorgrids=true,
    grid style={popline!50},
    title={\small\bfseries Les principaux partenaires commerciaux du Cameroun (exportations)},
]
\addplot[fill=popblue,draw=popdark] coordinates {(UE,30) (Chine,25) (Inde,10) (CEMAC,8) (USA,5) (Autres,22)};
\end{axis}
\end{tikzpicture}
\end{center}
\legende{Diagramme en barres modèle. Hauteur de chaque barre proportionnelle à la part du partenaire dans les exportations camerounaises (ordres de grandeur récents, en \%). Source : données douanières et INS Cameroun, valeurs arrondies.}
\end{popfigure}

\begin{exemplebox}[La Chine, premier fournisseur du Cameroun]
Exemple chiffré à retenir : la \cle{Chine} est devenue le \textbf{premier fournisseur} du Cameroun, devant la \textbf{France}, l'ancienne puissance coloniale. Elle fournit environ un cinquième des importations camerounaises (machines, textiles, matériel électrique, produits manufacturés), contre une part en recul pour la France. Ce renversement illustre la recomposition des flux mondiaux : le centre de gravité du commerce mondial se déplace vers l'\cle{Asie}, et l'Afrique s'insère désormais dans des échanges \emph{Sud--Sud} qui complètent les anciennes relations Nord--Sud.
\end{exemplebox}

\begin{attention}[Pièges du diagramme en barres]
\begin{itemize}
\item Ne jamais oublier de \textbf{légender les axes} (grandeur et unité : ici, \%).
\item Les barres doivent avoir la \textbf{même largeur} et être espacées régulièrement ; seule leur \emph{hauteur} varie.
\item Classer les barres dans un ordre logique (décroissant ou thématique) facilite la lecture et est apprécié par le correcteur.
\item Vérifier que la somme des parts fait bien 100 \% avant de tracer.
\end{itemize}
\end{attention}

\courssub{TP 3 — Schéma : la chaîne de valeur mondialisée du cacao camerounais}

\textbf{Consigne.} Le Cameroun est le cinquième producteur mondial de cacao (environ 290 000 tonnes par an). Représente par un schéma la \cle{chaîne de valeur} du cacao camerounais, du planteur au consommateur mondial, et fais apparaître la répartition inégale de la valeur ajoutée.

\begin{popfigure}
\begin{center}
\begin{tikzpicture}[node distance=0.55cm and 0.9cm,
  etape/.style={rectangle,rounded corners,draw=popdark,fill=popblueL,minimum width=2.5cm,minimum height=1.0cm,align=center,font=\scriptsize},
  val/.style={font=\scriptsize\itshape,text=popdark},
  fleche/.style={-{Stealth[length=2.5mm]},thick,poporange}]
\node[etape, align=center] (p){\textbf{Planteur}\\ camerounais\\ (Sud, Centre,\\ Littoral)};
\node[etape,right=of p, align=center] (c){\textbf{Collecteur /}\\ \textbf{coopérative}\\ (achat des fèves)};
\node[etape,right=of c, align=center] (e){\textbf{Exportateur}\\ (port de Douala,\\ fèves brutes)};
\node[etape,right=of e, align=center] (t){\textbf{Chocolatiers}\\ européens et\\ asiatiques (transformation)};
\node[etape,right=of t, align=center] (con){\textbf{Consommateur}\\ mondial\\ (tablettes, poudres)};
\draw[fleche] (p) -- (c);
\draw[fleche] (c) -- (e);
\draw[fleche] (e) -- node[above,val]{fèves exportées} (t);
\draw[fleche] (t) -- node[above,val]{produits finis} (con);
\node[below=0.15cm of p,val] {environ 5 \% du prix final};
\node[below=0.15cm of t,val] {l'essentiel de la valeur ajoutée};
\node[font=\small\bfseries,align=center] at ($(e)+(0,2.0)$) {La chaîne de valeur mondialisée du cacao camerounais};
\end{tikzpicture}
\end{center}
\legende{Schéma modèle de la filière cacao. Les flèches orange représentent la circulation du produit ; la valeur ajoutée se concentre à l'aval (transformation et distribution en Europe et en Asie), tandis que le planteur ne perçoit qu'une faible part du prix final (ordre de grandeur : 5 à 7 \%).}
\end{popfigure}

\textbf{Exploitation.} Ce schéma montre un mécanisme essentiel de la mondialisation : les pays du Sud, comme le Cameroun, fournissent des \cle{matières premières brutes} peu rémunératrices, tandis que la \cle{transformation} — là où se concentre la valeur ajoutée — a lieu dans les métropoles du Nord et d'Asie. La construction d'usines de transformation du cacao au Cameroun est une réponse politique à cette inégalité.

\courssub{Analyse croisée de documents et rédaction d'une conclusion argumentée}

Le dossier documentaire fourni comprend une \emph{carte des flux commerciaux mondiaux} et un \emph{tableau statistique} des échanges du Cameroun. L'épreuve du Baccalauréat technique demande souvent de croiser ces deux types de documents pour rédiger une conclusion.

\begin{methode}[Rédiger et justifier une conclusion à partir de documents]
\begin{enumerate}
\item \textbf{Dégager l'idée principale} que les documents illustrent ensemble (ex. : le Cameroun est inséré dans la mondialisation mais de façon inégale et dépendante).
\item \textbf{Citer les documents} explicitement : « le document 1 montre que... », « selon le document 2... ». Une conclusion sans référence aux documents n'est pas justifiée.
\item \textbf{Chiffrer} : appuyer chaque affirmation sur au moins une donnée précise (pourcentage, tonnage, date) tirée des documents.
\item \textbf{Croiser} : montrer ce que la carte révèle (la géographie des flux) et ce que le tableau confirme (les volumes et les partenaires).
\item \textbf{Conclure} par une phrase de synthèse qui répond exactement à la question posée, sans introduire d'élément nouveau.
\end{enumerate}
\end{methode}

\begin{exoresolu}[Conclusion argumentée à partir de deux documents]
Énoncé. À partir de la carte des routes maritimes mondiales (document 1) et du tableau des partenaires commerciaux du Cameroun (document 2), rédigez une conclusion de 8 à 10 lignes répondant à la question : \emph{comment le Cameroun est-il inséré dans les mécanismes de la mondialisation ?}
\tcblower
Solution détaillée. Le Cameroun est pleinement inséré dans les flux de la mondialisation, mais selon une position périphérique et dépendante. Le document 1 montre que le pays est raccordé aux grandes routes maritimes mondiales par la façade atlantique du golfe de Guinée : ses deux ports, Douala (port historique) et Kribi (port en eau profonde inauguré en 2018), relient le pays à l'Europe et à l'Asie, notamment par la route Europe--Asie via Suez. Le document 2 confirme cette insertion par les chiffres : l'Union européenne (environ 30 \%) et la Chine (environ 25 \%) absorbent à elles seules plus de la moitié des exportations camerounaises, tandis que la Chine est devenue le premier fournisseur du pays, devant la France. Cependant, le croquis de la filière cacao révèle les limites de cette insertion : le Cameroun exporte des matières premières brutes (cacao, pétrole, bois, coton) dont la transformation — et donc l'essentiel de la valeur ajoutée — se réalise à l'étranger, le planteur ne percevant qu'environ 5 \% du prix final. En conclusion, le Cameroun participe aux mécanismes de la mondialisation comme fournisseur de matières premières et marché pour les produits manufacturés ; son défi est de transformer sur place et de renforcer son rôle de hub régional, à l'image du port de Kribi.
\end{exoresolu}

\begin{exercice}[À faire en classe ou à la maison]
\begin{enumerate}
\item Refais le croquis du TP 1 sans le modèle, en 20 minutes, puis auto-évalue-toi avec la grille.
\item Construis un diagramme circulaire (ou en barres) des \emph{importations} du Cameroun à partir des données fournies par ton professeur.
\item Rédige une conclusion de 8 lignes croisant le schéma de la filière cacao et le diagramme des partenaires commerciaux, en citant et chiffrant les documents.
\end{enumerate}
\end{exercice}

\begin{corrige}[Exercice : éléments attendus]
\begin{enumerate}
\item Le croquis doit comporter : titre, nord, au moins trois routes maritimes (Atlantique Nord, Europe--Asie par Suez, Pacifique Nord), quatre ports minimum dont Douala et Kribi, les passages de Suez, Malacca et Panama, et une légende complète et ordonnée.
\item Le diagramme doit présenter des barres de même largeur, des axes légendés (partenaires ; part en \%), un titre, et des hauteurs strictement proportionnelles aux données.
\item La conclusion doit citer les deux documents, mobiliser au moins deux chiffres (par exemple 25 \% pour la Chine, 5 \% pour la part du planteur) et se terminer par une phrase de synthèse sur l'insertion inégale du Cameroun dans la mondialisation.
\end{enumerate}
\end{corrige}

\begin{aretenir}[L'essentiel du TP]
Toute production graphique au Baccalauréat technique obéit aux mêmes règles : \cle{titre}, \cle{légende} complète, \cle{orientation} et \cle{proportionnalité} des figurés. Le croquis des routes maritimes montre la hiérarchie des flux mondiaux et la place des ports camerounais (Douala, Kribi) ; le diagramme en barres chiffre la dépendance commerciale du Cameroun envers l'Union européenne et la Chine ; le schéma de la filière cacao révèle la concentration de la valeur ajoutée dans les métropoles du Nord. Croiser carte et tableau, citer et chiffrer les documents, puis conclure : telle est la démarche attendue à l'examen.
\end{aretenir}

\newpage

\coursec{Activité d'intégration}

\coursec{Les mécanismes de la mondialisation : activité d'intégration}

\courssub{Objectifs de l'activité}

À l'issue de cette activité d'intégration, l'élève doit être capable de :

\begin{itemize}
\item mobiliser les savoirs de la leçon (\cle{facteurs} de la mondialisation, \cle{flux}, \cle{acteurs} : FTN, métropoles, États) pour analyser un objet de la vie quotidienne : le téléphone mobile ;
\item lire et croiser des documents géographiques variés (carte, photographie, tableau statistique, données économiques) ;
\item construire un schéma simple montrant l'enchaînement des mécanismes de la mondialisation, du câble sous-marin au mobile money ;
\item rédiger et justifier une conclusion argumentée de 10 à 15 lignes répondant à une question problématisée, en citant précisément les documents utilisés.
\end{itemize}

Cette activité mobilise les deux savoir-faire officiels de l'unité : \emph{appliquer les notions de l'unité à des situations concrètes de la vie courante} et \emph{rédiger et justifier une démarche, une réponse ou une conclusion}.

\courssub{Situation}

\textbf{Le téléphone mobile camerounais : un concentré de mondialisation.}

À Yaoundé, Amina, élève en classe de Terminale technique, utilise chaque jour son smartphone : elle appelle sa tante à Douala, reçoit de l'argent de son frère installé en France par transfert mobile, commande parfois des articles sur des plateformes de commerce en ligne et suit des vidéos éducatives sur Internet. Sans le savoir, chacun de ces gestes quotidiens met en mouvement l'ensemble des mécanismes de la mondialisation étudiés dans la leçon. Pour comprendre comment, son professeur de Géographie constitue un dossier documentaire de quatre documents.

\textbf{Document 1 : les câbles sous-marins raccordant le Cameroun.} Le Cameroun est raccordé à la toile mondiale des câbles sous-marins de fibre optique par plusieurs câbles aboutissant à Douala et à Kribi : SAT-3 (Afrique du Sud -- Europe, mis en service en 2001), WACS (\emph{West Africa Cable System}, 2012), ACE (\emph{Africa Coast to Europe}, 2012) et SAIL (Afrique -- Brésil, 2018). Ces câbles transportent plus de 95\,\% du trafic Internet intercontinental mondial. Le débarquement de ces câbles a fait chuter le coût de la connexion et permis l'essor de l'Internet mobile : le Cameroun comptait environ 10 millions d'internautes au début des années 2020, soit plus d'un tiers de la population.

\textbf{Document 2 : photographie du port de Douala.} La photographie montre le terminal à conteneurs du port autonome de Douala, premier port du Cameroun, par lequel transitent environ 95\,\% des importations du pays. Les smartphones vendus à Yaoundé, fabriqués en Chine, en Corée du Sud ou au Vietnam, arrivent dans des conteneurs standardisés déchargés par des portiques. Le port de Douala traite environ 12 millions de tonnes de marchandises par an.

\textbf{Document 3 : parts de marché de la téléphonie mobile au Cameroun.}

\begin{center}
\begin{tabularx}{\linewidth}{|l|X|X|}\hline
\rowcolor{popblueL} \textbf{Opérateur} & \textbf{Origine du groupe} & \textbf{Part de marché (abonnés)} \\ \hline
MTN Cameroun & Afrique du Sud (groupe MTN) & environ 55\,\% \\ \hline
Orange Cameroun & France (groupe Orange) & environ 35\,\% \\ \hline
Nexttel (Viettel) & Vietnam (groupe Viettel) & environ 10\,\% \\ \hline
\end{tabularx}
\end{center}

Le Cameroun compte plus de 20 millions d'abonnés à la téléphonie mobile. Les trois opérateurs sont des filiales de firmes transnationales (FTN) dont les sièges se trouvent à Johannesburg, Paris et Hanoï.

\textbf{Document 4 : diaspora et mobile money.} La diaspora camerounaise, estimée à plusieurs centaines de milliers de personnes (Europe, Amérique du Nord, Afrique), envoie chaque année au pays plus de 300 milliards de FCFA de transferts de fonds. Une partie croissante de ces transferts transite par le mobile money : au Cameroun, les services MTN Mobile Money et Orange Money comptent plusieurs millions d'utilisateurs et traitent des dizaines de milliards de FCFA de transactions par an. Un transfert Paris--Yaoundé s'effectue désormais en quelques minutes.

\textbf{Question problématique :} le Cameroun est-il un acteur ou un spectateur de la mondialisation ?

\courssub{Consignes}

Le travail se fait par groupes de quatre élèves pour les questions 1 à 4, puis individuellement pour la question 5.

\begin{enumerate}
\item \textbf{Identifier les facteurs.} À partir des documents 1, 2 et 3, relevez les facteurs techniques (transports, télécommunications) et les facteurs économiques et politiques (libéralisation, FTN, ouverture des marchés) qui rendent possible l'usage quotidien du téléphone mobile au Cameroun. Présentez votre réponse dans un tableau à deux colonnes : « facteur » et « document qui le montre ».
\item \textbf{Décrire les flux.} Pour chacune des trois actions suivantes -- un appel téléphonique, un transfert d'argent de la diaspora, une commande en ligne d'un smartphone --, identifiez le ou les types de flux mobilisés (flux d'informations, flux de capitaux, flux de marchandises, flux humains) et précisez le sens du flux (origine $\rightarrow$ destination).
\item \textbf{Repérer les acteurs.} Dans chaque document, citez les acteurs de la mondialisation mis en jeu (FTN, État camerounais, métropoles, diaspora, consommateurs) et précisez leur rôle.
\item \textbf{Construire un schéma.} Réalisez un schéma simple intitulé « Du câble sous-marin au mobile money », faisant apparaître par des flèches légendées : le câble sous-marin, l'opérateur téléphonique, le téléphone de l'utilisateur, la diaspora, le port de Douala.
\item \textbf{Rédiger la conclusion (individuel).} En 10 à 15 lignes, répondez à la question : « Le Cameroun est-il un acteur ou un spectateur de la mondialisation ? » Votre texte doit citer au moins trois documents du dossier et suivre une démarche justifiée (position annoncée, arguments appuyés sur les documents, nuance, conclusion).
\end{enumerate}

\courssub{Ressources à mobiliser}

\begin{aretenir}[Notions utiles de la leçon]
\begin{itemize}
\item \textbf{Facteurs techniques :} révolution des transports (conteneurisation, portiques, hubs portuaires) et des télécommunications (câbles sous-marins de fibre optique, satellites, Internet, téléphonie mobile).
\item \textbf{Facteurs économiques et politiques :} libéralisation des échanges, ouverture des marchés, stratégies des FTN (implantation de filiales, conquête des marchés du Sud).
\item \textbf{Flux :} flux d'informations (données numériques), flux de capitaux (investissements, transferts de fonds), flux de marchandises (commerce conteneurisé), flux humains (migrations, diasporas).
\item \textbf{Acteurs :} les FTN (moteurs de la mondialisation), les États (régulateurs, aménageurs des infrastructures), les métropoles (commandement et concentration des fonctions), les réseaux de diaspora.
\end{itemize}
\end{aretenir}

\begin{methode}[Rédiger une conclusion justifiée en géographie]
\begin{enumerate}
\item Annoncer clairement sa position (une phrase de réponse directe à la question).
\item Développer deux ou trois arguments, chacun appuyé sur un document cité (« le document 3 montre que... »).
\item Introduire une nuance (limites, dépendances, inégalités).
\item Conclure par une phrase de synthèse qui reformule la position.
\end{enumerate}
\end{methode}

\courssub{Démarche et corrigé commenté}

\begin{exoresolu}[Corrigé commenté de l'activité]
\textbf{Question 1 : identifier les facteurs de la mondialisation.}

\emph{Étape 1 : repérer dans chaque document l'indice d'un facteur. Étape 2 : classer le facteur (technique, économique, politique).}

\begin{center}
\begin{tabularx}{\linewidth}{|l|X|}\hline
\rowcolor{popblueL} \textbf{Facteur identifié} & \textbf{Document et justification} \\ \hline
Câbles sous-marins de fibre optique (facteur technique) & Doc. 1 : SAT-3, WACS, ACE et SAIL transportent plus de 95\,\% du trafic Internet intercontinental ; ils rendent possibles l'Internet mobile et les appels. \\ \hline
Conteneurisation et port de commerce (facteur technique) & Doc. 2 : les smartphones importés arrivent en conteneurs au port de Douala (95\,\% des importations du pays). \\ \hline
Présence des FTN (facteur économique) & Doc. 3 : MTN, Orange et Viettel, groupes sud-africain, français et vietnamien, se partagent le marché camerounais. \\ \hline
Ouverture et libéralisation du marché (facteur politique) & Doc. 3 : l'État camerounais a autorisé plusieurs opérateurs étrangers à s'implanter, signe de l'ouverture du secteur des télécommunications. \\ \hline
\end{tabularx}
\end{center}

\emph{Commentaire :} on vérifie ici que les deux catégories de facteurs du programme (techniques d'une part, économiques et politiques d'autre part) sont toutes représentées dans la situation.

\tcblower

\textbf{Question 2 : décrire les flux.}

\begin{itemize}
\item \textbf{Un appel téléphonique :} flux d'informations (voix numérisée) circulant par les réseaux des opérateurs et, pour un appel international, par les câbles sous-marins (Doc. 1) ; sens : Yaoundé $\rightarrow$ Douala, ou Cameroun $\leftrightarrow$ étranger.
\item \textbf{Un transfert d'argent de la diaspora :} flux de capitaux (plus de 300 milliards de FCFA par an, Doc. 4) ; sens : Europe/Amérique du Nord $\rightarrow$ Cameroun ; il s'appuie sur un flux d'informations (la transaction mobile money) et suppose un flux humain préalable (la migration du frère d'Amina).
\item \textbf{Une commande en ligne d'un smartphone :} flux d'informations (commande sur Internet, Doc. 1), flux de capitaux (paiement), puis flux de marchandises (le téléphone fabriqué en Asie arrive en conteneur au port de Douala, Doc. 2) ; sens : Asie $\rightarrow$ Douala $\rightarrow$ Yaoundé.
\end{itemize}

\emph{Commentaire :} un même geste quotidien combine souvent plusieurs types de flux : c'est précisément ce qui illustre le fonctionnement intégré de la mondialisation.

\textbf{Question 3 : repérer les acteurs.}

\begin{itemize}
\item \textbf{FTN :} MTN, Orange, Viettel (Doc. 3) -- elles investissent, déploient les réseaux et captent le marché ; les fabricants asiatiques de smartphones (Doc. 2).
\item \textbf{État camerounais :} régulateur (licences aux opérateurs) et aménageur (raccordement aux câbles, Doc. 1 ; port de Douala, Doc. 2).
\item \textbf{Métropoles :} Douala, interface portuaire et ville de débarquement des câbles ; Johannesburg, Paris, Hanoï, sièges de commandement des FTN.
\item \textbf{Diaspora et consommateurs :} la diaspora alimente les flux de capitaux (Doc. 4) ; les 20 millions d'abonnés camerounais constituent le marché.
\end{itemize}

\textbf{Question 4 : le schéma « Du câble sous-marin au mobile money ».}

Le schéma attendu enchaîne : câble sous-marin (Douala/Kribi) $\rightarrow$ réseau de l'opérateur (FTN) $\rightarrow$ smartphone de l'utilisateur (importé via le port de Douala) ; en parallèle : diaspora (Europe) $\rightarrow$ transfert de fonds $\rightarrow$ mobile money $\rightarrow$ bénéficiaire au Cameroun. Chaque flèche est légendée par le type de flux (informations, capitaux, marchandises, hommes).

\begin{popfigure}
\begin{tikzpicture}[
  node distance=1.8cm and 2.2cm,
  box/.style={draw, rounded corners, fill=popblue!10, text width=2.2cm, align=center, font=\small, minimum height=1.2cm},
  arrow/.style={-Stealth, thick, popdark},
  fluxlabel/.style={font=\footnotesize, text=popdark, midway, fill=white, inner sep=1pt}
]
% Nœuds principaux
\node[box, fill=popblue!20, align=center] (cable){Câble sous-marin\\(Douala, Kribi)};
\node[box, right=of cable, align=center] (operateur){Opérateur téléphonique\\(FTN : MTN, Orange, Viettel)};
\node[box, below right=of operateur, align=center] (smartphone){Smartphone utilisateur\\(importé via port Douala)};
\node[box, fill=poporange!20, above right=of operateur, align=center] (diaspora){Diaspora\\ (Europe, Am. Nord)};
\node[box, fill=popgreen!20, below=of operateur, align=center] (mobilemoney){Mobile Money\\(MTN MoMo, Orange Money)};
\node[box, fill=poppurple!20, right=of mobilemoney, align=center] (benef){Bénéficiaire\\au Cameroun};
\node[box, fill=poppink!20, below=of smartphone, align=center] (port){Port de Douala\\(conteneurs, 12 Mt/an)};

% Flèches flux informations
\draw[arrow, popblue] (cable) -- node[fluxlabel, above, align=center]{Flux d'informations\\ (95\% trafic mondial)} (operateur);
\draw[arrow, popblue] (operateur) -- node[fluxlabel, right, align=center]{Flux d'informations\\ (voix, data)} (smartphone);
\draw[arrow, popblue] (diaspora) -- node[fluxlabel, above, align=center]{Flux d'informations\\ (ordre de virement)} (mobilemoney);

% Flèches flux marchandises
\draw[arrow, poporange] (port) -- node[fluxlabel, left, align=center]{Flux de marchandises\\ (smartphones Asie)} (smartphone);

% Flèches flux capitaux
\draw[arrow, popgreen] (diaspora) -- node[fluxlabel, left, align=center]{Flux de capitaux\\ (>300 Mds FCFA/an)} (mobilemoney);
\draw[arrow, popgreen] (mobilemoney) -- node[fluxlabel, below, align=center]{Flux de capitaux\\ (transfert instantané)} (benef);

% Flèche flux humains (migration)
\draw[arrow, poppurple, dashed] (benef) -- node[fluxlabel, left, align=center]{Flux humains\\ (migration, diaspora)} (diaspora);

% Légende
\node[draw, fill=popcream, rounded corners, font=\footnotesize, anchor=south west] at (current bounding box.south west) {
\textbf{Légende :} \textcolor{popblue}{---} Informations \quad \textcolor{poporange}{---} Marchandises \quad \textcolor{popgreen}{---} Capitaux \quad \textcolor{poppurple}{- - -} Humains
};
\end{tikzpicture}
\legende{Schéma simplifié des mécanismes de la mondialisation : du câble sous-marin au mobile money (flux imbriqués, acteurs multiples).}
\end{popfigure}

\textbf{Question 5 : exemple de conclusion rédigée (10 à 15 lignes).}

\emph{Le Cameroun participe pleinement aux flux de la mondialisation, mais il y occupe une position plus subie que dirigée : il est à la fois acteur et spectateur. Acteur d'abord, car le pays est intégré aux grands réseaux mondiaux : quatre câbles sous-marins le raccordent à l'Internet intercontinental (document 1), le port de Douala traite environ 12 millions de tonnes de marchandises et 95\,\% des importations (document 2), et plus de 20 millions de Camerounais sont abonnés au mobile (document 3). Acteur encore, car sa diaspora injecte plus de 300 milliards de FCFA par an dans l'économie nationale grâce au mobile money (document 4), ce qui fait des Camerounais de véritables producteurs de flux de capitaux. Mais le Cameroun reste aussi spectateur : les câbles, les réseaux et les téléphones appartiennent à des firmes transnationales dont les sièges de commandement sont à Johannesburg, Paris ou Hanoï (document 3), et les smartphones sont fabriqués en Asie (document 2). Les décisions se prennent donc loin de Douala et de Yaoundé. En conclusion, le Cameroun est bien inséré dans la mondialisation, mais dans une position dépendante : il consomme et transite plus qu'il ne commande.}

\emph{Commentaire sur la démarche :} la position est annoncée dès la première phrase ; trois documents sont cités explicitement ; la nuance (« acteur mais dépendant ») est argumentée ; la synthèse finale reformule la réponse. C'est la structure attendue au Baccalauréat technique.
\end{exoresolu}

\courssub{Bilan de l'activité}

Cette activité a permis de mettre en évidence trois idées essentielles de la leçon.

\begin{itemize}
\item \textbf{La mondialisation est concrète et quotidienne.} Un objet banal -- le téléphone mobile -- concentre tous les facteurs étudiés : câbles sous-marins et conteneurisation (facteurs techniques), FTN et ouverture des marchés (facteurs économiques et politiques). La mondialisation n'est pas une abstraction : elle tient dans la poche d'un élève de Terminale.
\item \textbf{Les flux sont imbriqués.} Un appel, un transfert d'argent ou une commande en ligne combinent flux d'informations, de capitaux, de marchandises et d'hommes. Le schéma « du câble sous-marin au mobile money » montre que ces flux empruntent des infrastructures précises (câbles, ports, réseaux) et relient des lieux précis (Douala, Paris, Johannesburg, l'Asie industrielle).
\item \textbf{L'intégration dans la mondialisation est inégale.} Le Cameroun est à la fois acteur (marché de 20 millions d'abonnés, diaspora productive de flux, interface portuaire) et spectateur (commandement détenu par des FTN étrangères, dépendance aux importations). Cette nuance, que l'élève a appris à formuler dans une conclusion rédigée et justifiée, est au cœur du raisonnement géographique attendu au Baccalauréat technique.
\end{itemize}

\begin{attention}[Pièges à éviter dans la rédaction finale]
Ne répondez pas « acteur » ou « spectateur » de façon absolue : le correcteur attend une réponse nuancée, appuyée sur les documents. Citez les documents explicitement (« le document 4 montre que... ») plutôt que de recopier des chiffres sans source. Enfin, respectez la longueur demandée (10 à 15 lignes) : une conclusion trop courte n'argumente pas, une conclusion trop longue sort du cadre de l'exercice.
\end{attention}

\newpage

\coursec{Méthodes et exercices résolus}

\courssub{Méthodes et exercices résolus}

\begin{methode}[Analyser un document (texte, graphique, tableau) sur les facteurs de la mondialisation]
    \textbf{Objectif :} Identifier, classer et expliquer les facteurs (techniques, économiques, politiques) à partir d'une source.
    \begin{enumerate}
        \item \textbf{Identifier la nature et la source du document} (titre, auteur, date, type : texte, courbe, histogramme, carte).
        \item \textbf{Repérer les informations clés} : mots-clés (conteneur, Internet, libéralisation, IDE, OMC), chiffres (dates, pourcentages, volumes), noms d'acteurs (FTN, États, OMC, FMI).
        \item \textbf{Classer les facteurs} en trois colonnes : \cle{Techniques} (transport, TIC), \cle{Économiques} (libre-échange, IDE, division internationale du travail), \cle{Politiques} (régionalisation, politiques d'ouverture, institutions internationales).
        \item \textbf{Établir les liens de causalité} : en quoi tel facteur technique (ex. conteneur) a permis tel phénomène économique (baisse coûts transport $\rightarrow$ hausse échanges) ?
        \item \textbf{Rédiger une synthèse structurée} : introduction (présentation doc + thèse), développement (3 paragraphes = 3 types de facteurs), conclusion (rôle combiné).
    \end{enumerate}
\end{methode}

\begin{exoresolu}[Analyse d'un ensemble documentaire sur l'accélération de la mondialisation depuis 1990]
    \textbf{Document 1 (Texte) :} « La généralisation du conteneur standardisé (normes ISO) dès les années 1970, puis la révolution des technologies de l'information et de la communication (TIC) dans les années 1990 (Internet, fibre optique, GPS) ont réduit de 90\% les coûts de transport maritime et quasi annulé les coûts de transmission de l'information. »
    \textbf{Document 2 (Graphique) :} Évolution du volume des exportations mondiales de marchandises (en milliards \$ courants) : 1990 : 3\,400 ; 2000 : 6\,400 ; 2010 : 15\,200 ; 2022 : 25\,000 (Source : OMC).
    \textbf{Document 3 (Tableau) :} Principaux accords de libéralisation : 1995 (création OMC), 1994 (ALENA), 2001 (entrée Chine OMC), 2021 (ZLECAf entrée en vigueur).
    \textbf{Consigne :} En vous appuyant sur les documents et vos connaissances, montrez comment les facteurs techniques, économiques et politiques se combinent pour expliquer l'accélération de la mondialisation depuis 1990.
    \tcblower
    \textbf{Solution détaillée}
    \begin{enumerate}
        \item \textbf{Présentation de l'ensemble documentaire :} Trois documents complémentaires (texte explicatif, série statistique évolutive, tableau chronologique) couvrant la période 1990-2022, sources officielles (OMC), thème : facteurs de l'accélération de la mondialisation.
        \item \textbf{Facteurs techniques (Doc 1) :} Le texte cite explicitement deux ruptures techniques majeures.
            \begin{itemize}
                \item \textbf{Transport :} Conteneurisation (normalisation ISO) $\rightarrow$ rupture de charge supprimée, manutention mécanisée, économies d'échelle (navires $>$ 20\,000 EVP). \textbf{Résultat :} Baisse de 90\% des coûts de transport maritime (Doc 1), rendant rentable l'échange de biens à faible valeur au kg.
                \item \textbf{TIC :} Internet, fibre optique, GPS $\rightarrow$ gestion en temps réel des chaînes logistiques mondiales, délocalisation des services (offshoring), finance dématérialisée (HFT). \textbf{Résultat :} Quasi-gratuité de l'information, synchronisation mondiale de la production.
            \end{itemize}
        \item \textbf{Facteurs économiques (Doc 1 + Doc 2) :} La baisse des coûts (technique) rend possible l'intensification des échanges (économique).
            \begin{itemize}
                \item Le graphique (Doc 2) montre une multiplication par \num{7,35} du volume des exportations entre 1990 (3\,400 Md\$) et 2022 (25\,000 Md\$).
                \item Taux de croissance moyen annuel (TCAM) : $\left(\frac{25000}{3400}\right)^{\frac{1}{32}}-1 \approx 6,4\%$ par an, bien supérieur à la croissance du PIB mondial ($\sim 3\%$). C'est la \cle{mondialisation en profondeur} (échanges croissent plus vite que la production).
                \item Moteur : Investissements directs à l'étranger (IDE) des FTN, fragmentation de la chaîne de valeur (DITN - Division Internationale du Travail Nouvelle).
            \end{itemize}
        \item \textbf{Facteurs politiques (Doc 3) :} L'encadrement institutionnel lève les barrières résiduelles.
            \begin{itemize}
                \item \textbf{Niveau mondial :} Création OMC (1995) $\rightarrow$ règlement des différends, principe de la nation la plus favorisée (NPF), cycles de négociations (Uruguay Round).
                \item \textbf{Niveau régional :} ALENA (1994), élargissement UE, ZLECAf (2021) $\rightarrow$ zones de libre-échange profondes (normes, services, investissement).
                \item \textbf{Niveau national :} Ouverture unilatérale (Chine entrée OMC 2001, ex-pays socialistes, Inde 1991) $\rightarrow$ intégration de 3 milliards de consommateurs/travailleurs supplémentaires.
            \end{itemize}
        \item \textbf{Synthèse / Conclusion :} La mondialisation post-1990 résulte d'une \textbf{boucle de rétroaction positive} : les innovations techniques (conteneur, TIC) abaissent les coûts de la distance $\rightarrow$ incitent les firmes à fragmenter la production (facteur économique : IDE, DITN) $\rightarrow$ les États accompagnent par la libéralisation (facteur politique : OMC, accords régionaux) $\rightarrow$ l'expansion du marché mondial finance de nouvelles innovations techniques. Les trois facteurs sont indissociables ; aucun n'aurait produit cet effet seul.
    \end{enumerate}
    \textbf{Conclusion :} Depuis 1990, la conjugaison de la révolution logistique et numérique, de la stratégie d'internationalisation des FTN et du cadre libéral multilatéral/régional a provoqué une croissance exponentielle des flux (x7,35 en valeur), marquant le passage à une mondialisation \emph{financiarisée} et \emph{numérique}.
\end{exoresolu}

\begin{methode}[Réaliser un croquis synthétique des flux et acteurs de la mondialisation]
    \textbf{Objectif :} Produire une carte schématique lisible, codée, respectant la sémiologie graphique (points, flèches, surfaces, couleurs).
    \begin{enumerate}
        \item \textbf{Analyser le sujet} : identifier les éléments obligatoires (Triade, pays émergents/PED, PMA, flux commerciaux/financiers/migratoires, sièges FTN, métropoles mondiales, axes maritimes).
        \item \textbf{Choisir le fond de carte} : planisphère projection Robinson ou Mercator simplifié (contours continents). Tracer l'équateur, les tropiques, le méridien de Greenwich.
        \item \textbf{Hiérarchiser l'information (3 niveaux)} :
            \begin{itemize}
                \item \textbf{Niveau 1 (Fond) :} Zones de puissance / intégration (Triade = hachures denses bleu/rouge/vert ; Émergents = hachures pointillées orange ; PMA = hachures légères jaune).
                \item \textbf{Niveau 2 (Flux) :} Flèches \textbf{proportionnelles} (épaisseur $\propto$ volume). Codes couleur : \textcolor{popblue}{Bleu} = marchandises, \textcolor{popgreen}{Vert} = capitaux/IDE, \textcolor{poporange}{Orange} = migrants/informations. Sens : Nord-Nord (épais), Nord-Sud / Sud-Nord (moyens), Sud-Sud (fins mais croissants).
            \end{itemize}
        \item \textbf{Niveau 3 (Acteurs ponctuels) :} Symboles ponctuels.
            \begin{itemize}
                \item $\blacksquare$ Carré noir = Siège social FTN (Top 100) : concentration Triade (NY, Tokyo, Paris, Francfort, Pékin).
                \item $\bullet$ Cercle rouge = Métropoles mondiales (commande : NY, Londres, Tokyo, Paris, Singapour, Shanghai, Dubaï).
                \item $\blacktriangle$ Triangle bleu = Grands ports mondiaux (Shanghai, Singapour, Ningbo, Rotterdam, Los Angeles).
            \end{itemize}
        \item \textbf{Légende obligatoire :} Titre précis, auteur/date, orientation (Nord), échelle qualitative (ex. « Epaisseur flèche $\propto$ volume échanges »), signature.
        \item \textbf{Soin graphique :} Règle, crayons de couleur, écriture fine horizontale, pas de surcharge.
    \end{enumerate}
\end{methode}

\begin{exoresolu}[Croquis : Les grands flux et acteurs de la mondialisation au début des années 2020]
    \textbf{Consigne :} Réaliser un croquis légendé intitulé « La mondialisation en 2023 : flux, acteurs et territoires ». Faire apparaître : la Triade, les pays émergents (BRICS+), les PMA, les 3 grands flux (marchandises, capitaux, informations), les 10 premières métropoles mondiales, les 5 premiers ports à conteneurs, les sièges des 20 premières FTN.
    \tcblower
    \textbf{Solution détaillée (Description du croquis attendu + Code TikZ simplifié pour l'enseignant)}
    \begin{popfigure}
        \begin{tikzpicture}[scale=0.55, every node/.style={font=\tiny}]
            % --- NIVEAU 1 : ZONES (Hachures simulées par remplissage semi-transparent) ---
            % Triade : USA, UE, Japon/Asie Est développée
            \fill[popblue!15] (-18,2) rectangle (-8,10); % AmNord + UE approx
            \fill[popblue!15] (12,2) rectangle (18,10);  % Japon/Corée/Chine côtière
            \node[align=center, popblue, font=\tiny\bfseries] at (-13,6) {TRIADE};
            
            % Émergents (BRICS+)
            \fill[poporange!15] (-18,-10) rectangle (-8,-2); % Brésil/AmSud
            \fill[poporange!15] (-5,-5) rectangle (5,0);    % Afrique subsaharienne / Moyen-Orient partiel
            \fill[poporange!15] (5,-5) rectangle (18,2);    % Chine/Inde/Asie SE
            \node[align=center, poporange, font=\tiny\bfseries] at (10,-2) {ÉMERGENTS\\ (BRICS+)};
            
            % PMA (reste)
            \fill[popgold!15] (-20,-10) rectangle (20,-5);  % Zone sud large
            
            % --- NIVEAU 2 : FLUX (Flèches proportionnelles) ---
            % Flux Merchandises (Bleu) - Nord-Nord fort, Nord-Sud, Sud-Sud
            \draw[popblue, very thick, -{Stealth[length=3mm]}] (-13,6) -- (-3,6); % Transatlantique
            \draw[popblue, very thick, -{Stealth[length=3mm]}] (15,6) -- (5,6);   % Transpacifique
            \draw[popblue, thick, -{Stealth[length=2mm]}] (-13,6) -- (-13,-2);   % Nord-Sud Amériques
            \draw[popblue, thick, -{Stealth[length=2mm]}] (-3,6) -- (-3,-2);     % Europe-Afrique
            \draw[popblue, thick, -{Stealth[length=2mm]}] (15,6) -- (10,-2);     % Asie-Est -> Asie SE
            \draw[popblue, thin, dashed, -{Stealth[length=1.5mm]}] (-13,-2) -- (10,-2); % Sud-Sud (croissant)
            
            % Flux Capitaux/IDE (Vert) - Sortie Triade vers Émergents + Retours
            \draw[popgreen!70!black, thick, -{Stealth[length=2mm]}, bend left=20] (-13,8) to (-5,-1);
            \draw[popgreen!70!black, thick, -{Stealth[length=2mm]}, bend right=20] (15,8) to (10,-1);
            \draw[popgreen!70!black, thin, -{Stealth[length=1.5mm]}] (-5,-1) -- (10,-1); % Sud-Sud IDE
            
            % Flux Informations/Migrants (Orange) - Concentrés métropoles
            \draw[poporange, thin, -{Stealth[length=1.5mm]}, bend left=10] (-13,9) to (15,9); % Fibre optique Nord-Nord
            \draw[poporange, thin, -{Stealth[length=1.5mm]}] (-13,4) -- (-13,-4); % Migration Sud-Nord Am
            \draw[poporange, thin, -{Stealth[length=1.5mm]}] (-3,4) -- (-3,-4);   % Migration Afr->Eur
            
            % --- NIVEAU 3 : ACTEURS PONCTUELS ---
            % Métropoles mondiales (Cercles rouges/roses)
            \foreach \x/\y/\name in {
                -15/7/New York, -3/8/Londres, -1/7/Paris, 2/7/Francfort,
                15/6/Tokyo, 13/4/Shanghai, 11/2/Singapour, 5/3/Dubaï,
                -15/2/São Paulo, -8/-2/Johannesburg
            } {
                \fill[poppink] (\x,\y) circle (2.5pt);
                \node[above=1pt, poppink, font=\tiny\bfseries] at (\x,\y) {\name};
            }
            
            % Ports (Triangles bleus)
            \foreach \x/\y/\name in {
                13/4/Shanghai, 11/2/Singapour, 12/5/Ningbo, -3/9/Rotterdam, -15/6/LosAngeles
            } {
                \fill[popblue] (\x,\y) -- ++(-2pt,-3pt) -- ++(4pt,0) -- cycle;
                \node[below=1pt, popblue, font=\tiny] at (\x,\y-3pt) {\name};
            }
            
            % Sièges FTN Top 20 (Carrés noirs) - Concentrés Triade + Pékin
            \foreach \x/\y in {
                -15/7, -14/6.5, -3/8, -2/7.5, -1/7, 2/7, 15/6, 14/5.5, 13/4, 11/2, 5/3
            } {
                \fill[popdark] (\x,\y) rectangle ++(3pt,3pt);
            }
            
            % Équateur / Tropiques / Méridien
            \draw[popmuted, dashed] (-20,0) -- (20,0) node[right] {Équateur};
            \draw[popmuted, dashed] (-20,4.6) -- (20,4.6) node[right] {Tropique Cancer};
            \draw[popmuted, dashed] (-20,-4.6) -- (20,-4.6) node[right] {Tropique Capricorne};
            \draw[popmuted, dashed] (0,-10) -- (0,10) node[above] {Méridien Greenwich};
            
            % Flèche Nord
            \draw[popdark, thick, -{Stealth[length=4mm]}] (18,-8) -- (18,-6) node[above] {N};
        \end{tikzpicture}
        \legende{Croquis schématique : La mondialisation en 2023. Flèches bleu = marchandises (épaisseur $\propto$ volume), vert = capitaux/IDE, orange = informations/migrations. Carrés noirs = Sièges FTN (échantillon Top 20). Cercles roses = Métropoles mondiales. Triangles bleus = 5 premiers ports. Zones colorées : Triade (bleu clair), Émergents (orange clair), PMA (jaune clair).}
    \end{popfigure}
    
    \textbf{Points de vigilance pour la notation (Bac) :}
    \begin{itemize}
        \item \textbf{Légende complète et hiérarchisée} (3 niveaux) = 4 pts.
        \item \textbf{Proportionnalité des flèches} (Nord-Nord > Nord-Sud > Sud-Sud) = 3 pts.
        \item \textbf{Localisation précise} des 10 métropoles et 5 ports = 3 pts.
        \item \textbf{Propreté / Esthétique} (règle, couleurs, écriture) = 2 pts.
        \item \textbf{Titre informatif} (Thème + Date + Échelle qualitative) = 1 pt. Total = 13/13 pts croquis.
    \end{itemize}
\end{exoresolu}

\begin{methode}[Calculer et interpréter un taux de variation (échanges, IDE, flux migratoires)]
    \textbf{Objectif :} Maîtriser le calcul de taux de variation (annuel ou sur période) et l'interpréter géographiquement (ralentissement, accélération, rupture).
    \begin{enumerate}
        \item \textbf{Identifier les bornes temporelles} : $V_0$ (valeur initiale, année de départ), $V_1$ (valeur finale, année d'arrivée), $n$ (nombre d'années).
        \item \textbf{Choisir la formule adaptée} :
            \begin{itemize}
                \item \textbf{Taux de variation global (sur la période) :} $t_{global} = \frac{V_1 - V_0}{V_0} \times 100$ (en \%).
                \item \textbf{Taux moyen annuel (TMA / TCAM - Taux de Croissance Annuel Moyen) :} $1 + t_{annuel} = \left(\frac{V_1}{V_0}\right)^{\frac{1}{n}} \Rightarrow t_{annuel} = \left(\frac{V_1}{V_0}\right)^{\frac{1}{n}} - 1$. Exprimé en \% ($ \times 100$).
            \end{itemize}
        \item \textbf{Calculer avec rigueur} : attention aux unités (milliards, millions), aux virgules/points, aux parenthèses sur la calculatrice.
        \item \textbf{Interpréter géographiquement} : comparer au taux de croissance du PIB mondial ($\sim 3\%$), au taux de croissance du commerce mondial (élasticité), aux sous-périodes (crise 2008, Covid 2020, reprise). Dire : « Le commerce croît 2 fois plus vite que la production $\rightarrow$ approfondissement de la DITN ».
        \item \textbf{Conclure} : lien avec le thème (mondialisation en profondeur, ralentissement récent « slowbalisation »).
    \end{enumerate}
\end{methode}

\begin{exoresolu}[Calcul et interprétation de l'évolution des IDE mondiaux (1990-2022)]
    \textbf{Données (Source CNUCED) :} Flux mondiaux d'Investissements Directs à l'Étranger (IDE) entrants.
    \begin{center}
        \begin{tabular}{|c|c|c|}\hline
        \textbf{Année} & \textbf{Flux IDE (Milliards \$ courants)} & \textbf{Part des pays en développement (\% )} \\ \hline
        1990 & 200 & 18 \\ \hline
        2000 & 1\,400 & 25 \\ \hline
        2010 & 1\,500 & 45 \\ \hline
        2022 & 1\,300 & 65 \\ \hline
        \end{tabular}
    \end{center}
    \textbf{Consignes :}
    \begin{enumerate}
        \item Calculer le taux de variation global des flux IDE mondiaux entre 1990 et 2022.
        \item Calculer le taux moyen annuel de croissance (TCAM) sur la même période.
        \item Calculer l'évolution de la part des pays en développement dans les flux entrants (en points de pourcentage et en \%\ de variation relative).
        \item Interpréter ces résultats au regard des mutations de la mondialisation (acteurs, géographie).
    \end{enumerate}
    \tcblower
    \textbf{Solution détaillée}
    \begin{enumerate}
        \item \textbf{Taux de variation global (1990-2022) :}
            $V_0 = 200$, $V_1 = 1300$, $n = 32$ ans.
            $t_{global} = \frac{1300 - 200}{200} \times 100 = \frac{1100}{200} \times 100 = \textbf{550\%}$.
            \emph{Les flux mondiaux d'IDE ont été multipliés par 6,5 en 32 ans.}
        
        \item \textbf{Taux moyen annuel (TCAM) :}
            $1 + t = \left(\frac{1300}{200}\right)^{\frac{1}{32}} = (6,5)^{0,03125}$.
            $\ln(6,5) \approx 1,8718$. $\frac{1,8718}{32} \approx 0,0585$. $e^{0,0585} \approx 1,0602$.
            $t_{annuel} \approx 0,0602 = \textbf{6,02\% par an}$.
            \emph{Note : Ce taux masque de fortes variations (bulle 2000, crise 2008, pic 2015-2016, chute Covid).}
        
        \item \textbf{Évolution de la part des PED :}
            \begin{itemize}
                \item \textbf{En points de pourcentage :} $65 - 18 = \textbf{+47 points}$.
                \item \textbf{En variation relative :} $\frac{65 - 18}{18} \times 100 = \frac{47}{18} \times 100 \approx \textbf{+261\%}$.
            \end{itemize}
        
        \item \textbf{Interprétation géographique :}
            \begin{itemize}
                \item \textbf{Explosion des flux (x6,5) :} Témoigne de la stratégie des FTN (DITN, contournement barrières tarifaires, proximité marchés). Croissance IDE ($+6\%$/an) > Croissance PIB mondial ($\sim 3\%$) $\rightarrow$ \cle{mondialisation par le capital} (production délocalisée plutôt qu'exportée).
                \item \textbf{Bascule vers les PED (Part : 18\% $\rightarrow$ 65\%) :} Preuve majeure de l'émergence. Les PED ne sont plus seulement \emph{récepteurs} passifs (usines bas coûts) mais deviennent \emph{émetteurs} d'IDE (FTN chinoises, indiennes, brésiliennes, sud-africaines : Huawei, Tata, Petrobras, MTN). Entrée Chine OMC (2001) = point de bascule.
                \item \textbf{Baisse absolue 2022 (1300 vs pic 2016 $\sim$ 2000) :} Signe de \cle{slowbalisation} / « démondialisation sélective » : raccourcissement chaînes de valeur (near-shoring, friend-shoring), tensions géopolitiques (US-Chine), résilience post-Covid.
            \end{itemize}
    \end{enumerate}
    \textbf{Conclusion :} La croissance soutenue des IDE (+6\%/an) et le renversement de leur géographie (PED majoritaires) illustrent le passage d'une mondialisation \emph{inter-industrielle Nord-Nord} à une mondialisation \emph{intra-industrielle Nord-Sud-Sud}, pilotée par les FTN, mais aujourd'hui freinée par la géopolitique.
\end{exoresolu}

\begin{methode}[Caractériser une Firme Transnationale (FTN) à partir de données chiffrées et de sa stratégie spatiale]
    \textbf{Objectif :} Démontrer qu'une firme est transnationale (critères quantitatifs) et analyser sa stratégie d'internationalisation (critères qualitatifs).
    \begin{enumerate}
        \item \textbf{Vérifier les critères quantitatifs (définition OCDE/ONUDI) :}
            \begin{itemize}
                \item Siège social dans un pays (pays d'origine).
                \item Filiales/affiliées dans \textbf{au moins 2 autres pays} (pays d'accueil).
                \item \textbf{Seuil de contrôle :} Participation $\ge 10\%$ du capital votant dans la filiale (norme FMI/OCDE).
                \item \textbf{Part de l'international :} Au moins 25\% (souvent 50\%+) du Chiffre d'Affaires (CA), des Effectifs, ou des Actifs hors du pays d'origine. \cle{Indice de transnationalité} (moyenne des 3 ratios).
            \end{itemize}
        \item \textbf{Analyser la stratégie spatiale (Modèle de Bartlett \& Ghoshal / Dunning OLI) :}
            \begin{itemize}
                \item \textbf{Recherche d'efficience (Efficiency-seeking) :} Délocalisation production vers bas coûts (main-d'œuvre, énergie, fiscalité) $\rightarrow$ \cle{Plateformes d'exportation} (ex. Mexique, Vietnam, Maroc, Europe de l'Est).
                \item \textbf{Recherche de marchés (Market-seeking) :} Implantation près des clients (contournement barrières, adaptation produit) $\rightarrow$ \cle{Production locale pour marché local} (ex. Chine, Inde, Brésil, USA pour firmes européennes).
                \item \textbf{Recherche de ressources/actifs stratégiques (Resource/Asset-seeking) :} Accès matières premières (pétrole, minerais, terres rares) ou technologies/marques (R\&D, rachats startups) $\rightarrow$ \cle{Amont (extractif) ou Aval (high-tech)}.
            \end{itemize}
        \item \textbf{Identifier le modèle d'organisation :}
            \begin{itemize}
                \item \textbf{Global :} Produit standardisé, décision centralisée (siège), économies d'échelle max (ex. Intel, Pfizer, Coca-Cola).
                \item \textbf{Multidomestique :} Produits adaptés, filiales autonomes, réactivité locale (ex. Nestlé, Unilever, banques de détail).
                \item \textbf{Transnational (idéal-type) :} Interdépendance filiales, flux bidirectionnels savoirs/produits, innovation distribuée (ex. Toyota, GE, Apple \emph{design US / prod Chine / R\&D Israël/Inde}).
            \end{itemize}
        \item \textbf{Évaluer l'impact territorial :} Emplois directs/indirects, transferts technologiques, effets d'entraînement (sous-traitants), risques (délocalisation brutale, évasion fiscale, pression environnementale).
    \end{enumerate}
\end{methode}

\begin{exoresolu}[Étude de cas : Apple Inc., archétype de la FTN « Transnationale » à l'ère numérique]
    \textbf{Données 2023 (Rapport annuel 10-K) :}
    \begin{itemize}
        \item CA Total : \SI{383}{milliards \$} ; CA Hors Amériques (Intl) : \SI{247}{milliards \$} (\textbf{64,5\%}).
        \item Effectifs : \SI{164 000}{employés} ; \num{60\%} hors USA (majorité Chine/Inde via sous-traitants Foxconn/Pegatron non consolidés, mais filiales retail/R\&D oui).
        \item Actifs : Majorité propriété intellectuelle (IP) détenue par filiales irlandaises (historique) / US maintenant ; Usines \textbf{zéro} en propriété directe (modèle \cle{fabless} / \cle{asset-light}).
        \item Présence : Vente dans \textbf{175+ pays} ; R\&D : USA (Cupertino), Chine, Inde, Israël, UK, Allemagne ; Assemblage : Chine (Zhengzhou « iPhone City »), Inde (Tamil Nadu), Brésil ; Composants : Corée (écrans, puces), Taïwan (TSMC puces), Japon, USA, Europe.
    \end{itemize}
    \textbf{Consigne :} Démontrer qu'Apple est une FTN. Caractériser sa stratégie spatiale et son modèle d'organisation. Quels sont les enjeux géopolitiques actuels pour cette firme ?
    \tcblower
    \textbf{Solution détaillée}
    \begin{enumerate}
        \item \textbf{Preuve du statut FTN (Critères quantitatifs) :}
            \begin{itemize}
                \item Siège : Cupertino, Californie (USA).
                \item Filiales contrôlées ($>10\%$) : \textbf{Des centaines} dans plus de \textbf{50 pays} (Apple Sales International Irlande, Apple China, Apple India, Apple Japan, Apple Operations Europe, FileMaker, Beats, Shazam, etc.).
                \item \textbf{Indice de transnationalité (estimation) :}
                    \begin{align*}
                        \text{Part CA Intl} &= 64,5\% \\
                        \text{Part Actifs Intl (IP + Cash offshore)} &\approx 80\%+ \\
                        \text{Part Effectifs Intl (directs)} &\approx 40\% \text{ (mais } >95\% \text{ si sous-traitance incluse)}
                    \end{align*}
                    Indice moyen $> 60\%$ $\gg$ Seuil 25\%. \textbf{Apple est une FTN majeure (Top 5 capitalisation mondiale).}
            \end{itemize}
        
        \item \textbf{Stratégie spatiale combinée (Modèle OLI - Dunning) :}
            \begin{itemize}
                \item \textbf{Asset-seeking (R\&D / Propriété Intellectuelle) :} Cœur de la valeur ajoutée (\textbf{60\%+ du prix iPhone}). Concentré au siège (Cupertino) + pôles d'excellence décentralisés (IA en UK/Israël/Inde, Puces en Allemagne/Autriche, Écrans au Japon/Corée). \textbf{Avantage de propriété (O) :} Marque, iOS, Silicon (puces M/A series), écosystème verrouillé.
                \item \textbf{Efficiency-seeking (Assemblage / Fabrication) :} \textbf{Externalisation totale (Fabless)}. Partenaires : Foxconn (Hon Hai), Pegatron, Luxshare, Tata (Inde).
                    \begin{itemize}
                        \item \textbf{Chine (Zhengzhou, Shenzhen) :} Écosystème complet fournisseurs, main-d'œuvre qualifiée massive, infrastructure logistique, rapidité (scale). \emph{Risque : concentration excessive.}
                        \item \textbf{Inde (Tamil Nadu, Karnataka) :} Subventions PLI (Production Linked Incentive), main-d'œuvre jeune/peu coûteuse, marché intérieur 1,4 Md hab. \emph{Stratégie « China+1 ».}
                        \item \textbf{Vietnam / Brésil :} Assemblage AirPods, iPad, Mac ; contournement tarifs US/Chine ; marché local.
                    \end{itemize}
                \item \textbf{Market-seeking (Vente / Retail / Services) :} \textbf{Apple Stores (526 magasins, 26 pays)} + App Store / iCloud / Apple Music / Apple Pay. Revenus Services (\SI{85}{Md\$} en 2023) = marge brute $\sim 70\%$, fidélisation, revenu récurrent. Implantation prioritaire métropoles mondiales (haut revenu).
            \end{itemize}
        
        \item \textbf{Modèle d'organisation : \textbf{Transnational (Hybride Global/Transnational)}.}
            \begin{itemize}
                \item \textbf{Global :} Produit unique mondial (iPhone 15 = même partout), OS unique (iOS), Supply Chain pilotée centralement (Tim Cook = ex-Supply Chain), Marque unique.
                \item \textbf{Transnational :} Innovation distribuée (puces conçues US/Israël/Allemagne, fabriquées Taïwan TSMC 3nm ; écrans Corée/Japon ; capteurs Sony/US). Flux de savoirs bidirectionnels : Retours terrain Chine/Inde $\rightarrow$ adaptation produits (Double SIM physique Chine, prise en charge langues locales, UPI paiement Inde).
                \item \textbf{Pas Multidomestique :} Pas d'autonomie stratégique filiales pays ; standardisation extrême.
            \end{itemize}
        
        \item \textbf{Enjeux géopolitiques actuels (Risques systémiques) :}
            \begin{enumerate}
                \item \textbf{Dépendance Taïwan (TSMC) :} 90\% puces avancées mondiales. Risque blocus/ invasion Chine $\rightarrow$ Rupture approvisionnement critique. \textbf{Réponse :} Investissements TSMC Arizona (USA), Japon, Allemagne (subventions CHIPS Act) ; diversification vers Samsung Foundry / Intel Foundry.
                \item \textbf{Tensions US-Chine :} Sanctions US (Entity List), restrictions export puces IA (H100, A100), loi anti-espionnage Chine. Apple coincé : Chine = 19\% CA + 95\% assemblage. \textbf{Réponse :} Accélération « India+Vietnam » (objectif 25-30\% prod Inde 2027), lobbying intense Washington.
                \item \textbf{Fiscalité / Régulation UE/USA :} DMA (Digital Markets Act) $\rightarrow$ ouverture App Store, paiement tiers, interopérabilité. Menace sur rente Services (30\% commission). Irlande / OCDE (Pilier 2, impôt min 15\%) $\rightarrow$ fin optimisation fiscale « Double Irish / Dutch Sandwich ».
                \item \textbf{Chaîne de valeur « Scope 3 » (Émissions indirectes) :} 99\% empreinte carbone = fournisseurs (Foxconn, TSMC, mines). Pression Net Zero 2030 (Apple) / 2050 (fournisseurs). Exigence énergie renouvelable fournisseurs.
            \end{enumerate}
    \end{enumerate}
    \textbf{Conclusion :} Apple incarne la FTN \emph{asset-light, fabless, écosystémique} du XXIe siècle : elle capte la rente de l'innovation (propriété intellectuelle, design, logiciel, marque) en externalisant le capital fixe et la main-d'œuvre bas coût, mais se heurte aujourd'hui à la \textbf{réalité géographique dure} (géopolitique Taïwan/Chine, régulation étatique, urgence climatique) qui l'oblige à « ré-territorialiser » partiellement sa chaîne de valeur (friend-shoring).
\end{exoresolu}

\begin{methode}[Rédiger un paragraphe argumenté structuré (Thèse - Arguments - Exemples - Conclusion partielle)]
    \textbf{Objectif :} Répondre à une question de dissertation ou d'analyse de documents par un développement rigoureux, standard attendu au Bac (Probatoire/Bac Tech).
    \begin{enumerate}
        \item \textbf{Analyser le sujet (Mots-clés, Verbes d'action, Bornes spatio-temporelles) :} Souligner : « Montrez que », « Expliquez pourquoi », « Dans quelle mesure », « Les acteurs », « Les territoires ».
        \item \textbf{Formuler la THÈSE (phrase unique, réponse directe au sujet) :} Ex: « Les métropoles mondiales sont les nœuds de commande de la mondialisation car elles concentrent les fonctions de décision, d'innovation et de finance. »
        \item \textbf{Construire 2 à 3 ARGUMENTS distincts (A1, A2, A3) :} Chaque argument = une idée force. Structure : \textbf{Argument} $\rightarrow$ \textbf{Explication (mécanisme)} $\rightarrow$ \textbf{Exemple chiffré / Cas concret (Preuve)}.
        \item \textbf{Utiliser le vocabulaire précis du cours :} \cle{Métropoles mondiales}, \cle{Fonctions de commande}, \cle{FTN}, \cle{IDE}, \cle{Flux financiers}, \cle{Aéroports hubs}, \cle{Câbles sous-marins}, \cle{Places boursières}, \cle{Sièges sociaux}, \cle{Réseaux}, \cle{Inégalités métropolisation}.
        \item \textbf{Rédiger la CONCLUSION PARTIELLE :} Résumer l'argumentaire, ouvrir sur une nuance ou le paragraphe suivant (ex: « Cependant, cette puissance est inégalement répartie et contestée... »).
        \item \textbf{Relire :} Vérifier l'absence de « je », le présent de vérité générale, la cohérence logique, l'orthographe des noms propres (NY, Londres, Tokyo, Shanghai, Singapour, Dubaï, Paris, Francfort, São Paulo, Johannesburg).
    \end{enumerate}
\end{methode}

\begin{exoresolu}[Rédaction argumentée : « Les États ont-ils perdu le contrôle de la mondialisation ? »]
    \textbf{Consigne :} Rédiger un paragraphe argumenté d'environ 20 lignes répondant à la question. Utiliser la méthode Thèse/Arguments/Exemples.
    \tcblower
    \textbf{Proposition de rédaction (Niveau Bac Excellent)}
    
    \textbf{Thèse :} Si la mondialisation a considérablement réduit la marge de manœuvre des États en les insérant dans des interdépendances contraignantes, ils n'ont pas « perdu le contrôle » : ils en redéfinissent les règles, en régulent les excès et en orientent les flux par des stratégies de puissance, devenant à la fois \textbf{soumis} et \textbf{acteurs stratégiques} du système.
    
    \textbf{Argument 1 : La contrainte structurelle — L'érosion de la souveraineté économique par les flux et les FTN.}
    L'ouverture des capitaux (libéralisation financière années 1980-90) et la mobilité des firmes transnationales (FTN) soumettent les politiques nationales à la \cle{concurrence fiscale et sociale} à la baisse (« course au fond »). Une hausse d'impôt sur les sociétés ou un salaire minimum trop élevé risque la fuite des IDE ou la délocalisation. \textbf{Exemple :} L'Irlande a attiré les sièges européens de Google, Apple, Pfizer par un IS à 12,5\% (avant accord OCDE 2021), privant la France ou l'Allemagne de recettes fiscales. Les agences de notation (Moody's, S\&P, Fitch) sanctionnent les écarts budgétaires, imposant l'austérité (Grèce 2010-2018).
    
    \textbf{Argument 2 : La réponse régulatrice — Les États reconstruisent la gouvernance multilatérale et régionale.}
    Loin de disparaître, l'État renforce le droit international pour encadrer la finance et le commerce. \textbf{Exemple :} L'\textbf{OCDE/G20} impose depuis 2023 un \textbf{impôt minimum mondial de 15\%} (Pilier 2) sur les bénéfices des FTN (CA $>$ 750 M€), mettant fin à l'optimisation agressive. L'\textbf{OMC} (malgré son blocage actuel) et les \textbf{accords de nouvelle génération} (CPTPP, UE-Mercosur, ZLECAf) intègrent normes sociales, environnementales, numériques. Les \textbf{banques centrales} (Fed, BCE, PBOC) coordonnent leurs politiques monétaires (swap lines) pour stabiliser le dollar, monnaie de réserve mondiale.
    
    \textbf{Argument 3 : L'offensive géopolitique — L'État redevient stratège : « friend-shoring », subventions, souveraineté.}
    La rivalité US-Chine (guerre technologique, puces, IA) et la pandémie ont relégué l'efficacité pure (coût minimum) au second plan derrière la \textbf{résilience} et la \textbf{sécurité}. \textbf{Exemple :} Le \textbf{CHIPS and Science Act} US (2022, \SI{52}{Md\$}) et le \textbf{European Chips Act} (\SI{43}{Md\$}) subventionnent massivement la relocalisation de la fabrication de semi-conducteurs (TSMC Arizona, Intel Magdeburg). La Chine impose le stockage local des données (Cybersecurity Law) et soutient ses champions (Huawei, SMIC) par l'État. La \cle{diplomatie économique} (Nouvelles Routes de la Soie, Partenariat Global UE) utilise l'aide, l'investissement, les normes comme armes d'influence.
    
    \textbf{Conclusion partielle :} L'État n'a pas disparu ; il a \textbf{muté} : du « État keynésien » gestionnaire de l'économie nationale au « \textbf{État stratège »} naviguant dans la mondialisation. Il subit les flux (contrainte) mais façonne les règles (régulation) et arme l'économie (puissance). La mondialisation n'est pas un destin subi, mais un \textbf{champ de bataille} où les États, inégaux (USA/Chine/UE vs PMA), négocient leur place.
    
    \textbf{Auto-évaluation (Grille Bac) :}
    \begin{itemize}
        \item Thèse claire, nuancée (ni tout puissant, ni impuissant) : 2/2.
        \item 3 arguments distincts (Contrainte, Régulation, Géopolitique) : 3/3.
        \item Exemples précis, datés, chiffrés (Irlande 12,5\%, OCDE 15\%, CHIPS Act 52 Md\$, ZLECAf) : 4/4.
        \item Vocabulaire expert (friend-shoring, course au fond, Pilier 2, souveraineté, interdépendance) : 3/3.
        \item Structure logique, connecteurs logiques (Si... cependant, Loin de..., Exemple:) : 3/3. \textbf{Total : 15/15.}
    \end{itemize}
\end{exoresolu}

\begin{methode}[Comparer les effets de la mondialisation sur des territoires inégaux (Nord / Émergents / PMA)]
    \textbf{Objectif :} Structurer une comparaison à double entrée (Acteurs / Territoires) ou (Dimensions : Éco / Social / Enviro) pour éviter la juxtaposition.
    \begin{enumerate}
        \item \textbf{Construire un tableau de comparaison (brouillon) :} Lignes = Dimensions (Croissance PIB, Emploi/Industrialisation, Inégalités, Environnement, Insertion flux, Dépendances). Colonnes = Groupes de pays (Triade / Pays Émergents BRICS+ / PMA / Pays en transition).
        \item \textbf{Identifier les contrastes majeurs (Ruptures) :}
            \begin{itemize}
                \item \textbf{Croissance :} Émergents ($+4$ à $+6\%$) > Monde ($+3\%$) > Triade ($+1$ à $+2\%$) > PMA (instable, souvent $<3\%$).
                \item \textbf{Industrialisation :} Émergents = \textbf{usine du monde} (part manuf. export $\uparrow$) ; Triade = \textbf{désindustrialisation relative} (part manuf. emploi $\downarrow$, tertiarisation) ; PMA = \textbf{primarisation} (export matières premières, import manuf.).
                \item \textbf{Inégalités :} \textbf{Entre pays :} Écart PIB/hab Triade/PMA se creuse (facteur 50 à 100). \textbf{Dedans pays :} Hausse inégalités quasi générale (Kuznets inversé ?) : 1\% top capte 38\% croissance mondiale 1995-2021 (Rapport Chancel/Wild).
                \item \textbf{Environnement :} Émergents = \textbf{atelier pollué} (export carbone incorporé) ; Triade = \textbf{consommateurs nets} (empreinte carbone > émissions territoriales) ; PMA = \textbf{victimes} (changement climatique, extraction prédatrice).
            \end{itemize}
        \item \textbf{Rédiger par dimensions (pas par pays) :} Paragraphe 1 : Croissance/Convergence ? Paragraphe 2 : Structures productives/Emploi. Paragraphe 3 : Inégalités sociales. Paragraphe 4 : Soutenabilité écologique.
        \item \textbf{Nuancer :} Hétérogénéité \emph{inside} groupes (Chine $\neq$ Inde $\neq$ Brésil ; Botswana $\neq$ Tchad ; Allemagne $\neq$ Grèce). Rôle des politiques publiques (Corée du Sud vs Argentine).
    \end{enumerate}
\end{methode}

\begin{exoresolu}[Tableau comparatif et analyse : Mondialisation et développement inégal]
    \textbf{Consigne :} Compléter le tableau ci-dessous avec des données indicatives (années 2020-2023) et rédiger une synthèse de 15 lignes sur les contrastes de développement.
    
    \begin{center}
        \begin{tabularx}{\linewidth}{|l|X|X|X|X|}\hline
        \rowcolor{popblueL}
        \textbf{Indicateur} & \textbf{Triade (USA, UE, Japon)} & \textbf{Pays Émergents (Chine, Inde, Brésil...)} & \textbf{PMA (Afrique subsaharienne, Haïti...)} & \textbf{Analyse / Tendance} \\ \hline
        Croissance PIB réel (\%/an) & 1,5 -- 2,5 & 4 -- 6 (Chine $\sim$5, Inde $\sim$7) & 3 -- 4 (volatil, pop $\uparrow$) & \textbf{Divergence} : Convergence \textbf{relative} Émergents vs Triade. PMA \textbf{décrochent} (croissance $<$ démographie). \\ \hline
        Part industrie dans PIB (\%/Emploi) & PIB $\sim$15-20\% / Emploi $\sim$15\% & PIB $\sim$30-40\% (Chine 38\%) / Emploi $\uparrow$ & PIB $\sim$10-15\% / Emploi $\sim$10\% & \textbf{Spécialisation} : Émergents = \textbf{Usine monde}. Triade = \textbf{Tertiarisation supérieure} (R\&D, Finance, Santé). PMA = \textbf{Primarisation}. \\ \hline
        Part export manuf. / total export & $> 70\%$ (haute tech) & $> 80\%$ (moyenne tech, assemblage) & $< 20\%$ (matières 1ères, agro) & \textbf{Division Int. Travail (DIT) :} Hiérarchique. Émergents montent en gamme (Chine puces, EV). PMA piégés \textbf{malédiction ressources}. \\ \hline
        IDE entrants (Mds \$ / Part monde) & $\sim$ 600 / 45\% (mais \textbf{sorties} massives) & $\sim$ 850 / 65\% (Chine $\sim$ 180, Inde $\sim$ 50, Brésil $\sim$ 40) & $\sim$ 25 / 2\% & \textbf{Bascule historique :} PED 1er récepteurs mondiaux. Triade = \textbf{Investisseurs nets} (sorties $>$ entrées). \\ \hline
        Coeff. Gini (Inégalités internes) & 0,30 -- 0,40 (UE) ; 0,49 (USA) & 0,35 -- 0,45 (Chine 0,38, Brésil 0,53, Inde 0,35) & 0,35 -- 0,50 (très variables) & \textbf{Hausse générale} (sauf AmLat récente). Top 1\% capte croissance. Filet social Triade amortit ; PMA/PED fragiles. \\ \hline
        Émissions CO2 / hab (t/an) & 8 -- 15 (USA 14, UE 6) & 4 -- 8 (Chine 8, Inde 2) & $< 1$ (sauf pétroliers) & \textbf{Injustice climatique :} Triade = \textbf{Responsabilité historique + Empreinte conso}. Émergents = \textbf{Émissions production} (export). PMA = \textbf{Vulnérabilité max}. \\ \hline
        \end{tabularx}
    \end{center}
    
    \tcblower
    \textbf{Synthèse rédigée (Modèle Bac) :}
    La mondialisation produit une \textbf{convergence sélective} : les pays émergents (Chine, Inde, Vietnam) rattrapent la Triade par l'industrialisation exportatrice (part manuf. export $>80\%$, IDE entrants majoritaires 65\%), opérant une \textbf{mise à niveau technologique} (Chine : puces, VE, IA). Inversement, les \textbf{PMA restent marginalisés} (2\% IDE mondiaux, export primaires $>80\%$, croissance souvent inférieure à la démographie), piégés dans une \textbf{spécialisation par le bas} (malédiction des ressources, manque d'infrastructures, instabilité). La Triade conserve la \textbf{commande haute valeur} (R\&D, finance, marques, brevets) et externalise la pollution (empreinte carbone conso $>$ émissions territoriales). Partout, les \textbf{inégalités internes s'accroissent} (Gini $\uparrow$), le capital captant la croissance au détriment du travail (part salaires PIB $\downarrow$), sauf politiques redistributives fortes (modèle nordique). La mondialisation creuse ainsi \textbf{trois fractures} : Nord/Émergents (réduction relative), Émergents/PMA (élargissement), et \emph{intra-nationale} (élites globalisées vs classes populaires exposées à la concurrence).
\end{exoresolu}

\begin{methode}[Analyser une carte / un croquis de la mondialisation (Lecture critique et description commentée)]
    \textbf{Objectif :} Passer de la description (je vois) à l'explication (cela signifie) pour un croquis de flux, de réseaux ou de puissance.
    \begin{enumerate}
        \item \textbf{Lecture d'identité (Légende + Titre) :} Quel thème ? (Flux marchandises, IDE, câbles internet, métropoles, FTN). Quelle date ? Quelle projection ? Quelle échelle ?
        \item \textbf{Description structurée (3 temps) :}
            \begin{itemize}
                \item \textbf{Les grands ensembles (Taches/Surfaces) :} Zones denses / vides. \emph{Ex: « La Triade (Amérique du Nord, Europe de l'Ouest, Asie du Nord-Est) forme un archipel de puissance dense... »}
                \item \textbf{Les lignes / Flux (Flèches) :} Directions (NS, EO, Sud-Sud), épaisseurs (volumes), nature (marchandises, capitaux, données). \emph{Ex: « Deux axes majeurs Est-Ouest : Transpacifique (Chine-US) et Transatlantique (UE-US), flux Nord-Nord dominants mais flux Nord-Sud/Sud-Sud en forte hausse... »}
                \item \textbf{Les points / Nœuds (Villes, Ports, Sièges) :} Hiérarchie (Métropoles mondiales $\alpha, \beta, \gamma$), concentration littorale/portuaire. \emph{Ex: « Concentration des sièges FTN et places boursières dans 10 métropoles (NY, Londres, Tokyo, Paris, Shanghai, Singapour, HK, Pékin, Francfort, Dubaï)... »}
            \end{itemize}
        \item \textbf{Explication / Interprétation (Le « Pourquoi » géographique) :} Relier formes observées aux mécanismes du cours.
            \begin{itemize}
                \item \textbf{Avantages comparatifs / Coûts :} Littoralisation (transport maritime 90\% volume), bas coûts (ZES, main-d'œuvre).
                \item \textbf{Effets d'agglomération / Externalités :} Clusters (Silicon Valley, Shenzhen, Bavière), métropoles (services avancés, aéroports hubs, câbles).
                \item \textbf{Politiques publiques :} ZLE, corridors logistiques (Nouvelles Routes de la Soie), accords commerciaux.
                \item \textbf{Contraintes physiques / Géopolitiques :} Détroits (Malacca, Hormuz, Suez, Panama = goulets d'étranglement), conflits (Mer Rouge 2024), sanctions.
            \end{itemize}
        \item \textbf{Synthèse critique :} La carte montre-t-elle une mondialisation \emph{archipélique} (concentration) vs \emph{diffuse} ? Une \emph{hiérarchie} figée ou \emph{mobile} (montée Asie) ? Des \emph{exclus} (Afrique intérieure, Asie centrale, Amazonie) ?
    \end{enumerate}
\end{methode}

\begin{exoresolu}[Analyse critique d'un croquis : « Les flux de conteneurs et les ports mondiaux 2023 »]
    \textbf{Support (Imaginaire type Bac) :} Carte monde projection Mercator. Flèches bleues proportionnelles flux conteneurs (EVP). Points rouges = Ports > 10 M EVP/an. Hachures bleues = Zones de production manufacturière dense. Hachures rouges = Zones de consommation finale dense.
    \textbf{Consigne :} Commenter ce croquis en montrant l'organisation spatiale du commerce maritime conteneurisé.
    \tcblower
    \textbf{Commentaire structuré (Modèle réponse Bac - 15-20 min)}
    
    \textbf{1. Introduction :} Ce croquis représente les \textbf{flux maritimes de conteneurs (EVP - Équivalent Vingt Pieds)} et les \textbf{ports pivots} en 2023. Il révèle une mondialisation \textbf{archipélique et hiérarchisée}, dominée par l'Asie, structurée par quelques corridors majeurs et des goulets d'étranglement stratégiques.
    
    \textbf{2. Description : Une géométrie en « Y » centrée sur l'Asie.}
    \begin{itemize}
        \item \textbf{Le pôle hégémonique : L'Asie de l'Est.} Concentration exceptionnelle des 10 premiers ports mondiaux (Shanghai \textbf{49 M EVP}, Singapour 39 M, Ningbo 35 M, Shenzhen 29 M, Busan 23 M, Guangzhou, Qingdao, Hong Kong, Tianjin). La Chine seule gère $\sim$ 30\% du trafic mondial. La façade Pacifique asiatique est une \textbf{usine-monde} continue (hachures bleues) du delta de la Rivière des Perles au delta du Yangtsé, jusqu'à la Corée et au Japon.
        \item \textbf{Les deux bras du « Y » : Les corridors Est-Ouest.}
            \begin{itemize}
                \item \textbf{Transpacifique (Bras Nord) :} Flèche la plus épaisse (Chine/Asie $\rightarrow$ Côte Ouest US : LA/Long Beach, Oakland). \textbf{Flux déséquilibré} : conteneurs pleins à l'aller (biens de conso), souvent vides au retour (déséquilibre commercial US/Chine $\sim$ 300 Md\$).
                \item \textbf{Asie-Europe via Suez (Bras Sud) :} Flux massif (Chine $\rightarrow$ Rotterdam, Anvers, Hambourg). Passage obligé par \textbf{détroits de Malacca} (1/3 trafic mondial) et \textbf{canal de Suez} (12\% commerce mondial, 30\% conteneurs). Vulnérabilité géopolitique (Houthis 2024 $\rightarrow$ détour Cap Bonne Espérance +10-14 jours).
            \end{itemize}
        \item \textbf{Le tronc du « Y » : L'axe intra-asiatique / Sud-Sud.} Flèches épaisses entre Chine, ASEAN (Singapour hub de transbordement \#2 mondial), Inde (Jawaharlal Nehru), Moyen-Orient (Dubaï \#10, Jebel Ali). Croissance la plus rapide ($+6-8\%$/an). Singapour = \textbf{pivot obligatoire} (transbordement 85\% trafic).
        \item \textbf{Les « vides » relatifs :} Afrique (seuls Durban, Tanger Med, Lagos émergents), Amérique du Sud (Santos, Buenaventura, Callao - flux faibles, export primaires), Océanie. Intérieurs continentaux déconnectés.
    \end{itemize}
    
    \textbf{3. Explication : Les mécanismes de cette organisation.}
    \begin{itemize}
        \item \textbf{Économie d'échelle / Navires géants (ULCV 24 000 EVP) :} Nécessitent ports en eau profonde, terminaux automatisés, arrière-pays logistique $\rightarrow$ \textbf{Concentration portuaire} (hub-and-spoke). Les « feeders » (navettes) redistribuent depuis Singapour/Busan/Tanger vers ports secondaires.
        \item \textbf{Division Internationale du Travail (DITN) :} Flux reflètent \textbf{fragmentation production}. Composants (Corée/Japon/Taïwan) $\rightarrow$ Assemblage (Chine/Vietnam) $\rightarrow$ Export biens finis (US/UE). Flux « vides » retour = coût logistique de la spécialisation.
        \item \textbf{Gouvernance / Infrastructure :} Ports = interface mer/terre. Investissements massifs États/Collectivités (dragage, rail, autoroutes, ZES). \textbf{Diplomatie des ports :} China Merchants Port / COSCO / PSA Singapore / DP World (Dubaï) / MSC (Suisse) / CMA CGM (France) contrôlent terminaux mondiaux (ex: Piraeus, Hambourg, Duala, Chancay Pérou).
    \end{itemize}
    
    \textbf{4. Conclusion critique :} Ce croquis fige une \textbf{mondialisation asymétrique} : l'Asie est le \textbf{cœur productif et logistique}, l'Occident le \textbf{cœur consommateur et financier}. Les \textbf{goulets d'étranglement} (Malacca, Suez, Panama) sont des \textbf{points de vulnérabilité systémique} (géopolitique, climatique). L'émergence de \textbf{nouvelles routes} (Arctique ? Corridor Moyen-Orient-Europe Inde-Émirats-Israël-Europe IMEC ? Port de Chancay Pérou-Chine direct ?) et la \textbf{robotisation/IA portuaire} redessineront cette carte d'ici 2030.
\end{exoresolu}

\begin{attention}[Les 5 erreurs qui coûtent des points au Bac (Probatoire / Baccalauréat Technique)]
    \begin{enumerate}
        \item \textbf{Confondre « Mondialisation » et « Internationalisation » :} L'internationalisation = échanges \textbf{entre} pays (export/import). La mondialisation = \textbf{intégration} des marchés (biens, capitaux, travail, information) + \textbf{fragmentation de la production} (DITN / chaînes de valeur mondiales) + \textbf{acteurs transnationaux} (FTN) qui échappent partiellement aux États. \emph{Piège :} Dire « la mondialisation c'est le commerce entre la France et l'Allemagne » $\rightarrow$ \textbf{Hors sujet / Note éliminatoire}.
        \item \textbf{Oublier le facteur POLITIQUE (ou ne citer que l'OMC) :} Les facteurs sont \textbf{TRIPLES} : Techniques (conteneur, TIC), Économiques (libre-échange, IDE, FTN), \textbf{Politiques} (Choix États : ouverture unilatérale Chine/Inde/Vietnam, Accords régionaux UE/ALENA/ZLECAf/CPTPP, Rôle FMI/Banque Mondiale/OMC, \textbf{Géopolitique actuelle : friend-shoring, sanctions, subventions vertes IRA/Green Deal}). \emph{Perte : 3 à 4 points sur une dissertation.}
        \item \textbf{Croquis : Flèches non proportionnelles, légende incomplète, pas de titre daté.} \emph{Erreurs fatales :} Flèche Sud-Sud plus grosse que Transpacifique ; Pas de distinction flux marchandises / capitaux / infos (même couleur) ; Pas d'échelle qualitative (« Épaisseur $\propto$ volume ») ; Oubli des métropoles / sièges FTN / ports ; Écriture penchée, coloriage brouillon. \emph{Perte : 5 à 7 points sur 13-15 du croquis.}
        \item \textbf{Calculs : Taux de variation global au lieu du TCAM (ou inversement) sans le dire ; Mauvaise formule ; Unités oubliées (Mds \$ vs \$ ; \% vs points de \%).} \emph{Exemple :} « La part des PED passe de 20\% à 40\%, elle a doublé » $\rightarrow$ \textbf{FAUX} (elle a augmenté de 20 points, soit +100\%). « Croissance de 500\% sur 30 ans = 16\%/an » $\rightarrow$ \textbf{FAUX} (c'est $\sim$6\%/an, effet composé). \emph{Perte : 2 à 3 points faciles.}
        \item \textbf{Rédaction : « Je pense que », « À mon avis », phrases sans verbe, vocabulaire vague (``pays riches`` au lieu de ``Triade / Pays développés / OCDE`` ; ``usines`` au lieu de ``filiales de FTN / plateformes d'exportation / ZES`` ; ``villes`` au lieu de ``métropoles mondiales / hubs logistiques / places financières``).} \emph{Exigence Bac Tech :} Langage \textbf{géographique précis}, \textbf{argumentation structurée} (Thèse/Arguments/Exemples), \textbf{nuance} (ni tout noir, ni tout blanc), \textbf{chiffres-clés} (même approximatifs : « 60\% flux maritimes », « 80\% IDE Triade $\rightarrow$ PED », « 10 métropoles $\alpha$ »). \emph{Perte : 4 à 6 points sur la dissertation/analyse.}
    \end{enumerate}
\end{attention}

\newpage

\coursec{Exercices}

\courssub{Je vérifie mes connaissances}

\begin{exercice}[QCM : les notions clés]
Pour chaque question, entoure la (ou les) bonne(s) réponse(s).
\begin{enumerate}
\item La mondialisation désigne :
\begin{enumerate}
\item la fermeture des économies nationales ;
\item l'intensification des échanges de biens, de capitaux, d'informations et de personnes à l'échelle planétaire ;
\item la disparition des États ;
\item la circulation mondiale des seules matières premières.
\end{enumerate}
\item Une FTN (firme transnationale) est :
\begin{enumerate}
\item une entreprise publique camerounaise ;
\item une entreprise qui implante des filiales dans plusieurs pays ;
\item une organisation internationale comme l'ONU ;
\item une banque centrale.
\end{enumerate}
\item La conteneurisation a permis :
\begin{enumerate}
\item d'augmenter le coût du transport maritime ;
\item de standardiser et de réduire les coûts du transport des marchandises ;
\item de supprimer les ports ;
\item de remplacer le transport aérien.
\end{enumerate}
\item La ZLECAf est :
\begin{enumerate}
\item une zone de libre-échange continentale africaine ;
\item une monnaie unique africaine ;
\item une organisation de pays producteurs de pétrole ;
\item un accord entre le Cameroun et la France.
\end{enumerate}
\end{enumerate}
\end{exercice}

\begin{exercice}[Vrai ou faux — justifier]
Indique si chaque affirmation est vraie ou fausse, puis justifie ta réponse en une ou deux phrases.
\begin{enumerate}
\item Les IDE (investissements directs à l'étranger) sont des placements de capitaux réalisés par des entreprises hors de leur pays d'origine.
\item Les flux mondiaux de marchandises sont répartis de manière égale entre toutes les régions du monde.
\item Le développement d'Internet et des télécommunications est un facteur technique de la mondialisation.
\item Les États n'ont plus aucun rôle dans le fonctionnement de la mondialisation.
\end{enumerate}
\end{exercice}

\begin{exercice}[Définitions]
Définis en deux ou trois lignes chacun des termes suivants : \cle{mondialisation}, \cle{FTN}, \cle{IDE}, \cle{conteneurisation}, \cle{ZLECAf}, \cle{métropole}.
\end{exercice}

\begin{exercice}[Calculs directs sur le commerce mondial]
La valeur des exportations mondiales de marchandises est passée d'environ 2 000 milliards de dollars en 1980 à environ 25 000 milliards de dollars en 2022.
\begin{enumerate}
\item Calcule le taux de croissance des exportations mondiales entre 1980 et 2022 (formule : $\dfrac{V_f - V_i}{V_i} \times 100$).
\item Par combien la valeur des exportations a-t-elle été multipliée ?
\item En 2022, l'Afrique représente environ 3 \% du commerce mondial de marchandises. Calcule la valeur approximative des exportations africaines en milliards de dollars.
\end{enumerate}
\end{exercice}

\courssub{Je m'entraîne}

\begin{exercice}[Les facteurs de la mondialisation]
\begin{enumerate}
\item Cite trois facteurs techniques de la mondialisation et explique en une phrase le rôle de chacun.
\item Cite deux facteurs économiques et deux facteurs politiques de la mondialisation.
\item Recopie et complète le tableau suivant en classant chaque élément dans la bonne colonne : conteneurisation ; libéralisation des échanges ; Internet ; FTN ; accords de libre-échange ; transport aérien.
\end{enumerate}
\begin{center}
\begin{tabularx}{\linewidth}{|X|X|X|}
\hline
\rowcolor{popblueL} Facteurs techniques & Facteurs économiques & Facteurs politiques \\
\hline
 & & \\
\hline
\end{tabularx}
\end{center}
\end{exercice}

\begin{exercice}[Analyse de données chiffrées]
Voici l'évolution de la valeur des exportations mondiales de marchandises (en milliers de milliards de dollars) :
\begin{center}
\begin{tabularx}{\linewidth}{|l|X|X|X|X|}
\hline
\rowcolor{popblueL} Année & 1980 & 1995 & 2010 & 2022 \\
\hline
Valeur (10$^3$ milliards \$) & 2 & 5 & 15 & 25 \\
\hline
\end{tabularx}
\end{center}
\begin{enumerate}
\item Calcule le taux de croissance entre 1995 et 2022, puis entre 1980 et 1995. Quelle période connaît la croissance la plus forte ?
\item Construis un diagramme en bâtons représentant cette évolution (1 cm pour 5 000 milliards de dollars en ordonnée).
\item Rédige un commentaire de 8 à 10 lignes : décris l'évolution, propose une explication à partir des facteurs de la mondialisation étudiés, et nuance en rappelant la place marginale de l'Afrique (environ 3 \% du commerce mondial).
\end{enumerate}
\end{exercice}

\begin{exercice}[Étude de cas : la chaîne de valeur du cacao camerounais]
Le Cameroun produit environ 290 000 tonnes de cacao par an (5\textsuperscript{e} producteur mondial). La répartition de la valeur ajoutée d'une tablette de chocolat vendue 2 000 FCFA en Europe est estimée ainsi : planteur camerounais 6 \%, transport et intermédiaires 12 \%, transformation (chocolatiers européens) 44 \%, distribution et taxes 38 \%.
\begin{enumerate}
\item Identifie les acteurs de cette chaîne de valeur et localise chaque étape (pays producteur / pays consommateur).
\item Calcule la somme perçue par le planteur camerounais sur une tablette vendue 2 000 FCFA.
\item Compare cette somme à celle perçue par la transformation. Que révèle cet écart sur la place du Cameroun dans la mondialisation ?
\item Propose une solution pour que le Cameroun capte une part plus importante de la valeur ajoutée.
\end{enumerate}
\end{exercice}

\begin{exercice}[Les acteurs de la mondialisation]
\begin{enumerate}
\item À l'aide d'exemples précis (une FTN présente au Cameroun, une métropole mondiale, un État), montre le rôle de chacun des trois grands acteurs de la mondialisation.
\item Pourquoi dit-on que les métropoles sont les « centres de commandement » de l'économie mondiale ? Donne deux exemples de métropoles mondiales et un exemple de métropole africaine.
\item Rédige un paragraphe de 10 lignes maximum montrant comment ces trois acteurs interagissent dans le fonctionnement de la mondialisation.
\end{enumerate}
\end{exercice}

\courssub{Je me prépare au Bac}

\begin{exercice}[Sujet type Baccalauréat technique : les flux mondiaux et la place du Cameroun]
\textbf{Documents mis à disposition :}
\begin{itemize}
\item Document 1 : carte des principaux flux mondiaux de marchandises (routes maritimes reliant l'Asie de l'Est, l'Europe de l'Ouest et l'Amérique du Nord ; ports secondaires sur la façade atlantique africaine : Douala, Lagos, Abidjan).
\item Document 2 : tableau des principales exportations du Cameroun.
\end{itemize}
\begin{center}
\begin{tabularx}{\linewidth}{|l|X|X|}
\hline
\rowcolor{popblueL} Produit & Part dans les exportations (\%) & Destinations principales \\
\hline
Pétrole brut & 40 & Europe, Asie \\
\hline
Cacao & 12 & Pays-Bas, Belgique \\
\hline
Bois & 10 & Chine, Europe \\
\hline
Coton & 6 & Asie \\
\hline
Autres (bananes, café, aluminium) & 32 & Divers \\
\hline
\end{tabularx}
\end{center}
\textbf{Travail à faire :}
\begin{enumerate}
\item Décris les grandes orientations des flux mondiaux de marchandises présentés dans le document 1 (trois pôles dominants, routes maritimes majeures).
\item À partir du document 2, caractérise les exportations camerounaises (nature des produits, destinations).
\item Dégage deux facteurs techniques et deux facteurs économiques qui expliquent l'intensification de ces flux.
\item Conclus en une dizaine de lignes sur la place du Cameroun dans la mondialisation : pays intégré, marginalisé ou intermédiaire ? Justifie ta réponse.
\end{enumerate}
\end{exercice}

\begin{exercice}[Construction de croquis : les grandes routes maritimes mondiales]
Sur un fond de planisphère muet fourni par l'enseignant (ou reproduit à main levée), réalise un croquis intitulé « Les grandes routes maritimes mondiales et la façade atlantique africaine ».
\begin{enumerate}
\item Place les trois pôles de la Triade (Amérique du Nord, Europe de l'Ouest, Asie de l'Est).
\item Trace les trois grandes routes maritimes : route de l'Atlantique Nord, route Europe--Asie par le canal de Suez, route du Pacifique Nord. Utilise des flèches d'épaisseur proportionnelle à l'importance du trafic.
\item Localise les détroits stratégiques (canal de Suez, détroit de Malacca, canal de Panama) et les ports de Douala et Lagos sur la façade atlantique africaine.
\item Rédige la légende (au moins quatre éléments), ajoute le titre, l'orientation et une échelle indicative.
\item Rédige un court texte d'accompagnement (5 à 8 lignes) expliquant ce que ce croquis révèle du fonctionnement de la mondialisation et de la position de l'Afrique.
\end{enumerate}
\end{exercice}

\begin{exercice}[Débat argumenté : la mondialisation, une chance pour le Cameroun ?]
\textbf{Consigne :} « La mondialisation est-elle une chance pour l'économie camerounaise ? » Rédige un texte argumenté de 25 à 30 lignes comprenant :
\begin{enumerate}
\item une introduction dans laquelle tu définis la mondialisation et tu poses clairement la problématique ;
\item deux arguments montrant les opportunités offertes au Cameroun (par exemple : accès aux marchés extérieurs, IDE et implantation de FTN, ZLECAf), chacun justifié par un exemple précis ;
\item deux arguments montrant les limites ou les risques (par exemple : exportation de matières premières brutes, faible part de la valeur ajoutée, concurrence des produits importés), chacun justifié par un exemple précis ;
\item une conclusion personnelle nuancée proposant une condition pour que le Cameroun tire mieux profit de la mondialisation.
\end{enumerate}
\textbf{Critères d'évaluation :} exactitude des notions, pertinence des exemples, cohérence de l'argumentation, qualité de la rédaction.
\end{exercice}

\newpage

\coursec{Corrigés des exercices}

\coursec{Leçon 04 : Les mécanismes de la mondialisation — Exercices corrigés}

\begin{corrige}[Exercice 1 — QCM : les notions clés]
\begin{enumerate}
\item Bonne réponse : \textbf{b}. La mondialisation est l'intensification des échanges de biens, de capitaux, d'informations et de personnes à l'échelle planétaire. Les réponses a et c sont fausses (les économies s'ouvrent et les États subsistent) ; la réponse d est fausse car les flux concernent aussi les services, les capitaux et les informations.
\item Bonne réponse : \textbf{b}. Une FTN est une entreprise qui implante des filiales dans plusieurs pays (exemple : TotalEnergies, MTN). Ce n'est ni une entreprise publique nationale, ni une organisation internationale, ni une banque centrale.
\item Bonne réponse : \textbf{b}. La conteneurisation standardise le transport des marchandises (boîtes de dimensions normalisées) et réduit fortement les coûts et les délais de manutention. Elle n'a ni supprimé les ports (au contraire, elle les a modernisés), ni remplacé le transport aérien.
\item Bonne réponse : \textbf{a}. La ZLECAf (Zone de libre-échange continentale africaine), entrée en vigueur en 2021, vise à supprimer progressivement les barrières douanières entre les pays africains.
\end{enumerate}
\end{corrige}

\begin{corrige}[Exercice 2 — Vrai ou faux — justifier]
\begin{enumerate}
\item \textbf{Vrai.} Les IDE sont des investissements réalisés par une entreprise hors de son pays d'origine (création ou rachat de filiales, usines). Exemple : une FTN qui construit une usine au Cameroun réalise un IDE.
\item \textbf{Faux.} Les flux sont très inégalement répartis : ils se concentrent entre les trois pôles de la Triade (Amérique du Nord, Europe de l'Ouest, Asie de l'Est), tandis que l'Afrique ne pèse qu'environ 3 \% du commerce mondial.
\item \textbf{Vrai.} Internet et les télécommunications permettent la circulation instantanée et quasi gratuite de l'information à l'échelle planétaire : c'est un facteur technique majeur de la mondialisation, au même titre que les transports.
\item \textbf{Faux.} Les États restent des acteurs essentiels : ils signent les accords commerciaux, fixent les règles (douanes, fiscalité), aménagent les infrastructures (ports, routes) et peuvent attirer ou encadrer les IDE.
\end{enumerate}
\end{corrige}

\begin{corrige}[Exercice 3 — Définitions]
\begin{itemize}
\item \cle{Mondialisation} : processus d'intensification et d'accélération des échanges de biens, de services, de capitaux, d'informations et de personnes à l'échelle de la planète, qui met les territoires en interdépendance.
\item \cle{FTN} (firme transnationale) : grande entreprise qui organise sa production et ses ventes à l'échelle mondiale grâce à des filiales implantées dans plusieurs pays (exemple : MTN, TotalEnergies).
\item \cle{IDE} (investissement direct à l'étranger) : placement de capitaux réalisé par une entreprise hors de son pays d'origine pour y créer, racheter ou développer une activité productive.
\item \cle{Conteneurisation} : transport des marchandises dans des conteneurs normalisés, facilement transférables du navire au train ou au camion, ce qui réduit les coûts et les délais du commerce maritime.
\item \cle{ZLECAf} : Zone de libre-échange continentale africaine, accord entré en vigueur en 2021 visant à créer un marché unique africain par la réduction des droits de douane entre pays du continent.
\item \cle{Métropole} : grande ville qui concentre des fonctions de commandement (sièges sociaux, banques, bourses, décision politique) et qui rayonne sur un vaste espace, parfois sur le monde entier.
\end{itemize}
\end{corrige}

\begin{corrige}[Exercice 4 — Calculs directs sur le commerce mondial]
\begin{enumerate}
\item Taux de croissance entre 1980 et 2022 :
\[ \frac{V_f - V_i}{V_i} \times 100 = \frac{25\,000 - 2\,000}{2\,000} \times 100 = \frac{23\,000}{2\,000} \times 100 = 1\,150\ \% \]
Les exportations mondiales ont crû d'environ \textbf{1 150 \%} entre 1980 et 2022.
\item Coefficient multiplicateur : $\dfrac{25\,000}{2\,000} = 12,5$. La valeur des exportations a été multipliée par \textbf{12,5}.
\item Part de l'Afrique : $25\,000 \times \dfrac{3}{100} = 750$ milliards de dollars. Les exportations africaines représentent environ \textbf{750 milliards de dollars}.
\end{enumerate}
\end{corrige}

\begin{attention}[Points de vigilance — Exercice 4]
Bien soustraire la valeur initiale avant de diviser (erreur fréquente : calculer $\frac{V_f}{V_i}\times 100 = 1\,250\ \%$, qui est le coefficient en pourcentage et non le taux de croissance). Ne pas confondre « multiplié par 12,5 » et « augmenté de 12,5 \% ».
\end{attention}

\begin{corrige}[Exercice 5 — Les facteurs de la mondialisation]
\begin{enumerate}
\item Trois facteurs techniques :
\begin{itemize}
\item la \cle{conteneurisation} : elle standardise le transport maritime et réduit fortement les coûts d'expédition des marchandises ;
\item le \cle{transport aérien} : il permet la circulation rapide des personnes et des marchandises à forte valeur à l'échelle planétaire ;
\item \cle{Internet} et les télécommunications : ils assurent la transmission instantanée des informations, des commandes et des capitaux.
\end{itemize}
\item Facteurs économiques : la \cle{libéralisation des échanges} (baisse des droits de douane) et le rôle moteur des \cle{FTN} qui organisent la production à l'échelle mondiale. Facteurs politiques : la fin de la Guerre froide (ouverture des anciens pays socialistes au commerce mondial) et la multiplication des \cle{accords de libre-échange} (OMC créée en 1995, ZLECAf en 2021).
\item Tableau complété :
\begin{center}
\begin{tabularx}{\linewidth}{|X|X|X|}
\hline
\rowcolor{popblueL} Facteurs techniques & Facteurs économiques & Facteurs politiques \\
\hline
Conteneurisation ; Internet ; transport aérien & Libéralisation des échanges ; FTN & Accords de libre-échange \\
\hline
\end{tabularx}
\end{center}
Remarque : les FTN sont à la fois des acteurs et un facteur économique de la mondialisation ; la libéralisation des échanges (facteur économique) est rendue possible par les accords politiques.
\end{enumerate}
\end{corrige}

\begin{corrige}[Exercice 6 — Analyse de données chiffrées]
\begin{enumerate}
\item Entre 1995 et 2022 : $\dfrac{25-5}{5}\times 100 = \dfrac{20}{5}\times 100 = 400\ \%$. Entre 1980 et 1995 : $\dfrac{5-2}{2}\times 100 = \dfrac{3}{2}\times 100 = 150\ \%$. La période 1995--2022 connaît la croissance la plus forte (400 \% contre 150 \%).
\item Diagramme en bâtons (échelle : 1 cm pour 5 000 milliards de dollars, soit des hauteurs de 0,4 cm ; 1 cm ; 3 cm ; 5 cm) :
\begin{popfigure}
\begin{center}
\begin{tikzpicture}[scale=1]
\draw[->,thick] (0,0) -- (7.2,0) node[right]{\footnotesize Année};
\draw[->,thick] (0,0) -- (0,5.6) node[above]{\footnotesize $10^3$ milliards \$};
\foreach \y/\lab in {1/5,2/10,3/15,4/20,5/25}{\draw (0,\y) -- (-0.08,\y) node[left]{\footnotesize \lab};}
\fill[popblue] (0.7,0) rectangle (1.5,0.4);
\fill[popblue] (2.3,0) rectangle (3.1,1);
\fill[popblue] (3.9,0) rectangle (4.7,3);
\fill[popblue] (5.5,0) rectangle (6.3,5);
\node[below] at (1.1,0){\footnotesize 1980};
\node[below] at (2.7,0){\footnotesize 1995};
\node[below] at (4.3,0){\footnotesize 2010};
\node[below] at (5.9,0){\footnotesize 2022};
\node[above] at (1.1,0.4){\footnotesize 2};
\node[above] at (2.7,1){\footnotesize 5};
\node[above] at (4.3,3){\footnotesize 15};
\node[above] at (5.9,5){\footnotesize 25};
\end{tikzpicture}
\end{center}
\legende{Évolution de la valeur des exportations mondiales de marchandises, 1980--2022 (en milliers de milliards de dollars). Échelle : 1 cm pour 5 000 milliards de dollars. Source : OMC, données arrondies.}
\end{popfigure}
\item Commentaire attendu (8 à 10 lignes) : la valeur des exportations mondiales de marchandises est passée de 2 000 à 25 000 milliards de dollars entre 1980 et 2022, soit une multiplication par 12,5. La croissance est d'abord modérée (150 \% entre 1980 et 1995), puis s'accélère fortement (400 \% entre 1995 et 2022). Cette explosion s'explique par les facteurs techniques de la mondialisation : la conteneurisation a abaissé les coûts du transport maritime, tandis qu'Internet a fluidifié les échanges d'informations et de capitaux. S'y ajoutent des facteurs économiques (stratégies mondiales des FTN, libéralisation des échanges) et politiques (création de l'OMC en 1995, intégration de la Chine au commerce mondial). Cependant, cette croissance profite inégalement aux territoires : l'Afrique ne représente qu'environ 3 \% du commerce mondial, soit 750 milliards de dollars, ce qui montre que la mondialisation est un processus sélectif qui marginalise certains espaces.
\end{enumerate}
\end{corrige}

\begin{attention}[Points de vigilance — Exercice 6]
Le diagramme doit respecter l'échelle indiquée (1 cm pour 5 000 milliards \$). Les barres doivent être de même largeur, espacées régulièrement. Le commentaire doit articuler description chiffrée (données du graphique) et explication (facteurs vus en cours).
\end{attention}

\begin{corrige}[Exercice 7 — Étude de cas : la chaîne de valeur du cacao camerounais]
\begin{enumerate}
\item Acteurs et localisation : le \textbf{planteur camerounais} produit la fève (pays producteur, Sud) ; les \textbf{transporteurs et intermédiaires} (acheteurs, exportateurs, port de Douala) assurent la collecte et l'acheminement ; les \textbf{chocolatiers européens} transforment la fève en chocolat (pays consommateur, Nord) ; la \textbf{distribution} (grandes surfaces) et les États (taxes) interviennent dans les pays consommateurs.
\item Part du planteur : $2\,000 \times \dfrac{6}{100} = 120$ FCFA. Le planteur perçoit \textbf{120 FCFA} par tablette.
\item Part de la transformation : $2\,000 \times \dfrac{44}{100} = 880$ FCFA. La transformation capte 880 FCFA, soit environ \textbf{7,3 fois plus} que le planteur ($\frac{880}{120} \approx 7,3$). Cet écart révèle que le Cameroun, inséré dans la mondialisation comme fournisseur de matière première brute, capte une part très faible de la valeur ajoutée : les activités les plus rémunératrices (transformation, distribution) sont réalisées dans les pays du Nord.
\item Solution : développer la \textbf{transformation locale} du cacao (usines de broyage et chocolateries au Cameroun), afin d'exporter des produits finis ou semi-finis à plus forte valeur ajoutée ; on peut aussi citer la formation des planteurs, l'organisation en coopératives et le soutien de l'État.
\end{enumerate}
\end{corrige}

\begin{corrige}[Exercice 8 — Les acteurs de la mondialisation]
\begin{enumerate}
\item \textbf{FTN} : MTN (téléphonie) ou TotalEnergies (pétrole et distribution), implantées au Cameroun par des filiales, y investissent (IDE), y organisent production et ventes, et relient le pays aux réseaux mondiaux. \textbf{Métropole} : New York concentre bourse, banques et sièges sociaux qui orientent les flux de capitaux mondiaux. \textbf{État} : l'État camerounais aménage le port de Kribi, signe des accords commerciaux et fixe les règles d'accueil des investissements.
\item Les métropoles sont des « centres de commandement » car elles concentrent les fonctions de décision : sièges sociaux des FTN, bourses, grandes banques, institutions politiques. Exemples de métropoles mondiales : New York, Londres, Tokyo (ou Paris). Métropole africaine : Johannesburg (ou Le Caire, Lagos).
\item Paragraphe attendu : les FTN organisent la production et les échanges à l'échelle planétaire en implantant leurs filiales là où les coûts sont faibles et les marchés porteurs. Elles dirigent leurs réseaux depuis les métropoles mondiales, où se prennent les décisions financières et stratégiques. Les États, loin d'avoir disparu, créent le cadre de ce fonctionnement : ils négocient les accords de libre-échange, construisent les infrastructures (ports, aéroports, télécommunications) et cherchent à attirer les IDE. Les trois acteurs sont donc interdépendants : les FTN ont besoin des États et des métropoles, les États ont besoin des investissements des FTN, et les métropoles concentrent le commandement de l'ensemble.
\end{enumerate}
\end{corrige}

\begin{corrige}[Exercice 9 — Sujet type Baccalauréat technique : les flux mondiaux et la place du Cameroun]
\textbf{Barème indicatif (sur 20) :} question 1 : 4 pts ; question 2 : 4 pts ; question 3 : 4 pts ; question 4 : 6 pts ; présentation et exactitude des notions : 2 pts.
\begin{enumerate}
\item Le document 1 montre que les flux mondiaux de marchandises s'organisent autour de \textbf{trois pôles dominants} (la Triade) : l'Asie de l'Est, l'Europe de l'Ouest et l'Amérique du Nord. Les routes maritimes majeures relient ces pôles entre eux : route de l'Atlantique Nord, route Europe--Asie par le canal de Suez, route du Pacifique Nord. Les ports africains (Douala, Lagos, Abidjan) n'apparaissent que comme des relais secondaires sur la façade atlantique.
\item Les exportations camerounaises sont dominées par les \textbf{matières premières brutes} : le pétrole brut pèse à lui seul 40 \% des exportations, suivi du cacao (12 \%), du bois (10 \%) et du coton (6 \%). Les destinations sont les pays développés et émergents (Europe, Chine, Asie), ce qui traduit une relation d'échange inégale : produits bruts contre produits manufacturés.
\item Facteurs techniques : la \textbf{conteneurisation}, qui a réduit les coûts du transport maritime, et le développement des \textbf{télécommunications et d'Internet}, qui fluidifient les commandes et les paiements. Facteurs économiques : la \textbf{libéralisation des échanges} (baisse des barrières douanières) et les stratégies mondiales des \textbf{FTN}, qui organisent les filières (pétrole, cacao) à l'échelle planétaire.
\item Conclusion attendue : le Cameroun est un pays \textbf{intermédiaire, inséré mais en position subordonnée} dans la mondialisation. Il est bien intégré aux flux mondiaux : il exporte pétrole, cacao, bois et coton vers l'Europe et l'Asie, et son port de Douala est connecté aux routes maritimes. Mais cette intégration est inégale : il vend des matières premières brutes à faible valeur ajoutée et importe des produits manufacturés coûteux. De plus, l'Afrique entière ne pèse qu'environ 3 \% du commerce mondial. Le Cameroun n'est donc ni pleinement intégré au cœur du système, ni totalement marginalisé : il occupe une position périphérique, qu'une transformation locale des produits pourrait améliorer.
\end{enumerate}
\end{corrige}

\begin{attention}[Points de vigilance — Exercice 9]
Exiger l'usage précis des documents (chiffres du tableau cités : 40 \% de pétrole). À la question 4, le candidat doit choisir une position et la \emph{justifier} par au moins deux arguments issus des documents. Éviter les réponses générales non appuyées sur les données.
\end{attention}

\begin{corrige}[Exercice 10 — Construction de croquis : les grandes routes maritimes mondiales]
\textbf{Barème indicatif (sur 20) :} localisation des trois pôles : 3 pts ; tracé des trois routes avec flèches proportionnelles : 5 pts ; détroits et ports africains : 4 pts ; légende, titre, orientation, échelle : 4 pts ; texte d'accompagnement : 4 pts.
\begin{enumerate}
\item Les trois pôles sont placés en hachures ou couleurs distinctes : Amérique du Nord (côte est surtout), Europe de l'Ouest (façade atlantique), Asie de l'Est (Japon, Chine littorale).
\item Routes tracées : \textbf{Atlantique Nord} (Europe de l'Ouest -- côte est des États-Unis, flèche très épaisse) ; \textbf{Europe--Asie} par la Méditerranée, le canal de Suez, l'océan Indien et le détroit de Malacca (flèche très épaisse, la plus fréquentée) ; \textbf{Pacifique Nord} (Asie de l'Est -- côte ouest de l'Amérique du Nord, flèche épaisse). La route passant par la façade atlantique africaine est tracée en flèche fine.
\item Détroits et canaux figurés par un symbole (étoile ou point) : canal de Suez, détroit de Malacca, canal de Panama ; ports de Douala et Lagos localisés par des points nommés sur le golfe de Guinée.
\item Légende (au moins quatre éléments) : 1. pôles de la Triade ; 2. route maritime majeure (trafic très dense) ; 3. route maritime secondaire ; 4. détroit ou canal stratégique ; 5. port africain. Le titre est placé en haut, la flèche du nord en haut à droite, l'échelle indicative (par exemple 1 cm = 4 000 km) en bas.
\item Texte d'accompagnement attendu : ce croquis montre que le commerce mondial de marchandises emprunte un petit nombre de routes maritimes concentrées entre les trois pôles de la Triade, qui produisent et consomment l'essentiel des richesses échangées. Les détroits et canaux (Suez, Malacca, Panama) sont des points de passage obligés, donc stratégiques. L'Afrique apparaît en marge de ces grands axes : ses ports, comme Douala et Lagos, sont des relais secondaires qui exportent surtout des matières premières vers les pôles. Le croquis révèle ainsi le fonctionnement sélectif de la mondialisation, qui concentre les flux et marginalise certains espaces.
\end{enumerate}
\end{corrige}

\begin{attention}[Points de vigilance — Exercice 10]
Un croquis sans titre, sans légende ou sans orientation est sanctionné. Les flèches doivent être réellement proportionnelles (Atlantique Nord et Suez plus épaisses que la façade africaine). Ne pas surcharger : chaque élément figuré doit figurer dans la légende, et réciproquement.
\end{attention}

\begin{corrige}[Exercice 11 — Débat argumenté : la mondialisation, une chance pour le Cameroun ?]
\textbf{Éléments de correction (texte de 25 à 30 lignes).}

\textbf{Introduction attendue :} définition de la mondialisation (intensification des échanges de biens, capitaux, informations et personnes à l'échelle planétaire) ; présentation du Cameroun, pays en développement ouvert sur l'extérieur ; problématique : la mondialisation est-elle une opportunité ou un risque pour l'économie camerounaise ? Annonce du plan (opportunités, puis limites, puis position nuancée).

\textbf{Opportunités (deux arguments exemplifiés) :}
\begin{itemize}
\item \emph{Accès aux marchés extérieurs} : le Cameroun écoule ses produits dans le monde entier — le pétrole brut (40 \% des exportations) vers l'Europe et l'Asie, le cacao vers les Pays-Bas et la Belgique —, ce qui procure des devises indispensables à l'économie.
\item \emph{IDE et implantation de FTN} : des firmes comme MTN ou TotalEnergies investissent au Cameroun, créent des emplois, transfèrent des technologies et modernisent les infrastructures (port en eau profonde de Kribi). La \cle{ZLECAf}, entrée en vigueur en 2021, ouvre en outre un marché continental de plus d'un milliard de consommateurs aux produits camerounais.
\end{itemize}

\textbf{Limites et risques (deux arguments exemplifiés) :}
\begin{itemize}
\item \emph{Exportation de matières premières brutes} : le Cameroun vend pétrole, cacao et bois non transformés ; sur une tablette de chocolat vendue 2 000 FCFA en Europe, le planteur camerounais ne capte que 6 \% (120 FCFA) contre 44 \% pour la transformation réalisée au Nord.
\item \emph{Concurrence des produits importés} : l'ouverture des frontières expose les entreprises et les paysans camerounais à la concurrence de produits manufacturés et agricoles importés moins chers, ce qui freine l'industrie locale ; l'Afrique entière ne pèse qu'environ 3 \% du commerce mondial.
\end{itemize}

\textbf{Conclusion attendue :} position personnelle nuancée — la mondialisation est une chance \emph{à condition} que le Cameroun transforme localement ses matières premières (chocolateries, scieries, raffinage), s'appuie sur la ZLECAf et forme sa main-d'œuvre, afin de capter une part plus grande de la valeur ajoutée au lieu de rester un simple fournisseur de produits bruts.

\textbf{Barème indicatif (sur 20) :} introduction (définition + problématique + plan) : 4 pts ; deux opportunités justifiées par des exemples : 6 pts ; deux limites justifiées par des exemples : 6 pts ; conclusion nuancée avec condition : 2 pts ; qualité de la rédaction : 2 pts.
\end{corrige}

\begin{attention}[Points de vigilance — Exercice 11]
Chaque argument doit être étayé par un exemple chiffré ou localisé (un argument sans exemple perd la moitié des points). La conclusion doit être personnelle et nuancée : une réponse manichéenne (« tout bon » ou « tout mauvais ») est insuffisante. Veiller à l'exactitude des notions : ne pas confondre ZLECAf et monnaie unique, ni FTN et organisation internationale.
\end{attention}

\newpage

\coursec{Fiche bilan}

\begin{popfigure}
\begin{tikzpicture}[mindmap, grow cyclic, every node/.style={concept, font=\small},
  root concept/.append style={concept color=popblue, font=\bfseries},
  level 1/.append style={level distance=3.4cm, sibling angle=60, font=\footnotesize},
  level 2/.append style={level distance=2.2cm, sibling angle=45, font=\scriptsize}]
\node[root concept]{Mécanismes de la mondialisation}
  child[concept color=poporange]{ node{Définition} 
    child{ node{Mise en relation des espaces} }
    child{ node{Interdépendances mondiales} } }
  child[concept color=popgreen]{ node{Facteurs techniques}
    child{ node{Transports (conteneur)} }
    child{ node{TIC, Internet} } }
  child[concept color=poppurple]{ node{Facteurs éco-politiques}
    child{ node{Libéralisation (OMC)} }
    child{ node{Déréglementation} } }
  child[concept color=poppink]{ node{Les flux}
    child{ node{Biens, capitaux} }
    child{ node{Hommes, informations} } }
  child[concept color=popgold]{ node{Acteurs : FTN}
    child{ node{Stratégies mondiales} }
    child{ node{Délocalisations} } }
  child[concept color=popblue!60]{ node{Acteurs : États, métropoles}
    child{ node{Gouvernance, IDE} }
    child{ node{Commandement mondial} } };
\end{tikzpicture}
\legende{Carte mentale de la leçon : les mécanismes de la mondialisation (définition, facteurs, flux, acteurs).}
\end{popfigure}

\begin{bilan}
\textbf{Définitions clés}
\begin{itemize}
  \item \cle{Mondialisation} : processus de mise en relation croissante des espaces, des économies et des sociétés à l'échelle de la planète, créant des interdépendances.
  \item \cle{Flux} : mouvements de biens, de capitaux, d'hommes et d'informations entre les espaces.
  \item \cle{FTN} (firme transnationale) : entreprise implantée dans plusieurs pays, organisant sa production à l'échelle mondiale.
  \item \cle{IDE} : investissements directs à l'étranger, principaux flux de capitaux.
  \item \cle{Métropole} : ville concentrant des fonctions de commandement économique et politique.
\end{itemize}
\textbf{Les mécanismes à connaître par cœur}
\begin{center}
\begin{tabularx}{\linewidth}{|l|X|}\hline
\rowcolor{popblueL} \textbf{Mécanisme} & \textbf{Rôle dans la mondialisation} \\ \hline
Transports (conteneur, aviation) & Baisse des coûts et des temps de transport \\ \hline
TIC et Internet & Circulation instantanée de l'information \\ \hline
Libéralisation (OMC, 1995) & Abaissement des barrières douanières \\ \hline
Stratégies des FTN & Délocalisations, division internationale du travail \\ \hline
\end{tabularx}
\end{center}
\textbf{Méthodes express}
\begin{itemize}
  \item Croquis : légende organisée (flux = flèches d'épaisseur proportionnelle, acteurs = symboles, espaces = aplat).
  \item Raisonnement : toujours relier un facteur technique à un effet économique concret (ex. conteneur $\rightarrow$ baisse des coûts $\rightarrow$ hausse des échanges).
  \item Exemples à mobiliser : OMC, Coca-Cola ou Samsung (FTN), Douala et Yaoundé (métropoles camerounaises).
\end{itemize}
\end{bilan}

\begin{aretenir}[Checklist avant le Bac]
\begin{itemize}
  \item $\square$ Je sais définir la mondialisation en une phrase complète.
  \item $\square$ Je cite les deux grandes catégories de facteurs : techniques et économico-politiques.
  \item $\square$ J'explique le rôle du conteneur et des TIC dans l'accélération des échanges.
  \item $\square$ Je connais le rôle de l'OMC dans la libéralisation des échanges.
  \item $\square$ Je distingue les quatre types de flux (biens, capitaux, hommes, informations).
  \item $\square$ Je définis une FTN et je cite au moins un exemple.
  \item $\square$ Je sais ce qu'est un IDE et qui en sont les principaux émetteurs.
  \item $\square$ J'explique le rôle des métropoles dans le commandement de l'économie mondiale.
  \item $\square$ Je sais que les États restent des acteurs (gouvernance, accords commerciaux).
  \item $\square$ Je sais construire un croquis avec titre, légende organisée et flèche du nord.
  \item $\square$ Je sais rédiger une conclusion justifiée avec un exemple camerounais.
\end{itemize}
\end{aretenir}

\findecours

\end{document}
